\begin{center}\rule{0.5\linewidth}{0.5pt}\end{center}

\textbf{UNIVERSIDAD NACIONAL DE SAN AGUSTÍN DE AREQUIPA}

Facultad de Ingeniería de Producción y Servicios

Escuela Profesional de Ingeniería de Sistemas

\begin{center}\rule{0.5\linewidth}{0.5pt}\end{center}

\textbf{SISTEMA INTEGRADO DE SUPERVISIÓN OPERATIVA CON INTELIGENCIA
ARTIFICIAL EXPLICABLE PARA LA DETECCIÓN DE ANOMALÍAS Y GENERACIÓN DE
REPORTES TRAZABLES EN EMPRESAS AGROEXPORTADORAS PERUANAS}

\begin{center}\rule{0.5\linewidth}{0.5pt}\end{center}

Tesis presentada por el Bachiller:

\textbf{Cozco Mauri, Yoset}

Para optar el Título Profesional de Ingeniero de Sistemas

Asesor:

\textbf{Dr.~Víctor Manuel Cornejo Aparicio}

\begin{center}\rule{0.5\linewidth}{0.5pt}\end{center}

\textbf{Arequipa -- Perú}

\textbf{2026}

\begin{center}\rule{0.5\linewidth}{0.5pt}\end{center}

\section{DEDICATORIA}\label{dedicatoria}

\emph{(Por completar --- dedicatoria personal del autor)}

\begin{center}\rule{0.5\linewidth}{0.5pt}\end{center}

\section{AGRADECIMIENTOS}\label{agradecimientos}

\emph{(Por completar --- agradecimientos al asesor, jurado, institución
colaboradora y familiares)}

\begin{center}\rule{0.5\linewidth}{0.5pt}\end{center}

\section{PRESENTACIÓN}\label{presentaciuxf3n}

La presente investigación tiene como propósito el diseño, implementación
y evaluación de un sistema integrado de supervisión operativa para
empresas agroexportadoras peruanas, que combina técnicas de aprendizaje
automático, detección de anomalías, explicabilidad algorítmica y
generación automática de reportes trazables mediante modelos de lenguaje
con recuperación de contexto. El estudio responde a la necesidad de
contar con herramientas de inteligencia artificial capaces de detectar
desviaciones operativas, explicar sus causas probables y documentar
alertas de manera comprensible para supervisores, responsables de
calidad, gestores logísticos y auditores internos.

El trabajo se organiza en cinco capítulos: el Capítulo I establece el
planteamiento del problema, objetivos e hipótesis; el Capítulo II
desarrolla el marco teórico con antecedentes y estado del arte; el
Capítulo III describe la propuesta metodológica; el Capítulo IV presenta
los resultados y discusión; y el Capítulo V contiene las conclusiones y
trabajos futuros.

\textbf{El Autor}

\begin{center}\rule{0.5\linewidth}{0.5pt}\end{center}

\begin{center}\rule{0.5\linewidth}{0.5pt}\end{center}

\section{RESUMEN}\label{resumen}

Esta tesis propone un sistema integrado de supervisión operativa para
empresas agroexportadoras peruanas, que combina predicción tabular
mediante modelos Gradient Boosting Decision Trees (GBDT), detección de
anomalías operativas mediante ensemble de algoritmos, explicabilidad
mediante SHAP (SHapley Additive exPlanations), y generación automática
de reportes trazables con Modelos de Lenguaje de Gran Tamaño (LLMs) en
arquitectura RAG (Retrieval-Augmented Generation).

El sistema aborda una brecha identificada en el contexto agroexportador:
los procesos de producción, acopio, almacenamiento, control de calidad,
logística, cumplimiento fitosanitario y comercialización internacional
suelen analizarse mediante fuentes fragmentadas, reportes manuales o
tableros aislados. Esta fragmentación dificulta la detección temprana de
anomalías y reduce la trazabilidad de las decisiones. La propuesta se
evalúa con métricas técnicas (PR-AUC, F1-Score, cobertura de
trazabilidad), evaluación de comprensión operativa y datos
públicos/sintéticos documentados del dominio agroexportador.

Las contribuciones principales son: (1) arquitectura modular de cuatro
capas que separa predicción, detección, explicación y reporte; (2)
integración de fuentes públicas oficiales y dataset sintético
agroexportador documentado mediante criterios de Datasheets for Datasets
(Gebru et al., 2021); (3) uso de SHAP para explicar alertas operativas a
nivel de variable; (4) generación de reportes mediante RAG restringido a
evidencias estructuradas, reduciendo el riesgo de alucinación; (5)
evaluación comparativa del sistema integrado frente a componentes
aislados en rendimiento, trazabilidad y tiempo de interpretación. La
Resolución SBS N° 053-2023 se considera como referencia nacional de
buenas prácticas para gestión de riesgo de modelos, mientras que el D.S.
N° 115-2025-PCM se adopta como marco peruano general de gobernanza y
supervisión humana en IA.

\textbf{Palabras clave}: supervisión operativa, detección de anomalías,
agroexportación, explicabilidad IA, modelos de lenguaje, gobernanza,
GBDT, reportes automáticos, trazabilidad, inteligencia artificial.

\begin{center}\rule{0.5\linewidth}{0.5pt}\end{center}

\section{ABSTRACT}\label{abstract}

This thesis proposes an integrated operational supervision system for
Peruvian agro-export companies, combining tabular prediction using
Gradient Boosting Decision Trees (GBDT), operational anomaly detection
through an ensemble of algorithms, explainability through SHAP (SHapley
Additive exPlanations), and traceable automatic report generation with
Large Language Models (LLMs) in a Retrieval-Augmented Generation (RAG)
architecture.

The system addresses an identified gap in agro-export operational
supervision: production, storage, quality control, logistics,
phytosanitary compliance, and international commercialization are
commonly analyzed through fragmented sources, manual reports, or
isolated dashboards. This fragmentation limits early anomaly detection
and weakens decision traceability. The proposal is evaluated using
technical metrics (PR-AUC, F1-Score, traceability coverage), operational
comprehension assessment, and documented public/synthetic agro-export
data.

The main contributions are: (1) a modular four-layer architecture
separating prediction, detection, explanation, and reporting; (2)
integration of official public sources and a documented synthetic
agro-export dataset; (3) SHAP-based explanation of operational alerts;
(4) evidence-restricted RAG reporting to reduce hallucination risk; and
(5) comparative evaluation of the integrated system against isolated
components in terms of detection performance, traceability, and
interpretation time. Peruvian Resolution SBS N° 053-2023 is used only as
a national reference for model risk management practices, while D.S. N°
115-2025-PCM is adopted as the general Peruvian AI governance framework.

\textbf{Keywords}: operational supervision, anomaly detection,
agro-export, AI explainability, language models, governance, GBDT,
automatic reports, traceability, artificial intelligence.

\begin{center}\rule{0.5\linewidth}{0.5pt}\end{center}

\begin{center}\rule{0.5\linewidth}{0.5pt}\end{center}

\section{ÍNDICE DE CONTENIDOS}\label{uxedndice-de-contenidos}

\begin{itemize}
\tightlist
\item
  DEDICATORIA
\item
  AGRADECIMIENTOS
\item
  PRESENTACIÓN
\item
  RESUMEN
\item
  ABSTRACT
\item
  ÍNDICE DE CONTENIDOS
\item
  ÍNDICE DE FIGURAS
\item
  ÍNDICE DE TABLAS
\item
  ÍNDICE DE FÓRMULAS
\item
  INTRODUCCIÓN
\item
  \textbf{CAPÍTULO I: PLANTEAMIENTO DEL PROBLEMA}

  \begin{itemize}
  \tightlist
  \item
    1.1 Descripción de la Realidad Problemática
  \item
    1.2 Problema Principal
  \item
    1.3 Objetivos
  \item
    1.4 Hipótesis de la Investigación
  \item
    1.5 Variables e Indicadores
  \item
    1.6 Viabilidad de la Investigación
  \item
    1.7 Justificación e Importancia
  \item
    1.8 Alcance
  \item
    1.9 Línea, Tipo y Nivel de la Investigación
  \item
    1.10 Técnicas e Instrumentos de Recolección de Información
  \end{itemize}
\item
  \textbf{CAPÍTULO II: MARCO TEÓRICO}

  \begin{itemize}
  \tightlist
  \item
    2.1 Antecedentes de la Investigación
  \item
    2.2 Estado del Arte
  \item
    2.3 Marco Conceptual
  \end{itemize}
\item
  \textbf{CAPÍTULO III: PROPUESTA METODOLÓGICA}

  \begin{itemize}
  \tightlist
  \item
    3.1 Arquitectura del Sistema Integrado
  \item
    3.2 Datasets de Validación y Benchmarks
  \item
    3.3 Configuración Experimental y Métricas
  \end{itemize}
\item
  \textbf{CAPÍTULO IV: RESULTADOS Y DISCUSIÓN}

  \begin{itemize}
  \tightlist
  \item
    4.1 Resultados Cuantitativos
  \item
    4.2 Resultados Cualitativos
  \item
    4.3 Discusión de Resultados
  \end{itemize}
\item
  \textbf{CAPÍTULO V: CONCLUSIONES Y TRABAJOS FUTUROS}

  \begin{itemize}
  \tightlist
  \item
    5.1 Conclusiones
  \item
    5.2 Limitaciones
  \item
    5.3 Trabajos Futuros
  \end{itemize}
\item
  CRONOGRAMA DE ACTIVIDADES
\item
  CONCLUSIONES (English)
\item
  RECOMENDACIONES
\item
  GLOSARIO DE TÉRMINOS
\item
  REFERENCIAS BIBLIOGRÁFICAS
\item
  ANEXOS
\end{itemize}

\begin{center}\rule{0.5\linewidth}{0.5pt}\end{center}

\section{ÍNDICE DE FIGURAS}\label{uxedndice-de-figuras}

\emph{(Por completar --- se insertarán las figuras del diagrama de
arquitectura del sistema, flujo del pipeline y resultados
experimentales)}

\begin{center}\rule{0.5\linewidth}{0.5pt}\end{center}

\section{ÍNDICE DE TABLAS}\label{uxedndice-de-tablas}

\begin{itemize}
\tightlist
\item
  Tabla 2.1 --- Comparativa de Sistemas de Supervisión con IA
\item
  Tabla 2.2 --- Resumen del Estado del Arte por Bloques Temáticos
\item
  Tabla 1.1 --- Variables e Indicadores
\item
  Tabla 1.2 --- Cronograma de Actividades
\item
  Tabla 1.3 --- Técnicas e Instrumentos de Recolección
\end{itemize}

\begin{center}\rule{0.5\linewidth}{0.5pt}\end{center}

\section{ÍNDICE DE FÓRMULAS}\label{uxedndice-de-fuxf3rmulas}

\begin{itemize}
\tightlist
\item
  Fórmula 1 --- Función objetivo GBDT:
  \(F^*(x) = \arg\min_F \mathbb{E}[L(y, F(x))]\)
\item
  Fórmula 2 --- Iteración GBDT:
  \(F_m(x) = F_{m-1}(x) + \nu \cdot h_m(x)\)
\item
  Fórmula 3 --- Local Outlier Factor: \(\text{LOF}_k(p)\)
\item
  Fórmula 4 --- Deep SVDD:
  \(\min_{W,R,c} R^2 + \frac{1}{\nu n}\sum \max(0, \|f(x_i;W)-c\|^2 - R^2)\)
\item
  Fórmula 5 --- Valor SHAP:
  \(\phi_i = \sum_{S \subseteq F \setminus \{i\}} \frac{|S|!(|F|-|S|-1)!}{|F|!}[f(S\cup\{i\})-f(S)]\)
\end{itemize}

\begin{center}\rule{0.5\linewidth}{0.5pt}\end{center}

\begin{center}\rule{0.5\linewidth}{0.5pt}\end{center}

\section{INTRODUCCIÓN}\label{introducciuxf3n}

La agroexportación peruana constituye un sector estratégico para la
economía nacional debido a su crecimiento sostenido, diversificación de
productos y participación en mercados internacionales exigentes. De
acuerdo con información oficial del Ministerio de Desarrollo Agrario y
Riego, al cierre de 2025 las agroexportaciones peruanas alcanzaron
ventas por USD 15 013 millones, con un crecimiento de 17.3\% respecto al
año anterior (MIDAGRI, 2026). Entre los principales productos exportados
destacaron arándanos, uvas, paltas, cacao y espárragos, lo que evidencia
la importancia económica y operativa de las cadenas agroexportadoras
peruanas.

En este contexto, las empresas agroexportadoras articulan procesos de
producción agrícola, acopio, almacenamiento, control de calidad,
cumplimiento fitosanitario, logística y comercialización internacional.
Cada uno de estos procesos genera datos que pueden revelar desviaciones
relevantes para la gestión operativa: cambios inusuales en precios,
variaciones de volumen, condiciones climáticas adversas, incumplimientos
de calidad, retrasos logísticos o patrones atípicos en el comportamiento
exportador. No obstante, la supervisión de estos procesos suele depender
de reportes manuales, hojas de cálculo, sistemas no integrados o
análisis posteriores a la ocurrencia del problema.

La inteligencia artificial ofrece herramientas adecuadas para abordar
esta brecha. Los modelos Gradient Boosting Decision Trees (GBDT) han
demostrado buen desempeño en datos tabulares estructurados (Grinsztajn
et al., 2022); los ensembles de detectores de anomalías permiten
identificar comportamientos atípicos de manera más robusta que un
detector individual (Han et al., 2022); la explicabilidad mediante
valores de Shapley (SHAP) convierte predicciones opacas en
justificaciones comprensibles (Lundberg \& Lee, 2017); y los modelos de
lenguaje con arquitectura RAG pueden transformar resultados
cuantitativos en reportes comprensibles siempre que se restrinja su
función a la generación narrativa basada en evidencias (Schneider et
al., 2025).

La presente investigación propone un sistema integrado de cuatro capas
que une predicción tabular, detección de anomalías, explicabilidad y
generación de reportes trazables en un flujo coherente de supervisión
operativa. El sistema se orienta al contexto agroexportador peruano y
busca mejorar la capacidad de detectar desviaciones operativas, explicar
sus posibles causas y documentar cada alerta de manera comprensible para
supervisores, responsables de calidad, gestores logísticos y auditores
internos. La Resolución SBS N° 053-2023 (SBS, 2023) se toma como
referencia nacional de buenas prácticas para gestión de riesgo de
modelos, sin asumirla como obligación directa para agroexportadoras; el
D.S. N° 115-2025-PCM (PCM, 2025) se emplea como marco peruano general
sobre gobernanza, transparencia y supervisión humana en inteligencia
artificial.

El documento se estructura de la siguiente manera: el Capítulo I plantea
el problema de investigación, define los objetivos, hipótesis, variables
e indicadores, y evalúa la viabilidad del proyecto. El Capítulo II
desarrolla el marco teórico, incluyendo antecedentes nacionales e
internacionales, estado del arte organizado en cinco debates
argumentativos, y el marco conceptual que fundamenta cada componente
técnico de la propuesta. El Capítulo III describe la arquitectura del
sistema propuesto, los datasets de validación y la configuración
experimental. El Capítulo IV presenta los resultados obtenidos y la
discusión de los hallazgos. Finalmente, el Capítulo V sintetiza las
conclusiones, limitaciones de la investigación y propuestas de trabajos
futuros.

\begin{center}\rule{0.5\linewidth}{0.5pt}\end{center}

\begin{center}\rule{0.5\linewidth}{0.5pt}\end{center}

\section{CAPÍTULO I: PLANTEAMIENTO DEL
PROBLEMA}\label{capuxedtulo-i-planteamiento-del-problema}

\subsection{1.1 Descripción de la Realidad
Problemática}\label{descripciuxf3n-de-la-realidad-problemuxe1tica}

\subsubsection{Contexto agroexportador y
empresarial}\label{contexto-agroexportador-y-empresarial}

En empresas agroexportadoras ---dedicadas a la producción, acopio,
procesamiento, empaque, control de calidad y comercialización
internacional de productos agrícolas--- la supervisión operativa
constituye una actividad crítica para detectar desviaciones productivas,
mermas, variaciones de precios, condiciones climáticas adversas,
retrasos logísticos e incumplimientos de estándares fitosanitarios.
Estas desviaciones afectan directamente la rentabilidad, continuidad
operativa y competitividad internacional de la empresa.

La magnitud económica del sector refuerza la necesidad de sistemas de
supervisión más oportunos y trazables. Según MIDAGRI, las
agroexportaciones peruanas superaron los USD 15 013 millones al cierre
de 2025, con crecimiento de 17.3\% respecto al año anterior (MIDAGRI,
2026). Este dinamismo incrementa la complejidad de las cadenas
agroexportadoras, que deben coordinar producción, calidad, sanidad,
logística y comercio exterior bajo condiciones cambiantes de clima,
demanda internacional y requisitos de destino.

En el contexto peruano, la transformación digital y la adopción de
sistemas inteligentes exigen mayores niveles de gobernanza tecnológica y
trazabilidad operativa. La Ley N.° 31814 y su reglamento aprobado
mediante D.S. N.° 115-2025-PCM (PCM, 2025) establecen un marco general
para promover el uso responsable de la inteligencia artificial, con
énfasis en transparencia, supervisión humana y gestión de riesgos.
Asimismo, la Resolución SBS N.° 053-2023 (SBS, 2023) se considera en
esta tesis como referencia nacional de buenas prácticas para gestión de
riesgo de modelos, validación y monitoreo, sin asumirla como obligación
directa para empresas agroexportadoras.

En escenarios donde una empresa agroexportadora integra datos de
producción, precios, clima, inventario, calidad, logística y
exportación, la supervisión manual resulta operativamente limitada. Sin
embargo, sistemas automatizados sin mecanismos de explicabilidad reducen
la confianza organizacional y dificultan la revisión interna de
decisiones. En consecuencia, surge la necesidad de sistemas integrados
capaces de detectar anomalías operativas en tiempo oportuno, explicar
sus causas probables y generar reportes trazables que permitan
justificar cada alerta dentro de la cadena agroexportadora.

\subsubsection{Problemas identificados}\label{problemas-identificados}

\begin{enumerate}
\def\labelenumi{\arabic{enumi}.}
\item
  \textbf{Falta de integración entre módulos operativos}: Los
  componentes de predicción, detección de anomalías y generación de
  reportes funcionan de forma independiente. Los hallazgos de un módulo
  no se comunican al siguiente con contexto semántico, lo que impide
  obtener una visión coherente del estado operativo de la empresa.
\item
  \textbf{Baja explicabilidad de las alertas automáticas}: Los modelos
  predictivos generan puntuaciones de riesgo sin justificar qué
  variables operativas (mermas, retrasos, desvíos de calidad)
  determinaron el resultado. Los gestores y auditores internos requieren
  explicaciones comprensibles para tomar decisiones correctivas.
\item
  \textbf{Reportería manual e ineficiente}: La detección de anomalías
  operativas ---en producción, inventario, logística o calidad--- exige
  la redacción manual de informes por parte de analistas, con riesgo de
  inconsistencia, errores de interpretación y demoras que comprometen la
  capacidad de respuesta empresarial.
\item
  \textbf{Falta de trazabilidad en la cadena de decisiones}: La
  gobernanza tecnológica y la auditoría interna requieren poder rastrear
  el fundamento de cada alerta desde el dato de origen (sensor, ERP,
  reporte de campo) hasta el informe ejecutivo. Los sistemas en silos no
  permiten esta trazabilidad end-to-end.
\item
  \textbf{Ausencia de validación cruzada entre dimensiones operativas}:
  Una anomalía detectada en producción no se correlaciona
  automáticamente con posibles irregularidades en inventario, logística
  o calidad. Esta falta de perspectiva multidimensional genera falsos
  positivos que saturan a los analistas y falsos negativos que pasan
  desapercibidos.
\item
  \textbf{Incapacidad de anticipar riesgos operativos}: Los sistemas
  reactivos detectan anomalías ya ocurridas, pero no predicen tendencias
  de riesgo emergentes. La ausencia de módulos predictivos integrados
  impide anticipar desviaciones estacionales, cuellos de botella
  logísticos o deterioro progresivo de indicadores de calidad.
\end{enumerate}

\subsection{1.2 Problema Principal}\label{problema-principal}

\textbf{¿Cómo mejorar la detección, explicación y documentación de
anomalías operativas en empresas agroexportadoras peruanas mediante un
sistema integrado de inteligencia artificial explicable que combine
predicción tabular, detección de anomalías, explicabilidad y generación
de reportes trazables?}

\subsubsection{Sub-problemas}\label{sub-problemas}

\begin{itemize}
\tightlist
\item
  ¿Qué variables operativas, climáticas, comerciales y fitosanitarias
  pueden utilizarse para caracterizar el comportamiento normal y anómalo
  de procesos agroexportadores peruanos?
\item
  ¿Qué arquitectura de inteligencia artificial permite integrar
  predicción tabular, detección de anomalías, explicabilidad y
  generación de reportes en un flujo operativo trazable?
\item
  ¿De qué manera la explicabilidad mediante SHAP contribuye a que
  supervisores operativos comprendan las causas probables de una alerta?
\item
  ¿Cómo generar reportes automáticos que sean comprensibles, accionables
  y trazables sin permitir que el modelo de lenguaje tome decisiones o
  invente información?
\item
  ¿Cómo evaluar si el sistema integrado mejora la trazabilidad,
  comprensión de alertas y tiempo de decisión frente al uso de
  componentes aislados?
\end{itemize}

\subsection{1.3 Objetivos}\label{objetivos}

\subsubsection{1.3.1 Objetivo Principal}\label{objetivo-principal}

Diseñar, implementar y evaluar un sistema integrado de supervisión
operativa basado en inteligencia artificial explicable para detectar
anomalías en datos agroexportadores, explicar los factores asociados
mediante SHAP y generar reportes trazables que apoyen la toma de
decisiones en empresas agroexportadoras peruanas.

\subsubsection{1.3.2 Objetivos
Específicos}\label{objetivos-especuxedficos}

\begin{enumerate}
\def\labelenumi{\arabic{enumi}.}
\item
  \textbf{Fuentes de datos y dominio}: Identificar y documentar fuentes
  de datos públicas y sintéticas aplicables a la supervisión operativa
  agroexportadora, considerando variables de precios, volúmenes, clima,
  comercio exterior y cumplimiento fitosanitario.
\item
  \textbf{Arquitectura y modularidad}: Diseñar una arquitectura modular
  de cuatro capas (predicción → detección → explicación → reporte) que
  integre modelos tabulares, detectores de anomalías, explicabilidad y
  generación de reportes trazables.
\item
  \textbf{Predicción y detección robusta}: Implementar modelos de
  predicción y detección de anomalías sobre datos agroexportadores
  públicos y sintéticos, utilizando algoritmos adecuados para datos
  tabulares y series temporales.
\item
  \textbf{Explicabilidad verificable}: Integrar SHAP (Lundberg \& Lee,
  2017) para identificar las variables que más contribuyen a cada alerta
  generada por el sistema.
\item
  \textbf{Generación de reportes trazables}: Diseñar un componente
  LLM+RAG que redacte explicaciones operativas basadas exclusivamente en
  evidencias estructuradas del sistema.
\item
  \textbf{Evaluación integrada}: Evaluar el sistema integrado mediante
  métricas técnicas, trazabilidad documental y prueba de
  comprensión/tiempo de decisión con usuarios o evaluadores simulados.
\end{enumerate}

\subsection{1.4 Hipótesis de la
Investigación}\label{hipuxf3tesis-de-la-investigaciuxf3n}

\textbf{Hipótesis General (H1)}: Un sistema integrado de predicción,
detección de anomalías, explicabilidad y generación de reportes
trazables mejora la trazabilidad de decisiones, la comprensión de
alertas y el tiempo de decisión de supervisores operativos frente al uso
de componentes aislados.

\textbf{Hipótesis Nula (H0)}: No existe diferencia significativa entre
el sistema integrado y los componentes aislados en trazabilidad de
decisiones, comprensión de alertas o tiempo de decisión de supervisores
operativos.

\textbf{Sub-hipótesis}:

\begin{itemize}
\tightlist
\item
  \textbf{H1a}: El uso combinado de modelos tabulares y detectores de
  anomalías permite identificar desviaciones operativas con mejor
  rendimiento que detectores individuales aplicados de forma aislada.
\item
  \textbf{H1b}: Las explicaciones SHAP incrementan la comprensión de las
  alertas por parte de supervisores operativos, al identificar variables
  relevantes y dirección de impacto.
\item
  \textbf{H1c}: Los reportes generados mediante RAG a partir de
  evidencias estructuradas presentan mayor trazabilidad y consistencia
  que reportes generados sin recuperación de contexto.
\item
  \textbf{H1d}: El sistema integrado reduce el tiempo requerido para
  interpretar una alerta operativa frente a un flujo basado en tablas,
  gráficos o salidas técnicas aisladas.
\end{itemize}

\subsection{1.5 Variables e Indicadores}\label{variables-e-indicadores}

\subsubsection{1.5.1 Variable
Independiente}\label{variable-independiente}

\textbf{Tipo de sistema de supervisión operativa (variable categórica)}:
- VI1: Sistema integrado (predicción tabular + detección de anomalías +
SHAP + LLM+RAG) - VI2: Componentes aislados (salidas técnicas
independientes por módulo)

\subsubsection{1.5.2 Variables
Dependientes}\label{variables-dependientes}

\textbf{VD1: Rendimiento de detección} - Indicadores: ROC-AUC,
Precisión, Recall, F1-Score, PR-AUC - Criterio de aceptación: superar el
baseline individual o justificar rendimiento equivalente con mayor
trazabilidad

\textbf{VD2: Calidad de explicabilidad} - Indicadores: cobertura top-k
SHAP, consistencia cualitativa, claridad de variables explicativas -
Criterio de aceptación: las variables principales deben permitir
explicar operativamente la alerta

\textbf{VD3: Calidad de reportes generados} - Indicadores: completitud,
consistencia, accionabilidad, coherencia textual y correspondencia con
evidencias - Criterio de aceptación: evaluación manual ≥ 4/5 en rúbrica
de reporte trazable

\textbf{VD4: Comprensión y tiempo de decisión del supervisor} -
Indicadores: tiempo-a-decisión (segundos), comprensión de alerta (Likert
1--5), decisión final correcta - Criterio de aceptación: reducción de
tiempo y mejora de comprensión respecto a componentes aislados

\textbf{VD5: Trazabilidad documental} - Indicadores: porcentaje de
alertas con dato, modelo, score, umbral, explicación, fuente recuperada
y reporte generado - Criterio de aceptación: ≥ 95\% de alertas con
campos de trazabilidad completos

\subsection{1.6 Viabilidad de la
Investigación}\label{viabilidad-de-la-investigaciuxf3n}

\subsubsection{1.6.1 Viabilidad Técnica}\label{viabilidad-tuxe9cnica}

\textbf{Disponibilidad de tecnologías}: El stack tecnológico es
completamente open-source y maduro: XGBoost (Chen \& Guestrin, 2016),
LightGBM (Ke et al., 2017) y CatBoost (Prokhorenkova et al., 2018) para
predicción tabular; PyOD (Zhao et al., 2019) para ensemble de anomalías
con acceso a Isolation Forest (Liu et al., 2008), LOF (Breunig et al.,
2000) y Deep SVDD (Ruff et al., 2018); SHAP (Lundberg \& Lee, 2017) para
explicabilidad; APIs de LLM (Anthropic Claude, OpenAI GPT-4) o modelos
locales (Llama 3) para generación de reportes.

\textbf{Datos disponibles}: Se contemplan tres niveles de datos. El
primer nivel corresponde a fuentes públicas oficiales: MIDAGRI para
agroexportaciones, precios y boletines sectoriales; SENAMHI para
variables climáticas; SENASA para requisitos fitosanitarios; SUNAT para
exportaciones; INEI para indicadores económicos; FAOSTAT y UN Comtrade
para validación internacional. El segundo nivel corresponde a un dataset
sintético agroexportador documentado, construido con variables
operativas plausibles y etiquetas de anomalía controladas. El tercer
nivel, opcional, corresponde a datos privados de una empresa
agroexportadora bajo acuerdo de confidencialidad. Como referencia
metodológica complementaria puede utilizarse el BAF Benchmark (Jesus et
al., 2022), no como evidencia directa del dominio agroexportador, sino
como benchmark tabular desbalanceado con drift temporal.

\textbf{Riesgos técnicos identificados}: La latencia de SHAP en datasets
grandes (\textgreater1M filas) puede mitigarse con los métodos de
aproximación TreeSHAP. La variabilidad en salidas de LLMs requiere
prompt engineering robusto y restricción mediante RAG. La mitigación
incluye pruebas piloto en subconjuntos de datos y benchmarking
iterativo.

\subsubsection{1.6.2 Viabilidad Operativa}\label{viabilidad-operativa}

\textbf{Timeline}: Fase 1 (meses 1--2): preparación de datos,
implementación de arquitectura base. Fase 2 (meses 2--3): entrenamiento
de modelos, validación experimental. Fase 3 (mes 4): test de usabilidad
con auditores voluntarios. Fase 4 (mes 5): análisis de resultados,
escritura y defensa.

\textbf{Presupuesto estimado}: Infraestructura GPU cloud y APIs LLM: USD
500--1,000. Stack open-source: USD 0. Incentivos para participantes del
test de usabilidad: USD 200--300. Total aproximado: USD 800--1,300.

\subsubsection{1.6.3 Viabilidad
Económica}\label{viabilidad-econuxf3mica}

La viabilidad económica se justifica por la relevancia del sector
agroexportador peruano y por el costo operativo asociado a decisiones
tardías ante desviaciones de precio, volumen, calidad, clima o
logística. MIDAGRI reportó agroexportaciones por USD 15 013 millones al
cierre de 2025 (MIDAGRI, 2026), por lo que incluso mejoras marginales en
detección temprana, trazabilidad y tiempo de respuesta pueden
representar beneficios operativos relevantes. En esta fase, los
beneficios económicos se tratarán como escenarios exploratorios y no
como resultados finales hasta contar con evaluación experimental y
supuestos documentados.

\subsection{1.7 Justificación e Importancia de la
Investigación}\label{justificaciuxf3n-e-importancia-de-la-investigaciuxf3n}

\subsubsection{1.7.1 Justificación
Teórica}\label{justificaciuxf3n-teuxf3rica}

La revisión sistemática de la literatura revela avances importantes en
modelos tabulares, detección de anomalías, explicabilidad y generación
de reportes mediante LLMs. Sin embargo, estos componentes suelen
estudiarse de forma aislada y en dominios distintos al agroexportador.
Trabajos de auditoría financiera o fraude contable, como AuditCopilot
(Kadir et al., 2025), se utilizan solo como antecedentes metodológicos
sobre automatización de reportes y detección de anomalías, no como eje
del dominio de aplicación. La brecha central de esta tesis es la
ausencia de una arquitectura integrada y trazable para supervisión
operativa agroexportadora peruana que combine fuentes públicas
oficiales, datos sintéticos documentados, predicción tabular, detección
de anomalías por ensemble, explicabilidad SHAP y reportes basados en
evidencia bajo restricción anti-alucinación.

\textbf{Aporte original específico}: Esta tesis constituye, en el
conocimiento del autor (verificado mediante búsqueda sistemática
documentada en \texttt{docs/busqueda-sistematica-gap.md}), la primera
arquitectura integrada de cuatro capas (GBDT + ensemble Isolation
Forest/LOF/ECOD + TreeSHAP + LLM-RAG con restricción anti-alucinación)
evaluada sobre datos agroexportadores peruanos públicos y sintéticos,
con trazabilidad documental diseñada conforme al D.S. N° 115-2025-PCM y
los principios del NIST AI RMF (NIST, 2023).

Esta tesis aporta a la literatura y a la práctica profesional cuatro
elementos diferenciados y verificables:

\begin{enumerate}
\def\labelenumi{\arabic{enumi}.}
\item
  \textbf{Aporte arquitectónico}: Un modelo conceptual integrado de
  cuatro capas para supervisión operativa agroexportadora, donde la
  separación estricta de responsabilidades entre detección determinista
  (GBDT + ensemble) y narración asistida (LLM-RAG anclado en vectores
  SHAP) constituye un patrón de diseño anti-alucinación replicable en
  otros dominios regulados.
\item
  \textbf{Aporte de datos abiertos}: Un protocolo público y reproducible
  de construcción y documentación de un dataset sintético
  agroexportador, calibrado con rangos plausibles de
  MIDAGRI/SENAMHI/SENASA y descrito según Datasheets for Datasets (Gebru
  et al., 2021), disponible para la comunidad como referencia
  metodológica.
\item
  \textbf{Aporte de evaluación}: Una metodología de evaluación dual que
  combina métricas técnicas (PR-AUC, F1, ROC-AUC) con métricas de
  utilidad operativa (tiempo-a-decisión, comprensión Likert,
  trazabilidad documental) y aplica pruebas estadísticas formales
  (Wilcoxon signed-rank, t-Student apareado) para contrastar las
  sub-hipótesis H1a--H1d.
\item
  \textbf{Aporte regulatorio}: Una primera traducción operativa de los
  principios del D.S. N° 115-2025-PCM al diseño de un sistema de IA
  empresarial peruano, mostrando cómo cada capa de la arquitectura puede
  mapear con requisitos de explicabilidad, supervisión humana y
  trazabilidad documental.
\end{enumerate}

\subsubsection{1.7.2 Justificación
Económica}\label{justificaciuxf3n-econuxf3mica}

La automatización inteligente de la supervisión operativa tiene impacto
económico potencial al reducir el tiempo de análisis, mejorar la
detección temprana de desviaciones y facilitar la documentación de
decisiones. En empresas agroexportadoras, las alertas oportunas sobre
precios, volúmenes, clima, mermas, calidad o logística pueden apoyar
decisiones correctivas antes de que la desviación se convierta en
pérdida operativa o incumplimiento comercial. La escalabilidad del
sistema permite adaptarlo a empresas de distintos tamaños mediante
fuentes públicas, datos internos o datos sintéticos documentados.

\subsubsection{1.7.3 Justificación
Social}\label{justificaciuxf3n-social}

El Decreto Supremo N° 115-2025-PCM, reglamento de la Ley N° 31814 de
Inteligencia Artificial del Perú, proporciona un marco nacional para
promover el uso responsable de la IA, incluyendo transparencia,
supervisión humana y gestión de riesgos (PCM, 2025). A nivel
internacional, el Reglamento (UE) 2024/1689 ---EU AI Act--- refuerza la
importancia de documentar sistemas de IA, especialmente cuando sus
resultados afectan decisiones relevantes (Parlamento Europeo y Consejo,
2024). El sistema propuesto incorpora estos principios mediante
explicabilidad, trazabilidad documental y revisión humana de reportes.

Adicionalmente, una detección temprana de anomalías operativas protege a
las empresas agroexportadoras de tamaño medio ---que representan la
mayoría del sector exportador peruano--- frente a pérdidas acumuladas
por mermas, incumplimientos de calidad y fallas logísticas que, sin
sistemas de alerta temprana, solo se visibilizan en los estados
financieros al cierre del período.

\subsubsection{1.7.4 Importancia}\label{importancia}

\textbf{Nivel académico}: Contribución a los campos de ML interpretable,
detección de anomalías, supervisión operativa y gobernanza de IA
aplicada a cadenas agroexportadoras.

\textbf{Nivel profesional}: Guía de referencia para empresas
agroexportadoras que busquen incorporar IA explicable en procesos de
control operativo, calidad, logística y toma de decisiones.

\textbf{Nivel institucional}: Fortalece el posicionamiento de la UNSA en
investigación aplicada en IA responsable y establece vínculos de
colaboración con el sector financiero regional.

\subsection{1.8 Alcance}\label{alcance}

\textbf{Alcance temático}: Predicción con GBDT, detección de anomalías
mediante ensemble, explicabilidad SHAP, generación de reportes con
LLM+RAG, documentación de datasets y trazabilidad de alertas operativas.
\textbf{Excluye}: modelos de Deep Learning puro para datos tabulares
como propuesta principal, implementación productiva en tiempo real,
reemplazo de la decisión humana y análisis legal profundo de regulación
sectorial.

\textbf{Alcance geográfico}: Empresas agroexportadoras peruanas, con
énfasis en productos representativos como arándanos, uvas, paltas, cacao
y espárragos. La investigación utiliza fuentes públicas oficiales y
dataset sintético documentado; los datos privados de empresa quedan como
extensión futura u opcional.

\textbf{Alcance temporal}: Evaluación en dataset estático o
semiestático, sin monitoreo en producción. Estudio en 5 meses.
Evaluación de comprensión y tiempo de decisión con supervisores,
auditores internos o evaluadores simulados según disponibilidad.

\subsection{1.9 Línea, Tipo y Nivel de la
Investigación}\label{luxednea-tipo-y-nivel-de-la-investigaciuxf3n}

\subsubsection{1.9.1 Línea de
Investigación}\label{luxednea-de-investigaciuxf3n}

La presente investigación se enmarca en la línea de \textbf{Inteligencia
Artificial e Ingeniería de Software Aplicada} de la Escuela Profesional
de Ingeniería de Sistemas de la Universidad Nacional de San Agustín de
Arequipa. Específicamente, se inscribe en el área de sistemas
inteligentes para la toma de decisiones en organizaciones empresariales,
con énfasis en gobernanza tecnológica y conformidad regulatoria.

\subsubsection{1.9.2 Marco
Epistemológico}\label{marco-epistemoluxf3gico}

La investigación adopta un enfoque \textbf{post-positivista} (Creswell
\& Creswell, 2018): asume que los fenómenos operativos del dominio
agroexportador (precios, volúmenes, mermas, condiciones climáticas,
cumplimiento fitosanitario, tiempos logísticos) son medibles de manera
objetiva mediante variables cuantitativas y rangos plausibles
documentables. Al mismo tiempo, se reconoce que la evaluación de
utilidad operativa de un sistema de supervisión asistida por IA
incorpora componentes subjetivos ---comprensión, confianza,
accionabilidad del reporte--- que requieren triangulación entre métricas
cuantitativas (PR-AUC, tiempo-a-decisión) y métricas cualitativas
(escalas Likert, rúbricas de evaluación). Esta postura justifica la
combinación metodológica utilizada y excluye afirmaciones de causalidad
en sentido experimental estricto, sustituyéndolas por afirmaciones sobre
\textbf{efecto diferencial} entre condiciones experimentales controladas
(sistema integrado vs.~sistema con componentes aislados).

\subsubsection{1.9.3 Tipo de
Investigación}\label{tipo-de-investigaciuxf3n}

La investigación es de tipo \textbf{aplicada}, dado que utiliza
conocimiento teórico y metodológico existente en machine learning,
detección de anomalías, explicabilidad algorítmica y modelos de lenguaje
para diseñar, implementar y evaluar un sistema concreto que resuelve un
problema identificado en el contexto empresarial agroexportador peruano.
No se busca generar nuevos algoritmos de base, sino integrar y validar
una arquitectura que produce valor operativo y regulatorio verificable.

\subsubsection{1.9.4 Nivel de
Investigación}\label{nivel-de-investigaciuxf3n}

El nivel de la investigación es \textbf{explicativo-evaluativo}. Se
parte de un diagnóstico descriptivo del problema (sistemas de
supervisión fragmentados y sin trazabilidad), se propone una solución
arquitectónica basada en literatura existente, y se diseñan experimentos
controlados que permiten evaluar la relación entre integración de
componentes y mejoras en trazabilidad, comprensión operativa y tiempo de
decisión. La evaluación combina métricas cuantitativas (PR-AUC,
F1-Score, cobertura de trazabilidad, tiempo-a-decisión) y cualitativas
(escala Likert con supervisores o evaluadores). Las conclusiones
explicativas se derivan del análisis ablativo (Experimento E5) que aísla
la contribución de cada capa al rendimiento global, distinguiendo así
qué componente arquitectónico es responsable de cada mejora observada.

\subsubsection{1.9.5 Diseño de la
Investigación}\label{diseuxf1o-de-la-investigaciuxf3n}

El diseño es \textbf{cuasi-experimental con grupos contrabalanceados}
para la evaluación de utilidad operativa (VD4): cada participante del
estudio de usabilidad evalúa ambas condiciones (sistema integrado y
componentes aislados) en orden aleatorizado, controlando por el efecto
de aprendizaje mediante diseño within-subjects. Para las variables
técnicas (VD1, VD2, VD3, VD5) el diseño es \textbf{comparativo en
condiciones controladas}: el mismo dataset, mismas particiones
train/test temporales (sin solapamiento) y misma semilla aleatoria se
aplican a todas las configuraciones experimentales (E1--E5), variando
únicamente la condición evaluada (detector individual vs.~ensemble, sin
SHAP vs.~con SHAP, sin RAG vs.~con RAG, sin pipeline integrado vs.~con
pipeline integrado).

\subsection{1.10 Técnicas e Instrumentos de Recolección de
Información}\label{tuxe9cnicas-e-instrumentos-de-recolecciuxf3n-de-informaciuxf3n}

\subsubsection{1.10.1 Técnicas}\label{tuxe9cnicas}

Las técnicas de recolección de información utilizadas en esta
investigación son:

\begin{enumerate}
\def\labelenumi{\arabic{enumi}.}
\item
  \textbf{Revisión sistemática de literatura}: Búsqueda estructurada en
  bases de datos académicas (IEEE Xplore, ACM Digital Library, arXiv,
  Google Scholar, Scopus) utilizando términos de búsqueda como ``anomaly
  detection ensemble'', ``GBDT tabular data'', ``SHAP explainability'',
  ``RAG generated reports'', ``AI governance Peru'' y ``agricultural
  anomaly detection''. Se aplican criterios de inclusión: publicaciones
  entre 2017 y 2026, en inglés o español, con evaluación empírica o
  propuesta metodológica verificable.
\item
  \textbf{Análisis documental}: Revisión de fuentes oficiales nacionales
  (MIDAGRI, SENASA, SENAMHI, SUNAT, INEI), regulaciones nacionales sobre
  IA (Ley N° 31814 y D.S. N° 115-2025-PCM) y marcos internacionales de
  gobernanza (EU AI Act 2024, NIST AI RMF 1.0). La Resolución SBS N°
  053-2023 se considera solo como referencia metodológica de gestión de
  riesgo de modelos.
\item
  \textbf{Experimentación controlada}: Diseño de experimentos
  comparativos con el sistema integrado versus componentes aislados,
  evaluados sobre datos públicos/sintéticos agroexportadores. El BAF
  Benchmark (Jesus et al., 2022) puede emplearse solo como benchmark
  metodológico complementario para datos tabulares desbalanceados.
\item
  \textbf{Evaluación con usuarios o evaluadores simulados}: Prueba de
  comprensión y tiempo de decisión con supervisores, auditores internos,
  estudiantes avanzados o evaluadores simulados mediante protocolo de
  tareas cronometradas y cuestionario post-tarea.
\end{enumerate}

\subsubsection{1.10.2 Instrumentos}\label{instrumentos}

\begin{longtable}[]{@{}
  >{\raggedright\arraybackslash}p{(\columnwidth - 4\tabcolsep) * \real{0.3333}}
  >{\raggedright\arraybackslash}p{(\columnwidth - 4\tabcolsep) * \real{0.3333}}
  >{\raggedright\arraybackslash}p{(\columnwidth - 4\tabcolsep) * \real{0.3333}}@{}}
\toprule\noalign{}
\begin{minipage}[b]{\linewidth}\raggedright
Instrumento
\end{minipage} & \begin{minipage}[b]{\linewidth}\raggedright
Propósito
\end{minipage} & \begin{minipage}[b]{\linewidth}\raggedright
Variable que mide
\end{minipage} \\
\midrule\noalign{}
\endhead
\bottomrule\noalign{}
\endlastfoot
Fuentes públicas oficiales & Construir contexto, variables y rangos
plausibles del dominio agroexportador & VD1, VD5 \\
Dataset sintético agroexportador & Evaluación controlada de predicción,
anomalías y reportes & VD1, VD2, VD3 \\
BAF Benchmark (Jesus et al., 2022) & Benchmark metodológico
complementario para datos tabulares desbalanceados & VD1 \\
Cuestionario de comprensión (Likert 1--5) & Evaluar claridad de
explicaciones y reportes para supervisores & VD2, VD3, VD4 \\
Registro de tiempo-a-decisión & Medir segundos requeridos para
interpretar y decidir sobre una alerta & VD4 \\
Checklist de trazabilidad & Verificar dato, modelo, score, umbral,
explicación, fuente y reporte por alerta & VD5 \\
\end{longtable}

\subsection{1.11 Cronograma de
Actividades}\label{cronograma-de-actividades}

\begin{longtable}[]{@{}lccccc@{}}
\toprule\noalign{}
Actividad & Mes 1 & Mes 2 & Mes 3 & Mes 4 & Mes 5 \\
\midrule\noalign{}
\endhead
\bottomrule\noalign{}
\endlastfoot
Revisión bibliográfica y marco teórico & ✓ & & & & \\
Obtención y preprocesamiento de datos & ✓ & ✓ & & & \\
Implementación Capa 1 (GBDT + TFT) & & ✓ & & & \\
Implementación Capa 2 (Ensemble anomalías) & & ✓ & ✓ & & \\
Implementación Capa 3 (SHAP) & & & ✓ & & \\
Implementación Capa 4 (LLM+RAG) & & & ✓ & & \\
Integración del pipeline completo & & & ✓ & & \\
Experimentos comparativos (baselines) & & & & ✓ & \\
Prueba de usabilidad con auditores & & & & ✓ & \\
Análisis de resultados y discusión & & & & ✓ & \\
Redacción del documento final & & & & & ✓ \\
Revisión del asesor y correcciones & & & & & ✓ \\
Presentación y defensa & & & & & ✓ \\
\end{longtable}

\subsection{1.12 Limitaciones de la
Investigación}\label{limitaciones-de-la-investigaciuxf3n}

La presente investigación adopta un enfoque transparente respecto a sus
limitaciones, declarando de forma explícita las amenazas a la validez
que podrían afectar la interpretación o generalización de los
resultados. Esta declaración es un requisito de rigor académico y
constituye una práctica estándar en la literatura de IA responsable.

\subsubsection{1.12.1 Limitaciones de validez
externa}\label{limitaciones-de-validez-externa}

\textbf{Dataset sintético}: La evaluación principal se realiza sobre un
dataset sintético agroexportador documentado con criterios de Datasheets
for Datasets (Gebru et al., 2021). Aunque las distribuciones y
mecanismos de inyección de anomalías se calibran con rangos plausibles
tomados de fuentes oficiales (MIDAGRI, SENAMHI, SENASA, SUNAT), un
dataset sintético no reproduce todas las dinámicas operativas, sesgos de
medición y patrones de drift presentes en datos reales de una empresa
agroexportadora. Por lo tanto, los resultados deben interpretarse como
evidencia de viabilidad arquitectónica y comportamiento esperado bajo
condiciones controladas, no como predictores directos del rendimiento en
producción.

\textbf{Generalización a otros sectores}: La arquitectura se diseña y
evalúa para el dominio agroexportador peruano. Su extensión a otros
sectores (minería, manufactura, retail) requiere recalibración de
variables, ajustes en el dataset y nueva validación experimental antes
de afirmar transferibilidad.

\textbf{Generalización a otros países}: El marco regulatorio que orienta
el diseño (D.S. N° 115-2025-PCM, Ley N° 31814) corresponde al contexto
peruano. Aunque los principios subyacentes (transparencia, supervisión
humana, gestión de riesgos) son compatibles con marcos internacionales
como el EU AI Act y el NIST AI RMF, la aplicabilidad directa del sistema
en otras jurisdicciones requiere una revisión de conformidad específica.

\subsubsection{1.12.2 Limitaciones de validez
interna}\label{limitaciones-de-validez-interna}

\textbf{Tamaño de muestra del estudio de usabilidad}: La evaluación de
VD4 (comprensión y tiempo de decisión) requiere participantes humanos
con perfil de supervisor operativo, analista de calidad o auditor
interno. La disponibilidad realista en el contexto de una tesis
individual es de 15 a 20 participantes mediante un diseño within-subject
(cada participante evalúa ambas condiciones en orden aleatorizado). Con
este tamaño, los resultados de VD4 deben reportarse como exploratorios,
calculando intervalos de confianza y tamaño de efecto (Cohen's d) en
lugar de afirmar significancia estadística con potencia plena.

\textbf{Sesgo del evaluador}: Los participantes en el estudio de
usabilidad pueden conocer la procedencia del sistema integrado, lo cual
puede introducir un sesgo de expectativa. Se mitiga mediante: (a) orden
contrabalanceado de presentación de las condiciones, (b) cuestionario
post-tarea con preguntas no directivas, y (c) registro automático de
tiempo-a-decisión (no dependiente de auto-reporte).

\textbf{Variabilidad del LLM}: El módulo de generación de reportes
utiliza un modelo de lenguaje cuyas respuestas presentan variación
estocástica entre ejecuciones, incluso con el mismo prompt. Para mitigar
este efecto se fija el parámetro de temperatura (temperature = 0.2), se
documenta la versión exacta del modelo utilizado y se reportan los
resultados como promedio sobre al menos 3 generaciones por alerta.

\subsubsection{1.12.3 Limitaciones de validez de
constructo}\label{limitaciones-de-validez-de-constructo}

\textbf{Definición de ``anomalía operativa''}: La etiqueta
\texttt{etiqueta\_anomalia} del dataset sintético se construye mediante
un protocolo de inyección controlada, lo que define la anomalía
operativamente a partir de reglas predeterminadas. En la práctica
empresarial, la categorización de un registro como ``anómalo'' depende
de criterios contextuales no siempre formalizables. El sistema, por
tanto, evalúa su capacidad de detectar desviaciones según reglas
declaradas, no como reemplazo del juicio experto.

\textbf{Métrica de comprensión (VD2)}: La ``claridad de variables
explicativas'' se mide mediante una escala Likert 1--5 que captura
percepción subjetiva del evaluador. Esta métrica está expuesta a sesgos
de halo y aquiescencia. Se complementa con métricas objetivas (cobertura
top-k SHAP) para triangulación.

\subsubsection{1.12.4 Limitaciones técnicas y de
recursos}\label{limitaciones-tuxe9cnicas-y-de-recursos}

\textbf{Dependencia de APIs comerciales de LLM}: La generación de
reportes puede emplear servicios comerciales (Anthropic Claude, OpenAI
GPT-4) cuya versión, costo y disponibilidad pueden variar durante el
horizonte experimental. Se documenta la versión exacta utilizada en cada
experimento y se evalúa la posibilidad de réplica con modelos locales
(Llama 3) como verificación cruzada.

\textbf{Recursos computacionales}: La investigación se ejecuta en
infraestructura GPU cloud limitada por presupuesto académico. Esto
restringe la exploración exhaustiva de hiperparámetros y limita el
tamaño del dataset experimental a un rango medio (2,000--5,000
registros). La escalabilidad a millones de registros queda como trabajo
futuro.

\subsection{1.13 Declaración de Intereses y Aspectos
Éticos}\label{declaraciuxf3n-de-intereses-y-aspectos-uxe9ticos}

Se declara que el autor de esta investigación no mantiene relación
contractual, comercial o financiera con empresas agroexportadoras
específicas que pudiera condicionar la independencia metodológica de los
resultados. La investigación se desarrolla en el marco de la formación
de pregrado en la Universidad Nacional de San Agustín de Arequipa y no
cuenta con financiamiento externo.

El uso de datos en esta tesis se restringe a fuentes públicas oficiales
(MIDAGRI, SENAMHI, SENASA, SUNAT, INEI, FAOSTAT, UN Comtrade) y a un
dataset sintético generado por el autor con un protocolo documentado. No
se utilizan datos personales, datos de empresas bajo confidencialidad ni
datos que pudieran exponer información sensible de actores del sector.

El estudio de usabilidad con participantes humanos (VD4) sigue los
principios de consentimiento informado, voluntariedad, anonimato y
derecho de retiro establecidos en la Declaración de Helsinki y adaptados
al contexto de investigación en ingeniería. El protocolo completo,
formulario de consentimiento y cuestionario figuran en el Anexo A.

El uso de herramientas de IA durante la elaboración de la tesis
(búsqueda bibliográfica, revisión de redacción, apoyo en codificación)
se documenta en el Anexo D, conforme a las prácticas emergentes de
transparencia académica sobre el uso de IA generativa en investigación.

\begin{center}\rule{0.5\linewidth}{0.5pt}\end{center}

\begin{center}\rule{0.5\linewidth}{0.5pt}\end{center}

\section{CAPÍTULO II: MARCO
TEÓRICO}\label{capuxedtulo-ii-marco-teuxf3rico}

\subsection{2.1 Antecedentes de la
Investigación}\label{antecedentes-de-la-investigaciuxf3n}

Los antecedentes de esta investigación se organizan desde dos
perspectivas. La primera corresponde a trabajos técnicos sobre datos
tabulares, detección de anomalías, explicabilidad y generación de
reportes, incluso cuando fueron desarrollados en dominios financieros o
contables. Estos trabajos se utilizan como soporte metodológico. La
segunda corresponde al dominio de aplicación de esta tesis: supervisión
operativa agroexportadora, donde la contribución principal consiste en
adaptar e integrar dichas técnicas para detectar, explicar y documentar
desviaciones en procesos agroexportadores peruanos.

El desarrollo de sistemas de predicción, detección de anomalías y
generación de reportes en datos empresariales ha seguido una trayectoria
de especialización creciente, marcada por tres tendencias paralelas: el
auge de los modelos basados en árboles para datos tabulares, la
proliferación de benchmarks sistemáticos y la emergencia de los modelos
de lenguaje como capa de interpretación. Los siguientes antecedentes
fueron seleccionados por su proximidad metodológica con la propuesta de
esta investigación, aunque varios provienen de dominios financieros o
contables y se emplean aquí solo como soporte técnico transferible.

\subsubsection{2.1.1 Kadir et al.~(2025) --- AuditCopilot: LLMs para
Reportes de
Anomalías}\label{kadir-et-al.-2025-auditcopilot-llms-para-reportes-de-anomaluxedas}

Kadir et al.~(2025) desarrollaron AuditCopilot (Kadir et al., 2025), un
sistema de auditoría contable que integra LLMs con detección de
anomalías en asientos de doble entrada para generar explicaciones
automáticas en lenguaje natural. El sistema implementa un pipeline de
tres etapas ---detección de irregularidades, interpretación contextual
con LLM ajustado y generación de narrativas--- evaluado sobre un corpus
de asientos contables sintéticos y reales. Los resultados reportan
mejoras en la tasa de detección y reducción del tiempo de revisión, con
valoración positiva de auditores en pruebas de aceptabilidad.

La relevancia de este antecedente para la presente tesis es
metodológica: confirma la viabilidad de combinar detección de anomalías
con generación de reportes LLM. No obstante, su dominio es contable, por
lo que no se adopta como evidencia agroexportadora. Esta tesis traslada
el principio de generación narrativa a un contexto operativo, separando
estrictamente la detección de la redacción mediante RAG sobre scores y
vectores SHAP.

\subsubsection{2.1.2 Park (2024) --- Framework Multi-Agente LLM para
Anomalías
Financieras}\label{park-2024-framework-multi-agente-llm-para-anomaluxedas-financieras}

Park (2024) propuso un framework de múltiples agentes LLM especializados
para validar alertas de anomalías en el mercado bursátil (S\&P 500)
(Park, 2024). La arquitectura organiza cuatro agentes ---conversión de
datos, análisis estadístico, verificación cruzada y consolidación--- que
se comunican mediante prompts estructurados y alcanzan mejores tasas de
verdaderos positivos que un LLM único. La especialización de agentes
demuestra ser superior a la generalización en la validación de señales
financieras.

Este trabajo aporta a la literatura evidencia de que los LLMs en
arquitecturas especializadas pueden mejorar la calidad del análisis
automatizado. Sin embargo, opera en mercados de alta frecuencia, un
dominio alejado del contexto agroexportador. Esta tesis aplica
únicamente el principio de especialización de roles ---LLM como
intérprete, no como detector--- en un sistema de supervisión operativa
con trazabilidad y generación restringida a datos verificados (Schneider
et al., 2025).

\subsubsection{2.1.3 Autores varios (2025) --- Ensemble GBDT+SHAP en
Datos Tabulares
Críticos}\label{autores-varios-2025-ensemble-gbdtshap-en-datos-tabulares-cruxedticos}

En el trabajo publicado en el \emph{Journal of Risk and Financial
Management} (2025), los autores diseñaron un framework integrado de
detección de fraude en estados financieros combinando Stacking Ensemble
de XGBoost, LightGBM y CatBoost con explicabilidad SHAP (JRFM, 2025). El
ensemble alcanza PR-AUC = 0.93 y F1-Score = 0.83, superando a TabNet y
FT-Transformer, con un SHAP Stability Index = 0.87 que certifica la
coherencia forense de las explicaciones ---requisito indispensable en
auditoría.

Este antecedente respalda la decisión arquitectónica de combinar GBDT y
SHAP en datos tabulares críticos. La diferencia clave es que dicho
trabajo se limita a detección de fraude en estados financieros; la
presente investigación adapta la lógica de modelos tabulares explicables
al contexto agroexportador, incorporando fuentes públicas, dataset
sintético documentado y generación de reportes LLM+RAG para supervisión
operativa.

\subsubsection{2.1.4 Han et al.~(2022) --- ADBench: Benchmark para
Detección de
Anomalías}\label{han-et-al.-2022-adbench-benchmark-para-detecciuxf3n-de-anomaluxedas}

Han et al.~(2022) publicaron ADBench (Han et al., 2022), un benchmark
sistemático que evalúa 30 algoritmos de detección de anomalías en 57
datasets reales y sintéticos bajo tres niveles de supervisión ---no
supervisado, semisupervisado y supervisado. El hallazgo central es que
no existe un algoritmo universalmente superior: el rendimiento depende
del tipo de anomalía, la distribución de datos y el nivel de etiquetado.
Isolation Forest y ECOD muestran consistencia en escenarios no
supervisados, y los ensembles de múltiples detectores superan
sistemáticamente a los detectores individuales en escenarios de alta
variabilidad distribucional.

ADBench justifica formalmente la estrategia de ensemble adoptada en esta
tesis y proporciona la metodología experimental de referencia para el
Capítulo III. La librería PyOD (Zhao et al., 2019), compatible con todos
los algoritmos evaluados en ADBench, asegura la reproducibilidad directa
de los resultados.

\subsubsection{2.1.5 Grinsztajn et al.~(2022) --- GBDT vs.~Deep Learning
en Datos
Tabulares}\label{grinsztajn-et-al.-2022-gbdt-vs.-deep-learning-en-datos-tabulares}

Grinsztajn et al.~(2022) realizaron un benchmark sistemático en 45
datasets tabulares comparando GBDT contra FT-Transformer, TabNet y MLP
(Grinsztajn et al., 2022). El resultado es contundente: en datasets con
menos de 50,000 muestras, los GBDT superan a cualquier modelo de Deep
Learning en el 95\% de los casos. Los autores identifican tres
propiedades estructurales de los datos tabulares que favorecen a los
árboles: robustez ante features no informativas, orientación no
invariante a rotaciones e irregularidades en la función objetivo.

Este trabajo cierra el debate GBDT versus Deep Learning para el tamaño
de dataset típico en entornos empresariales medianos y justifica de
manera irrefutable la elección de XGBoost y LightGBM como backbone del
módulo de predicción de esta tesis. Es el argumento bibliográfico
central de la primera batalla del estado del arte (§2.2.1).

\subsubsection{2.1.6 Lim et al.~(2021) --- Temporal Fusion Transformer
para Forecasting
Interpretable}\label{lim-et-al.-2021-temporal-fusion-transformer-para-forecasting-interpretable}

Lim et al.~(2021) propusieron el Temporal Fusion Transformer (TFT) (Lim
et al., 2021), una arquitectura que combina codificación LSTM, selección
de variables mediante mecanismo de gating, atención multi-cabezal
interpretable y predicción por cuantiles para forecasting
multi-horizonte con covariables exógenas. TFT supera a LSTMs, N-BEATS y
Transformers vanilla en cuatro de seis datasets de referencia, con mapas
de atención legibles por analistas de negocio.

TFT es seleccionado como arquitectura del módulo de forecasting por su
doble ventaja: rendimiento predictivo e interpretabilidad incorporada.
En el contexto de supervisión operativa, la capacidad de justificar qué
períodos temporales y qué covariables (indicadores de producción,
calendarios logísticos) fundamentan la predicción es un requisito
funcional equivalente en importancia a la precisión numérica.

\subsubsection{2.1.7 Zhao et al.~(2019) --- PyOD: Librería Estándar para
Detección de
Outliers}\label{zhao-et-al.-2019-pyod-libreruxeda-estuxe1ndar-para-detecciuxf3n-de-outliers}

Zhao et al.~(2019) desarrollaron PyOD (Zhao et al., 2019), una librería
unificada en Python que implementa más de 40 algoritmos de detección de
outliers con una API compatible con scikit-learn. Cubre métodos basados
en proximidad (LOF), proyección (PCA), ensembles (Isolation Forest) y
redes neuronales (Deep SVDD, AutoEncoder). Con más de 7,000 estrellas en
GitHub y adopción en publicaciones de NeurIPS, ICDM e ICML, PyOD es la
infraestructura técnica de referencia para implementar el ensemble de
detección de anomalías de esta tesis, garantizando reproducibilidad
directa con los 30 algoritmos de ADBench (Han et al., 2022).

\begin{center}\rule{0.5\linewidth}{0.5pt}\end{center}

\begin{center}\rule{0.5\linewidth}{0.5pt}\end{center}

\subsection{2.2 Estado del Arte}\label{estado-del-arte}

El estado del arte se organiza en torno a cinco debates fundamentales de
la literatura que la presente propuesta debe resolver o posicionarse
explícitamente. Cada sub-sección presenta el debate, los trabajos
relevantes y la posición de esta tesis. La Tabla 2.1 sintetiza todas las
referencias relevantes al final de la sección.

\subsubsection{2.2.1 GBDT versus Deep Learning para Datos Tabulares
Empresariales y
Agroexportadores}\label{gbdt-versus-deep-learning-para-datos-tabulares-empresariales-y-agroexportadores}

El desarrollo de modelos para datos tabulares ha seguido una trayectoria
diferente a la de visión computacional y procesamiento de lenguaje
natural: el Deep Learning no ha conseguido desplazar a los modelos
basados en árboles como estándar de facto en datos estructurados. Chen y
Guestrin (2016) introdujeron XGBoost como sistema escalable de gradient
boosting con regularización L1/L2, manejo nativo de valores faltantes y
paralelización por columnas, estableciéndolo como el baseline universal
con más de 45,000 citas en la literatura científica. Ke et al.~(2017) lo
extendieron con LightGBM, que incorpora Gradient-based One-Side Sampling
(GOSS) e histogramas para lograr velocidades de entrenamiento hasta 20
veces superiores con rendimiento comparable. Prokhorenkova et al.~(2018)
resolvieron el problema de target leakage en variables categóricas con
Ordered Boosting, siendo especialmente relevante en datos contables con
alta cardinalidad (cuentas, departamentos, centros de costo).

El auge del Deep Learning motivó intentos de adaptar estas arquitecturas
a datos tabulares. Gorishniy et al.~(2021) propusieron FT-Transformer,
el primer Transformer robusto para tablas mediante feature embeddings,
que en algunos benchmarks iguala pero raramente supera a los GBDT. Arik
y Pfister (2021) desarrollaron TabNet, que combina selección secuencial
de features con atención interpretable, argumentando que puede ofrecer
tanto rendimiento como interpretabilidad en un solo modelo. Sin embargo,
el estudio seminal de Grinsztajn et al.~(2022) zanjó empíricamente este
debate: en datasets tabulares de hasta 50,000 muestras, los GBDT superan
a los modelos de Deep Learning en la inmensa mayoría de los escenarios
empíricos. Esta robusta evidencia respalda la decisión arquitectónica de
la presente tesis, cuyo conjunto de datos transaccionales se sitúa en un
volumen óptimo (hasta 10,000 registros sintéticos) para maximizar el
rendimiento de los árboles de decisión. Los autores identifican tres
propiedades estructurales de los datos tabulares que favorecen a los
árboles: robustez ante features no informativas, orientación no
invariante a rotaciones y presencia de irregularidades en la función
objetivo ---todas características presentes en los registros
transaccionales de auditoría.

En dominios empresariales con datos tabulares heterogéneos, esta
evidencia respalda el uso de GBDT como primera opción antes de recurrir
a arquitecturas neuronales complejas. En el caso agroexportador, los
registros combinan variables numéricas, categóricas, temporales y
contextuales ---producto, zona, volumen, precio, clima, destino,
cumplimiento y logística---, por lo que los modelos basados en árboles
son una base técnica adecuada para capturar relaciones no lineales y
manejar variables de distinta naturaleza.

\textbf{Posición de esta tesis}: XGBoost y LightGBM constituyen el
backbone del módulo de predicción tabular. TabNet y FT-Transformer se
evalúan como baselines comparativos, no como propuesta principal, dado
que la evidencia empírica no justifica su adopción en el contexto de
tamaño del dataset empresarial analizado.

\subsubsection{2.2.2 Detector Único versus Ensemble para Detección de
Anomalías}\label{detector-uxfanico-versus-ensemble-para-detecciuxf3n-de-anomaluxedas}

El campo de la detección de anomalías cuenta con una historia de más de
dos décadas de métodos en competencia. Breunig et al.~(2000)
establecieron el Local Outlier Factor (LOF) como referencia para
detectar anomalías locales mediante densidad relativa al vecindario
k-NN, un enfoque sensible a variaciones locales que permite identificar
transacciones con patrones de comportamiento heterogéneos. Liu et
al.~(2008) revolucionaron el campo con Isolation Forest, que aísla
anomalías por particionamiento aleatorio sin necesidad de definir
perfiles de normalidad, con complejidad O(n) que lo hace viable en
millones de transacciones diarias. Ruff et al.~(2018) extendieron la
detección a espacios de representación profundos con Deep SVDD,
capturando patrones no lineales en los datos mediante redes neuronales.
Li et al.~(2022) propusieron ECOD, un detector moderno libre de
parámetros basado en distribución empírica acumulada que supera a 11
baselines en datasets no supervisados, eliminando el riesgo de
sobreajuste al proceso de calibración.

El hallazgo central de Han et al.~(2022) en ADBench ---57 datasets, 30
algoritmos, tres niveles de supervisión--- establece que no existe un
algoritmo universalmente superior: el rendimiento depende fuertemente
del tipo de anomalía, la distribución de los datos y el nivel de
etiquetado disponible. Esta conclusión teórica valida la estrategia de
ensemble como la opción más robusta para entornos de producción donde la
distribución de anomalías es desconocida a priori. La librería PyOD
(Zhao et al., 2019) proporciona la infraestructura técnica para
implementar este ensemble de manera estandarizada y reproducible.

\textbf{Posición de esta tesis}: El ensemble Isolation Forest + LOF +
Deep SVDD (coordinado mediante PyOD) es más robusto que cualquier
detector individual. Esta decisión está respaldada por ADBench (Han et
al., 2022) como fundamento teórico.

\subsubsection{2.2.3 LLM como Detector versus LLM como Generador de
Reportes}\label{llm-como-detector-versus-llm-como-generador-de-reportes}

El surgimiento de los LLMs ha generado propuestas de integración en
sistemas empresariales con distintos roles. Hegselmann et al.~(2023)
demostraron con TabLLM que los LLMs pueden clasificar datos tabulares en
configuración zero/few-shot mediante serialización a texto, con
rendimiento no trivial incluso sin ajuste fino. Park (2024) llevó esta
lógica más lejos con un framework multi-agente donde LLMs especializados
validan alertas de anomalías. Estos antecedentes muestran potencial
metodológico, aunque no resuelven por sí mismos el problema de
trazabilidad operativa agroexportadora.

Sin embargo, existe evidencia sustancial de que usar LLMs como
detectores o tomadores de decisiones introduce riesgos inaceptables. El
survey sobre alucinaciones en LLMs (Maynez et al., 2026) documenta que
los modelos pueden generar razonamiento coherente en forma pero
incorrecto en contenido, con alta confianza aparente. Este riesgo es
especialmente importante en reportes operativos, donde una cifra o causa
inventada puede inducir decisiones equivocadas.

La arquitectura RAG (Schneider et al., 2025 (Schneider et al., 2025))
ofrece una solución al anclar las respuestas del LLM a bases de
conocimiento verificadas ---en este caso, scores, umbrales, vectores
SHAP y fuentes agroexportadoras recuperadas--- reduciendo el espacio de
alucinación al forzar al modelo a narrar únicamente lo que los datos
cuantitativos establecen. El LLM no infiere anomalías; las narra con
evidencias como fundamento.

\textbf{Posición de esta tesis}: El LLM se restringe estrictamente a la
capa de generación de reportes mediante RAG. La detección,
cuantificación y explicación son realizadas por modelos y evidencias
estructuradas (GBDT + ensemble + SHAP). Esta separación se alinea con
principios de transparencia, supervisión humana y trazabilidad
promovidos por marcos como el D.S. N° 115-2025-PCM (PCM, 2025), el EU AI
Act (Parlamento Europeo y Consejo, 2024) y el NIST AI RMF (NIST, 2023).

\subsubsection{2.2.4 Sistemas Aislados versus Sistema Integrado de
Supervisión Operativa
Continua}\label{sistemas-aislados-versus-sistema-integrado-de-supervisiuxf3n-operativa-continua}

La revisión de la literatura evidencia una fragmentación sistemática en
los sistemas de supervisión asistida por IA. Los trabajos pueden
agruparse en cuatro categorías según el módulo que abordan: (1) sistemas
de predicción tabular (Chen \& Guestrin, 2016; Ke et al., 2017;
Prokhorenkova et al., 2018); (2) sistemas de forecasting y series
temporales (Lim et al., 2021; Challu et al., 2022); (3) sistemas de
detección de anomalías (Liu et al., 2008; Han et al., 2022); y (4)
sistemas de generación de reportes con LLMs (Kadir et al., 2025; Park,
2024).

Trabajos como AuditCopilot (Kadir et al., 2025) logran una integración
parcial al combinar detección de anomalías con generación de reportes
LLM, pero operan en dominio contable y no abordan supervisión
agroexportadora. El framework de Park (2024) integra múltiples LLMs pero
opera en mercados financieros de alta frecuencia. AuditMAI
(Waltersdorfer et al., 2024) propone una infraestructura conceptual para
auditoría continua de sistemas de IA. La Tabla 2.2 resume
comparativamente los sistemas más cercanos a la propuesta de esta tesis
desde una perspectiva metodológica.

\textbf{Posición de esta tesis}: Esta investigación cierra la brecha de
integración al proponer y evaluar una arquitectura modular de cuatro
capas que combina predicción tabular, detección de anomalías,
explicabilidad SHAP y generación de reportes LLM-RAG con restricción
anti-alucinación, aplicada a supervisión operativa agroexportadora. BAF
(Jesus et al., 2022) se utiliza solo como benchmark metodológico
complementario; la validación principal se orienta a datos
agroexportadores públicos y sintéticos.

\subsubsection{2.2.5 Contexto Regulatorio Internacional versus
Perú}\label{contexto-regulatorio-internacional-versus-peruxfa}

La mayoría de marcos de gobernanza de IA de la literatura operan en
contextos regulatorios de EE.UU. (NIST AI RMF (NIST, 2023)), Europa (EU
AI Act (Parlamento Europeo y Consejo, 2024), GDPR) o Asia. Estos marcos
coinciden en principios relevantes para esta tesis: documentación,
transparencia, supervisión humana, gestión de riesgos y trazabilidad.

En el contexto peruano, el marco regulatorio ha madurado
significativamente en 2023--2025. La Resolución SBS N° 053-2023
establece lineamientos de gobernanza, trazabilidad y explicabilidad para
modelos de riesgo en entidades supervisadas por la SBS (SBS, 2023), por
lo que se adopta aquí solo como referencia de buenas prácticas. El
Decreto Supremo N° 115-2025-PCM, reglamento de la Ley N° 31814,
proporciona un marco nacional general para promover el uso responsable
de la inteligencia artificial (PCM, 2025). A nivel internacional, el EU
AI Act (Parlamento Europeo y Consejo, 2024) refuerza obligaciones de
transparencia y documentación para sistemas de IA.

\textbf{Posición de esta tesis}: Esta investigación diseña un sistema de
supervisión operativa agroexportadora que incorpora principios de
gobernanza, trazabilidad, documentación y supervisión humana. El D.S. N°
115-2025-PCM se adopta como marco peruano general de IA responsable,
mientras que la Resolución SBS N° 053-2023 se utiliza solo como
referencia nacional de gestión de riesgo de modelos.

\subsubsection{2.2.6 Síntesis y Tabla del Estado del
Arte}\label{suxedntesis-y-tabla-del-estado-del-arte}

La revisión sistemática de los bloques temáticos permite identificar la
brecha de investigación central: \textbf{no existe en la literatura
revisada un sistema orientado al contexto agroexportador peruano que
integre de manera modular, con evaluación reproducible y trazabilidad
explícita, los cuatro componentes}: predicción tabular, detección de
anomalías, explicabilidad SHAP y generación de reportes LLM-RAG basada
en evidencias. Esta tesis propone y evalúa dicha integración para
supervisión operativa agroexportadora.

\textbf{Tabla 2.1 --- Comparativa de Sistemas de Supervisión con IA}

\begin{longtable}[]{@{}
  >{\raggedright\arraybackslash}p{(\columnwidth - 10\tabcolsep) * \real{0.1667}}
  >{\raggedright\arraybackslash}p{(\columnwidth - 10\tabcolsep) * \real{0.1667}}
  >{\raggedright\arraybackslash}p{(\columnwidth - 10\tabcolsep) * \real{0.1667}}
  >{\raggedright\arraybackslash}p{(\columnwidth - 10\tabcolsep) * \real{0.1667}}
  >{\raggedright\arraybackslash}p{(\columnwidth - 10\tabcolsep) * \real{0.1667}}
  >{\raggedright\arraybackslash}p{(\columnwidth - 10\tabcolsep) * \real{0.1667}}@{}}
\toprule\noalign{}
\begin{minipage}[b]{\linewidth}\raggedright
Característica
\end{minipage} & \begin{minipage}[b]{\linewidth}\raggedright
\textbf{Esta tesis}
\end{minipage} & \begin{minipage}[b]{\linewidth}\raggedright
AuditCopilot (Kadir et al., 2025)
\end{minipage} & \begin{minipage}[b]{\linewidth}\raggedright
Park 2024 (Park, 2024)
\end{minipage} & \begin{minipage}[b]{\linewidth}\raggedright
AuditMAI (Waltersdorfer et al., 2024)
\end{minipage} & \begin{minipage}[b]{\linewidth}\raggedright
(JRFM, 2025)
\end{minipage} \\
\midrule\noalign{}
\endhead
\bottomrule\noalign{}
\endlastfoot
Predicción tabular GBDT & ✅ XGB+LGBM+CatBoost & ❌ & ❌ Solo LLMs & ❌
& ✅ Stacking \\
Benchmark público reproducible & ✅ Datos públicos + sintéticos
agroexportadores; BAF complementario & ❌ Dataset propio & ❌ S\&P 500 &
❌ Conceptual & ⚠️ Dataset propio \\
Forecasting series temporales (TFT) & ✅ & ❌ & ❌ & ❌ & ❌ \\
Ensemble de anomalías (ADBench) & ✅ IF+LOF+SVDD & ⚠️ Parcial & ❌ & ❌
& ❌ \\
Explicabilidad SHAP & ✅ TreeSHAP & ❌ & ❌ & ❌ & ✅ SHAP+Anchor \\
Generación LLM de reportes & ✅ RAG determinista & ✅ LLM narrativo & ✅
Multi-agente & ❌ & ❌ \\
Restricción anti-alucinación (RAG+SHAP) & ✅ & ❌ & ❌ & --- & --- \\
Marco de gobernanza explícito & ✅ DS115+NIST+EU AI Act; SBS como
referencia & ❌ & ❌ & ❌ & ❌ \\
Contexto peruano & ✅ Agroexportación peruana & ❌ & ❌ & ❌ & ❌ \\
Evaluación con supervisores/evaluadores & ✅ & ❌ & ❌ & ❌ & ❌ \\
Dominio & Supervisión operativa agroexportadora & Asientos contables &
Mercados bursátiles & Auditoría de IA & Fraude en EEFF \\
\end{longtable}

\textbf{Tabla 2.2 --- Resumen del Estado del Arte}

\begin{longtable}[]{@{}
  >{\centering\arraybackslash}p{(\columnwidth - 6\tabcolsep) * \real{0.1316}}
  >{\raggedright\arraybackslash}p{(\columnwidth - 6\tabcolsep) * \real{0.2895}}
  >{\centering\arraybackslash}p{(\columnwidth - 6\tabcolsep) * \real{0.1316}}
  >{\raggedright\arraybackslash}p{(\columnwidth - 6\tabcolsep) * \real{0.4474}}@{}}
\toprule\noalign{}
\begin{minipage}[b]{\linewidth}\centering
N°
\end{minipage} & \begin{minipage}[b]{\linewidth}\raggedright
Autor(es)
\end{minipage} & \begin{minipage}[b]{\linewidth}\centering
Año
\end{minipage} & \begin{minipage}[b]{\linewidth}\raggedright
Aporte principal
\end{minipage} \\
\midrule\noalign{}
\endhead
\bottomrule\noalign{}
\endlastfoot
1 & Chen \& Guestrin (Chen \& Guestrin, 2016) & 2016 & XGBoost: gradient
boosting escalable con regularización L1/L2, benchmark universal para
datos tabulares \\
2 & Ke et al.~(Ke et al., 2017) & 2017 & LightGBM: GOSS + histogramas,
20× más rápido que XGBoost con precisión comparable \\
3 & Prokhorenkova et al.~(Prokhorenkova et al., 2018) & 2018 & CatBoost:
Ordered Boosting elimina target leakage en variables categóricas de alta
cardinalidad \\
4 & Gorishniy et al.~(Gorishniy et al., 2021) & 2021 & FT-Transformer:
primer Transformer robusto para datos tabulares mediante feature
embeddings \\
5 & Arik \& Pfister (Arik \& Pfister, 2021) & 2021 & TabNet: atención
secuencial interpretable para tablas, combina rendimiento e
interpretabilidad \\
6 & Grinsztajn et al.~(Grinsztajn et al., 2022) & 2022 & GBDT supera a
DL en el 95\% de datasets ≤50K muestras; cierra el debate en contexto
empresarial \\
7 & Lim et al.~(Lim et al., 2021) & 2021 & TFT: forecasting
multi-horizonte con gating de covariables e interpretabilidad
incorporada \\
8 & Liu et al.~(Liu et al., 2008) & 2008 & Isolation Forest: aislamiento
aleatorio O(n), sin perfil de normalidad, escalable a millones de
registros \\
9 & Breunig et al.~(Breunig et al., 2000) & 2000 & LOF: densidad local
relativa k-NN, detecta anomalías locales heterogéneas \\
10 & Ruff et al.~(Ruff et al., 2018) & 2018 & Deep SVDD: detección
one-class en espacio latente profundo para patrones no lineales \\
11 & Han et al.~(Han et al., 2022) & 2022 & ADBench: benchmark de 30
algoritmos en 57 datasets; ensembles son más robustos que detectores
únicos \\
12 & Lundberg \& Lee (Lundberg \& Lee, 2017) & 2017 & SHAP: valores
Shapley con consistencia axiomática; TreeSHAP exacto para GBDT \\
13 & Kadir et al.~(Kadir et al., 2025) & 2025 & AuditCopilot:
LLM+detección en asientos contables; antecedente metodológico para
reportes automáticos \\
14 & Park (Park, 2024) & 2024 & Framework multi-agente LLM para validar
anomalías en mercados bursátiles \\
15 & Schneider et al.~(Schneider et al., 2025) & 2025 & RAG avanzado
para BI organizacional; arquitectura anti-alucinación base de esta
tesis \\
16 & SBS Perú (SBS, 2023) & 2023 & Resolución N° 053-2023: referencia
nacional de buenas prácticas para gestión de riesgo de modelos \\
\end{longtable}

\begin{center}\rule{0.5\linewidth}{0.5pt}\end{center}

\begin{center}\rule{0.5\linewidth}{0.5pt}\end{center}

\subsection{2.3 Marco Conceptual}\label{marco-conceptual}

\subsubsection{2.3.1 Reconocimiento de Patrones y Aprendizaje
Automático}\label{reconocimiento-de-patrones-y-aprendizaje-automuxe1tico}

El reconocimiento de patrones es la disciplina de la inteligencia
artificial que busca identificar regularidades, estructuras o relaciones
en datos a partir de ejemplos históricos. En el aprendizaje automático,
este proceso se formaliza mediante modelos que aprenden una función
\(f: X \rightarrow Y\) a partir de un conjunto de entrenamiento
\(\{(x_i, y_i)\}_{i=1}^n\), con el objetivo de generalizar hacia
instancias no observadas.

Se distinguen tres paradigmas principales: \textbf{aprendizaje
supervisado}, donde el modelo aprende de etiquetas explícitas (e.g.,
operación normal / operación anómala); \textbf{aprendizaje no
supervisado}, que identifica patrones sin etiquetas y es fundamental en
la detección de anomalías; y \textbf{aprendizaje semi-supervisado}, que
combina ambos cuando el etiquetado es costoso o escaso, escenario
habitual en procesos agroexportadores donde no todas las desviaciones
quedan registradas formalmente.

En el contexto de esta investigación, el reconocimiento de patrones
opera en dos dimensiones complementarias: la detección de registros
operativos anómalos (patrones puntuales) y la identificación de
secuencias temporales de comportamiento irregular (patrones colectivos),
requiriendo tanto modelos de clasificación supervisada como detectores
de anomalías no supervisados.

\subsubsection{2.3.2 Datos Tabulares en Sistemas Agroexportadores
Empresariales}\label{datos-tabulares-en-sistemas-agroexportadores-empresariales}

Los datos tabulares son la forma predominante de almacenamiento en los
sistemas de información empresarial: cada fila representa una instancia
operativa y cada columna una variable. En agroexportación, una instancia
puede describir un lote, día, producto, zona, envío o registro de
mercado; las columnas pueden incluir precio, volumen, temperatura,
precipitación, humedad, destino, cumplimiento fitosanitario, días
logísticos, merma y etiqueta de anomalía.

Estas propiedades estructurales explican por qué los GBDT son adecuados
para este dominio (Grinsztajn et al., 2022): manejan variables numéricas
y categóricas heterogéneas, toleran valores faltantes, capturan
relaciones no lineales y requieren menor ingeniería de features que
arquitecturas neuronales complejas. Cuando las etiquetas de anomalía son
escasas o desbalanceadas, el rendimiento debe evaluarse con métricas
orientadas a precisión-recall (PR-AUC), F1-Score y análisis de falsos
positivos.

\subsubsection{2.3.3 Gradient Boosting Decision Trees
(GBDT)}\label{gradient-boosting-decision-trees-gbdt}

Los Gradient Boosting Decision Trees (GBDT) son una familia de
algoritmos de aprendizaje supervisado que construyen modelos predictivos
mediante la combinación secuencial de múltiples árboles de decisión
débiles. El enfoque fue formalizado por Friedman (2001) como ``Greedy
Function Approximation'', donde cada árbol nuevo se ajusta para corregir
los errores residuales del conjunto de árboles previos mediante descenso
de gradiente en el espacio funcional de la función de pérdida.

La formulación matemática central de GBDT busca encontrar una función
\(F(x)\) que minimice la pérdida esperada:

\[F^*(x) = \arg\min_{F} \mathbb{E}_{y,x}[L(y, F(x))]\]

donde \(L\) es la función de pérdida (e.g., entropía cruzada para
clasificación, MSE para regresión) y la minimización se realiza
iterativamente añadiendo árboles de regresión \(h_m(x)\) con pesos de
aprendizaje \(\nu\):

\[F_m(x) = F_{m-1}(x) + \nu \cdot h_m(x)\]

\textbf{XGBoost} (Chen \& Guestrin, 2016) introduce mejoras clave sobre
el GBDT estándar: regularización L1 y L2 en la función objetivo para
controlar la complejidad del modelo, manejo nativo de valores faltantes
mediante aprendizaje automático de la dirección de ramificación, y
paralelización por columnas en lugar de por filas, lo que habilita el
procesamiento en datasets de alta dimensión.

\textbf{LightGBM} (Ke et al., 2017) acelera el entrenamiento mediante
dos innovaciones: Gradient-based One-Side Sampling (GOSS), que retiene
las muestras con mayor gradiente y descarta aleatoriamente las de menor
gradiente, preservando la distribución sin pérdida estadística
significativa; y Exclusive Feature Bundling (EFB), que agrupa features
mutuamente excluyentes para reducir dimensionalidad efectiva. El
resultado es una aceleración de hasta 20× sobre XGBoost con precisión
comparable.

\textbf{CatBoost} (Prokhorenkova et al., 2018) resuelve el problema del
target leakage en variables categóricas mediante Ordered Boosting:
calcula las estadísticas de objetivo para cada categoría usando
únicamente las observaciones previas en un orden aleatorio permutado,
evitando que la información del objetivo filtre hacia las features de
entrada durante el entrenamiento. Esta propiedad es especialmente
relevante en datos contables, donde variables como ``código de cuenta''
o ``centro de costo'' tienen alta cardinalidad.

La justificación empírica para elegir GBDT sobre Deep Learning en datos
tabulares empresariales está sólidamente documentada por (Grinsztajn et
al., 2022): en 45 datasets con hasta 50,000 muestras, los GBDT superan a
arquitecturas neuronales complejas (como FT-Transformer o TabNet) en la
inmensa mayoría de los casos. Dado que el dataset operativo de esta
tesis comprende un volumen controlado de hasta 10,000 registros
transaccionales sintéticos (Anexo C), se sitúa exactamente en el rango
donde los modelos basados en árboles maximizan su ventaja comparativa.
Esta superioridad es atribuible a tres propiedades estructurales de los
datos tabulares que los árboles aprovechan mejor que las redes
neuronales.

\subsubsection{2.3.4 Detección de Anomalías y Estrategia de
Ensemble}\label{detecciuxf3n-de-anomaluxedas-y-estrategia-de-ensemble}

La detección de anomalías es el problema de identificar observaciones
que se desvían significativamente del comportamiento esperado del
conjunto de datos. La literatura distingue tres tipos fundamentales de
anomalías (Han et al., 2022): (a) \textbf{puntuales} --- instancias
individuales anómalas (e.g., una transacción de monto atípico); (b)
\textbf{contextuales} --- instancias que son anómalas en un contexto
particular pero no en general (e.g., un cargo nocturno inusual para un
perfil de usuario); y (c) \textbf{colectivas} --- secuencias de
instancias que son anómalas en conjunto aunque cada una individualmente
no lo sea (e.g., un patrón de micro-transacciones).

\textbf{Isolation Forest} (Liu et al., 2008) se basa en el principio de
que las anomalías son ``pocas y diferentes'': son más fáciles de aislar
que los puntos normales mediante particionamiento aleatorio del espacio.
Un árbol de aislamiento selecciona aleatoriamente una feature y un valor
de corte; la anomalía será aislada en pocas particiones (camino corto),
mientras que los puntos normales requieren muchas particiones (camino
largo). El score de anomalía es el inverso de la longitud promedio del
camino de aislamiento, normalizada según la longitud esperada para un
punto normal en un conjunto de tamaño \(n\). La complejidad es O(n log
n) en entrenamiento y O(n) en inferencia.

\textbf{Local Outlier Factor (LOF)} (Breunig et al., 2000) cuantifica el
grado de anomalía de cada punto en función de la densidad de su
vecindario local respecto a la densidad de sus vecinos. El score LOF
para el punto \(p\) se define como:

\[\text{LOF}_k(p) = \frac{\sum_{o \in N_k(p)} \frac{\text{lrd}_k(o)}{\text{lrd}_k(p)}}{|N_k(p)|}\]

donde \(\text{lrd}_k\) es la densidad de alcanzabilidad local. Un valor
LOF \textgreater\textgreater{} 1 indica que \(p\) tiene densidad local
mucho menor que sus vecinos, lo que lo caracteriza como anomalía. LOF es
sensible a variaciones locales de densidad, lo que lo hace
complementario a Isolation Forest en datasets heterogéneos.

\textbf{Deep SVDD} (Ruff et al., 2018) extiende el Support Vector Data
Description al espacio de representación de redes neuronales: entrena
una red para mapear los datos normales hacia el interior de una
hipersfera mínima en el espacio latente. Las anomalías se detectan como
puntos que caen fuera o lejos de esta hipersfera. La función objetivo
minimiza el volumen de la hipersfera:

\[\min_{W, R, c} R^2 + \frac{1}{\nu n} \sum_{i=1}^{n} \max(0, \|f(x_i; W) - c\|^2 - R^2)\]

donde \(f(x_i; W)\) es la representación de la red neuronal, \(c\) es el
centro de la hipersfera y \(R\) es su radio.

\textbf{ECOD} (Li et al., 2022) calcula el score de anomalía como la
probabilidad acumulada de observar un punto tan extremo como \(x\) bajo
la distribución empírica del dataset, estimada mediante funciones de
distribución acumulada (ECDF) multivariadas. Su ventaja principal es que
no tiene hiperparámetros que calibrar, eliminando el riesgo de
sobreajuste y simplificando el despliegue en producción.

La estrategia de \textbf{ensemble} consolida las puntuaciones de
múltiples detectores mediante aggregation functions (promedio de scores,
votación por mayoría o meta-clasificación). El fundamento teórico lo
proporciona ADBench (Han et al., 2022): no existe un algoritmo
universal, y el ensemble reduce la varianza del estimador de anomalía
agregando perspectivas complementarias. PyOD (Zhao et al., 2019)
implementa esta estrategia con la clase \texttt{LSCP} (Locally Selective
Combination in Parallel Outlier Ensembles) y otras técnicas de
combinación estándar.

\subsubsection{2.3.5 Forecasting de Series Temporales con
Transformers}\label{forecasting-de-series-temporales-con-transformers}

Las series temporales agroexportadoras presentan tres desafíos que los
modelos de forecasting deben resolver: tendencia no estacionaria,
estacionalidad múltiple (diaria, semanal, mensual, anual) y dependencia
de covariables exógenas (clima, calendario agrícola, demanda
internacional, precios y condiciones logísticas). Los modelos
estadísticos clásicos como ARIMA capturan relaciones lineales con
eficacia, pero presentan limitaciones en la modelización de
no-linealidades y horizontes largos.

\textbf{Temporal Fusion Transformer (TFT)} (Lim et al., 2021) propone
una arquitectura especializada que combina cuatro mecanismos: (1)
codificación LSTM para dependencias secuenciales locales; (2) selección
de variables con mecanismo de gating (GLU --- Gated Linear Unit) que
identifica automáticamente las covariables más informativas; (3)
atención multi-cabezal interpretable que pondera los pasos temporales
según su relevancia predictiva; y (4) red de cuantiles para cuantificar
la incertidumbre de la predicción. TFT acepta tres tipos de entradas:
features estáticas conocidas (e.g., producto, zona, destino),
covariables futuras conocidas (e.g., calendario agrícola, campañas,
feriados) y covariables históricas observadas (e.g., precio, volumen,
clima o merma pasada).

El debate sobre la efectividad de los Transformers en series temporales
es relevante para esta tesis. Zeng et al.~(2023) argumentan que DLinear
---un modelo lineal simple--- supera a los Transformers en múltiples
benchmarks, atribuyendo la limitación de los Transformers al hecho de
que el mecanismo de self-attention es permutation-invariant y destruye
el orden temporal de las secuencias. Sin embargo, este argumento ha sido
rebatido sucesivamente: Nie et al.~(2023) demuestran que la tokenización
por patches ---agrupando segmentos temporales antes de aplicar
atención--- preserva el orden local y supera a DLinear en la mayoría de
benchmarks de horizonte largo. Liu et al.~(2024) proponen invertir la
tokenización: en lugar de tokenizar por timestamp, tokenizan por
variable, aplicando self-attention entre variables en lugar de entre
tiempos, obteniendo SOTA en 7 datasets multivariados.

\textbf{Posición de esta tesis respecto al debate}: TFT se considera por
su interpretabilidad incorporada ---el mecanismo de gating y los mapas
de atención son legibles por analistas--- más que exclusivamente por su
rendimiento predictivo. En el contexto agroexportador, la capacidad de
justificar qué períodos temporales y qué covariables fundamentan la
predicción es un requerimiento funcional comparable en importancia a la
precisión numérica.

N-HiTS (Challu et al., 2022) ofrece una alternativa no-Transformer para
forecasting de horizonte largo, con interpolación jerárquica multi-tasa
que reduce la complejidad computacional respecto a N-BEATS. Chronos
(Ansari et al., 2024) representa el paradigma emergente de los
foundation models para series temporales, basado en T5, que logra
performance zero-shot competitivo en múltiples datasets; sin embargo, su
opacidad y dependencia de infraestructura de gran escala limitan su
aplicación directa cuando se requiere trazabilidad operativa.

\subsubsection{2.3.6 Explicabilidad mediante Valores de Shapley
(SHAP)}\label{explicabilidad-mediante-valores-de-shapley-shap}

La explicabilidad en sistemas de IA se clasifica en dos grandes
categorías: \textbf{inherente} ---modelos cuya estructura es
intrínsecamente interpretable, como los árboles de decisión--- y
\textbf{post-hoc} ---métodos aplicados a cualquier modelo después del
entrenamiento para interpretar sus predicciones. SHAP y LIME son los dos
métodos post-hoc agnósticos más adoptados en la literatura.

\textbf{LIME} (Local Interpretable Model-agnostic Explanations) (Ribeiro
et al., 2016) genera explicaciones locales construyendo un modelo lineal
sustituto en el vecindario de la instancia a explicar, ponderando las
muestras según su proximidad al punto de interés. LIME es rápido y
flexible, pero produce explicaciones inestables: pequeñas perturbaciones
de la instancia pueden cambiar significativamente la explicación, un
problema crítico en contextos forenses.

\textbf{SHAP} (Lundberg \& Lee, 2017) fundamenta las explicaciones en
los valores de Shapley de la teoría de juegos cooperativos. El valor
SHAP de la feature \(i\) para la predicción \(f(x)\) cuantifica la
contribución marginal media de \(i\) a la predicción, promediando sobre
todas las coaliciones posibles de features:

\[\phi_i = \sum_{S \subseteq F \setminus \{i\}} \frac{|S|!(|F|-|S|-1)!}{|F|!} [f(S \cup \{i\}) - f(S)]\]

donde \(F\) es el conjunto de todas las features y \(S\) es una
coalición de features sin \(i\). Esta formulación garantiza cuatro
propiedades axiomáticas: (a) \textbf{eficiencia} --- la suma de todos
los valores SHAP iguala la diferencia entre la predicción y el valor
esperado; (b) \textbf{simetría} --- features con contribución idéntica
reciben el mismo valor; (c) \textbf{dummy} --- features sin efecto
tienen valor cero; y (d) \textbf{aditividad} --- los valores SHAP son
consistentes al combinar modelos.

SHAP resuelve las limitaciones de LIME al garantizar consistencia: si un
modelo cambia la predicción al aumentar la contribución de una feature,
el valor SHAP de esa feature nunca disminuye (Lundberg \& Lee, 2017).
\textbf{TreeSHAP} extiende este cálculo con un algoritmo exacto en
O(TLD²) para modelos basados en árboles ---donde T es el número de
árboles, L es el número de hojas por árbol y D es la profundidad
máxima--- haciendo el cálculo computacionalmente viable para GBDT en
producción.

En el contexto de supervisión operativa, la estabilidad de las
explicaciones permite verificar que el modelo asigna importancias
consistentes a variables semejantes. Un índice alto de estabilidad
fortalece la confianza en el sistema, porque evita que alertas similares
reciban justificaciones contradictorias.

La integración SHAP+LLM de esta tesis opera como sigue: los vectores
SHAP de una alerta operativa (una lista de pares variable→contribución
cuantitativa) se incorporan como contexto verificado en el RAG, y el LLM
genera la narración del informe sin posibilidad de inventar cifras que
no estén en esos vectores o en las fuentes recuperadas.

\subsubsection{2.3.7 Modelos de Lenguaje y Arquitectura RAG para
Generación de
Reportes}\label{modelos-de-lenguaje-y-arquitectura-rag-para-generaciuxf3n-de-reportes}

Los Modelos de Lenguaje de Gran Tamaño (LLMs) son sistemas entrenados
mediante autoregresión en corpus masivos de texto para aprender
distribuciones probabilísticas sobre secuencias de tokens. Su capacidad
de generalización les permite realizar tareas de reasoning, traducción,
resumen y generación de texto con calidad próxima a la humana en
configuraciones zero-shot y few-shot.

\textbf{In-context learning} permite guiar el comportamiento del LLM
mediante ejemplos incluidos directamente en el prompt, sin necesidad de
ajuste fino (fine-tuning). TabLLM (Hegselmann et al., 2023) demostró que
mediante serialización de datos tabulares a texto descriptivo, los LLMs
pueden realizar clasificación sobre datos estructurados con rendimiento
no trivial en zero-shot, ampliando el espectro de aplicación de estos
modelos más allá del texto no estructurado.

Sin embargo, el uso de LLMs como agentes de decisión autónoma introduce
el riesgo de \textbf{alucinaciones}: el modelo puede generar
afirmaciones coherentes en forma pero incorrectas en contenido (Ji et
al., 2023; Maynez et al., 2026). La literatura distingue al menos dos
tipos: (a) alucinaciones intrínsecas, en las que el texto generado
contradice la información del contexto recuperado; y (b) alucinaciones
extrínsecas, en las que el modelo inventa información no presente en el
contexto. En particular, las ``alucinaciones numéricas'' ---valores
específicos de métricas, porcentajes o fechas que no corresponden a los
datos reales (Barclays Research, 2025)--- son especialmente peligrosas
en reportes operativos, porque pueden inducir decisiones equivocadas
pese a la apariencia de precisión cuantitativa.

\textbf{Retrieval-Augmented Generation (RAG)} (Lewis et al., 2020;
Schneider et al., 2025) reduce este riesgo al separar el conocimiento
factual del modelo generativo: en lugar de que el LLM ``recuerde''
información de su entrenamiento, el sistema recupera documentos o datos
relevantes de una base de conocimiento externa verificada y los incluye
en el contexto del prompt. El LLM entonces genera texto fundamentado en
esos datos recuperados, no en su memoria paramétrica. Es importante
señalar que RAG \textbf{reduce significativamente pero no elimina} el
riesgo de alucinación; persisten casos de alucinación intrínseca
(faithful hallucination) en los que el modelo genera afirmaciones que se
desvían del contexto recuperado. Técnicas avanzadas como GraphRAG
incorporan grafos de conocimiento para recuperación semántica más rica,
mientras que Self-RAG permite al modelo verificar la pertinencia de los
documentos recuperados antes de usarlos.

En la arquitectura de esta tesis, la ``base de conocimiento'' del RAG
son los vectores SHAP de la alerta analizada, las métricas del ensemble
de detección, las fuentes agroexportadoras recuperadas y las reglas de
reporte definidas. El LLM recibe ese contexto verificado y genera el
informe narrativo sin acceso a conocimiento adicional no validado.
Adicionalmente se aplican dos controles complementarios: (a) plantillas
de prompt estructurado con campos obligatorios (dato, modelo, score,
umbral, explicación SHAP, fuente recuperada), y (b) validación posterior
del reporte contra los vectores SHAP de entrada para detectar
discrepancias numéricas. Este diseño permite que cada afirmación del
reporte pueda trazarse hasta una fuente, score, umbral o variable
explicativa.

La evaluación de calidad de los reportes generados puede utilizar
\textbf{ROUGE} (Recall-Oriented Understudy for Gisting Evaluation)
cuando exista un texto de referencia. Sin embargo, para esta tesis se
prioriza una rúbrica operativa de completitud, consistencia,
accionabilidad y correspondencia con evidencias, porque la calidad de un
reporte de supervisión depende no solo de similitud textual, sino de su
utilidad para la toma de decisiones.

\subsubsection{2.3.8 Gobernanza de IA y
MLOps}\label{gobernanza-de-ia-y-mlops}

El despliegue de sistemas de IA en entornos empresariales críticos
requiere un marco de gobernanza que trascienda el rendimiento técnico.
Sculley et al.~(2015) documentaron la ``deuda técnica oculta'' en
sistemas de ML: más del 95\% del código de un sistema ML de producción
no es el modelo en sí, sino la infraestructura de ingesta, validación,
features, servicio y monitoreo. Los pipelines con alto acoplamiento
entre componentes generan ``entanglement'' que dificulta el
mantenimiento y aumenta el riesgo de regresiones silenciosas.

\textbf{MLOps} (Kreuzberger et al., 2022) establece el conjunto de
prácticas para gestionar el ciclo de vida completo de los modelos ML en
producción: integración y entrega continua (CI/CD) para modelos,
monitoreo de data drift y model drift, automatización del
reentrenamiento, y trazabilidad de versiones de datos, código y modelos.
En el contexto de supervisión operativa agroexportadora, MLOps permite
reproducir qué modelo generó una alerta, con qué datos de entrada, bajo
qué versión y con qué umbral.

El \textbf{NIST Artificial Intelligence Risk Management Framework (AI
RMF 1.0)} (NIST, 2023) proporciona cuatro funciones de gestión de riesgo
para sistemas de IA: (1) \textbf{Govern} --- establecer políticas y
roles de responsabilidad; (2) \textbf{Map} --- identificar el contexto
de despliegue y los riesgos asociados; (3) \textbf{Measure} --- evaluar
los riesgos con métricas verificables; y (4) \textbf{Manage} ---
implementar controles y mitigaciones. La arquitectura modular de esta
tesis es diseñada para que cada capa corresponda a responsabilidades
verificables bajo este framework.

\textbf{Datasheets for Datasets} (Gebru et al., 2021) propone una
plantilla de documentación estandarizada para datasets que detalla:
motivación de recolección, proceso de recolección, composición,
preprocesamiento aplicado, distribución permitida y consideraciones
éticas. Esta práctica se aplicará al dataset sintético agroexportador y
a las fuentes públicas utilizadas, garantizando que los resultados
reportados en esta tesis sean reproducibles y que las limitaciones de
cada fuente estén identificadas antes de evaluar el sistema.

\textbf{Model Cards} (Mitchell et al., 2019) extiende la documentación
al nivel del modelo, especificando para quién fue entrenado, en qué
condiciones, cuáles son sus limitaciones conocidas y cómo debe usarse de
manera responsable. En esta tesis, se elaboran Model Cards para los
modelos XGBoost/LightGBM, detectores de anomalías y el componente
LLM+RAG, en conformidad con los principios de documentación del NIST AI
RMF (NIST, 2023).

El contexto peruano e internacional consolida la necesidad de este
framework de gobernanza. El D.S. N° 115-2025-PCM (PCM, 2025), el NIST AI
RMF (NIST, 2023) y el EU AI Act (Parlamento Europeo y Consejo, 2024)
refuerzan principios comunes: transparencia, documentación, supervisión
humana y gestión de riesgos. Estos principios son adoptados como
referencia de diseño para esta tesis. No se afirma cumplimiento formal
con ninguno de estos marcos ---tal afirmación requeriría auditoría
regulatoria externa--- sino conformidad de diseño con sus principios, en
particular: (a) transparencia mediante SHAP y Model Cards, (b)
supervisión humana mediante revisión obligatoria de cada reporte antes
de su uso operativo, (c) gestión de riesgos mediante umbrales calibrados
y validación cruzada, y (d) trazabilidad documental mediante logs
completos por alerta. La aplicabilidad regulatoria efectiva del sistema
a una empresa específica depende de su clasificación de riesgo bajo el
reglamento correspondiente y queda fuera del alcance de esta tesis.

\subsubsection{2.3.9 Supervisión Operativa, Trazabilidad e Inteligencia
Artificial}\label{supervisiuxf3n-operativa-trazabilidad-e-inteligencia-artificial}

La supervisión operativa en agroexportación exige monitorear procesos
que combinan producción, acopio, calidad, sanidad, logística y comercio
exterior. Las anomalías en este dominio no necesariamente corresponden a
fraude; pueden representar variaciones atípicas de precio, caídas de
volumen, condiciones climáticas adversas, mermas elevadas,
incumplimientos fitosanitarios o retrasos logísticos. Por ello, el
sistema propuesto se orienta a detectar desviaciones relevantes para la
toma de decisiones, no a sustituir procesos de investigación legal o
auditoría financiera.

La \textbf{supervisión operativa continua} busca reemplazar ciclos de
revisión tardíos por monitoreo frecuente y documentado de indicadores.
En este enfoque, cada alerta debe registrar el dato de origen, el modelo
aplicado, el score calculado, el umbral utilizado, las variables
explicativas y el reporte generado. Esta trazabilidad permite que un
supervisor operativo comprenda por qué el sistema marcó un evento como
anómalo y qué evidencia respalda la recomendación.

La integración de IA en supervisión operativa plantea el problema de la
confianza en decisiones automáticas. Esta exigencia convierte a la
explicabilidad (SHAP), la documentación de datasets (Datasheets), la
documentación de modelos (Model Cards) y los logs de decisión en
componentes funcionales del sistema. En el marco peruano, el D.S. N°
115-2025-PCM (PCM, 2025) proporciona una base general para el uso
responsable de IA; la Resolución SBS N° 053-2023 (SBS, 2023) se conserva
solo como referencia nacional de buenas prácticas para gestión de riesgo
de modelos.

\begin{center}\rule{0.5\linewidth}{0.5pt}\end{center}

\begin{center}\rule{0.5\linewidth}{0.5pt}\end{center}

\section{CAPÍTULO III: PROPUESTA
METODOLÓGICA}\label{capuxedtulo-iii-propuesta-metodoluxf3gica}

\subsection{3.1 Arquitectura del Sistema
Integrado}\label{arquitectura-del-sistema-integrado}

La arquitectura propuesta se divide en cuatro módulos secuenciales,
diseñados para maximizar trazabilidad, interpretabilidad y utilidad
operativa en procesos agroexportadores:

\begin{itemize}
\tightlist
\item
  \textbf{Módulo de Predicción Tabular (Capa 1):} Utiliza algoritmos
  GBDT como núcleo predictivo, priorizando XGBoost (Chen \& Guestrin,
  2016) y LightGBM (Ke et al., 2017) por su robustez ante datos
  tabulares con variables heterogéneas (Grinsztajn et al., 2022). El
  módulo puede estimar valores esperados de precio, volumen, merma o
  riesgo operativo.
\item
  \textbf{Módulo de Detección de Anomalías (Capa 2):} Emplea detectores
  como Isolation Forest (Liu et al., 2008), LOF (Breunig et al., 2000) y
  ECOD (Li et al., 2022), orquestados mediante PyOD (Zhao et al., 2019),
  para identificar comportamientos atípicos en variables
  agroexportadoras. Se selecciona ECOD sobre Deep SVDD (Ruff et al.,
  2018) ---considerado en la revisión bibliográfica del Capítulo II---
  por tres razones: (a) ECOD no requiere ajuste de hiperparámetros, lo
  cual elimina el riesgo de sobreajuste en la calibración; (b) su
  fundamento basado en funciones de distribución empírica acumulada es
  interpretable estadísticamente para auditores, mientras que Deep SVDD
  opera sobre representaciones latentes opacas; y (c) su complejidad
  computacional es lineal, apropiada para el tamaño medio del dataset
  experimental (2,000--5,000 registros). Deep SVDD se mantiene como
  referencia conceptual del Capítulo II por su valor histórico, pero no
  entra al ensemble final.
\item
  \textbf{Módulo de Explicabilidad (Capa 3):} SHAP (Lundberg \& Lee,
  2017) genera explicaciones locales por alerta, identificando qué
  variables ---precio, volumen, clima, destino, cumplimiento o merma---
  contribuyen al score del sistema.
\item
  \textbf{Módulo de Reportes LLM+RAG (Capa 4):} Un LLM restringido a
  evidencias estructuradas mediante RAG (Schneider et al., 2025) redacta
  reportes operativos trazables. El LLM no decide si existe una
  anomalía; solo traduce scores, umbrales y explicaciones SHAP a
  lenguaje comprensible.
\end{itemize}

\begin{verbatim}
┌─────────────────────────────────────────────────────────┐
│ CAPA 4: Reporte Automatizado (LLM + RAG)               │
│ Entrada: anomalías + vectores SHAP                      │
│ Salida: reporte auditado en lenguaje natural (MD/PDF)  │
└─────────────────────────────────────────────────────────┘
                         ↑
┌─────────────────────────────────────────────────────────┐
│ CAPA 3: Explicabilidad (SHAP / TreeSHAP)               │
│ Entrada: predicciones + datos originales               │
│ Salida: vectores Shapley + SHAP Stability Index        │
└─────────────────────────────────────────────────────────┘
                         ↑
┌─────────────────────────────────────────────────────────┐
│ CAPA 2: Detección de Anomalías (Ensemble)              │
│ Métodos: IF + LOF + ECOD (PyOD)                        │
│ Salida: score anomalía + método que detectó            │
└─────────────────────────────────────────────────────────┘
                         ↑
┌─────────────────────────────────────────────────────────┐
│ CAPA 1: Predicción tabular agroexportadora             │
│ Modelo: XGBoost / LightGBM                             │
│ Entrada: precio, volumen, clima, calidad, logística    │
│ Salida: valor esperado + riesgo operativo              │
└─────────────────────────────────────────────────────────┘
\end{verbatim}

\subsection{3.2 Fuentes de Datos, Dataset Sintético y
Benchmarks}\label{fuentes-de-datos-dataset-sintuxe9tico-y-benchmarks}

Para validar la robustez del sistema sin depender obligatoriamente de
datos privados, se emplearán tres niveles de información:

\begin{enumerate}
\def\labelenumi{\arabic{enumi}.}
\tightlist
\item
  \textbf{Fuentes públicas oficiales del dominio agroexportador}:
  MIDAGRI para agroexportaciones, precios y productos; SENAMHI para
  variables climáticas; SENASA para requisitos fitosanitarios; SUNAT
  para exportaciones; INEI para indicadores económicos; FAOSTAT y UN
  Comtrade para validación internacional.
\item
  \textbf{Dataset sintético agroexportador documentado}: conjunto de
  registros simulados con variables como fecha, producto, zona, volumen,
  precio, temperatura, precipitación, humedad, destino, cumplimiento
  fitosanitario, días logísticos, merma, etiqueta de anomalía y tipo de
  anomalía. Este dataset se documentará con criterios de Datasheets for
  Datasets (Gebru et al., 2021).
\item
  \textbf{Benchmark metodológico complementario}: BAF Benchmark (Jesus
  et al., 2022) podrá utilizarse solo para contrastar comportamiento de
  modelos en datos tabulares desbalanceados con drift temporal, sin
  presentarlo como validación directa del dominio agroexportador.
\end{enumerate}

\subsection{3.3 Configuración Experimental y
Métricas}\label{configuraciuxf3n-experimental-y-muxe9tricas}

\subsubsection{3.3.1 Métricas por variable
dependiente}\label{muxe9tricas-por-variable-dependiente}

\begin{itemize}
\tightlist
\item
  \textbf{Métricas de predicción y detección (VD1)}: PR-AUC (métrica
  principal para datasets desbalanceados), ROC-AUC, F1-Score, Precision
  y Recall con umbral óptimo determinado por el punto de máxima F1 sobre
  el conjunto de validación.
\item
  \textbf{Métricas de explicabilidad (VD2)}: cobertura top-k SHAP
  (porcentaje de alertas en las que las k=5 variables principales
  explican ≥80\% de la magnitud absoluta del score), consistencia
  cualitativa de variables explicativas y claridad operativa (Likert
  1--5 evaluada por revisores con perfil agroexportador).
\item
  \textbf{Métricas de calidad de reportes (VD3)}: rúbrica operativa de
  cinco dimensiones (completitud, consistencia numérica, accionabilidad,
  coherencia textual, correspondencia con evidencias) evaluada por dos
  revisores independientes con cálculo de Kappa de Cohen para
  confiabilidad inter-evaluador. Adicionalmente ROUGE-1/ROUGE-L cuando
  exista referencia humana.
\item
  \textbf{Métricas de comprensión y decisión (VD4)}: tiempo-a-decisión
  (segundos, medido automáticamente desde la apertura de la alerta hasta
  el envío del veredicto del evaluador), comprensión de alerta (Likert
  1--5) y decisión final correcta (sí/no respecto a la etiqueta del
  dataset).
\item
  \textbf{Métricas de trazabilidad (VD5)}: porcentaje de alertas con
  todos los campos completos (dato de origen, versión de dataset,
  modelo, score, umbral, explicación SHAP, fuente recuperada por RAG y
  reporte generado).
\end{itemize}

\subsubsection{3.3.2 División del dataset y
semilla}\label{divisiuxf3n-del-dataset-y-semilla}

Para evitar fuga de información temporal en variables agroexportadoras
con estacionalidad, se aplica una \textbf{división cronológica} y no
aleatoria: - \textbf{Train}: primeros 70\% de registros ordenados por
fecha. - \textbf{Validation}: 10\% siguiente, para selección de
hiperparámetros. - \textbf{Test}: 20\% final, evaluado solo al cierre
del entrenamiento.

Todas las ejecuciones experimentales fijan \texttt{np.random.seed(42)} y
\texttt{random.seed(42)}. Cada experimento se repite con cinco semillas
adicionales (43, 44, 45, 46, 47) para reportar media ± desviación
estándar de cada métrica.

\subsubsection{3.3.3 Diseño experimental: condiciones y experimentos
E1--E5}\label{diseuxf1o-experimental-condiciones-y-experimentos-e1e5}

La evaluación se organiza en cinco experimentos cuya condición
experimental aísla un componente arquitectónico distinto:

\begin{longtable}[]{@{}
  >{\raggedright\arraybackslash}p{(\columnwidth - 10\tabcolsep) * \real{0.1667}}
  >{\raggedright\arraybackslash}p{(\columnwidth - 10\tabcolsep) * \real{0.1667}}
  >{\raggedright\arraybackslash}p{(\columnwidth - 10\tabcolsep) * \real{0.1667}}
  >{\raggedright\arraybackslash}p{(\columnwidth - 10\tabcolsep) * \real{0.1667}}
  >{\raggedright\arraybackslash}p{(\columnwidth - 10\tabcolsep) * \real{0.1667}}
  >{\raggedright\arraybackslash}p{(\columnwidth - 10\tabcolsep) * \real{0.1667}}@{}}
\toprule\noalign{}
\begin{minipage}[b]{\linewidth}\raggedright
Exp.
\end{minipage} & \begin{minipage}[b]{\linewidth}\raggedright
Nombre
\end{minipage} & \begin{minipage}[b]{\linewidth}\raggedright
Condición experimental
\end{minipage} & \begin{minipage}[b]{\linewidth}\raggedright
Condición de control
\end{minipage} & \begin{minipage}[b]{\linewidth}\raggedright
Variable observada
\end{minipage} & \begin{minipage}[b]{\linewidth}\raggedright
Hipótesis
\end{minipage} \\
\midrule\noalign{}
\endhead
\bottomrule\noalign{}
\endlastfoot
E1 & Rendimiento de detección & Ensemble IF + LOF + ECOD & Isolation
Forest individual & VD1: PR-AUC, F1 & H1a \\
E2 & Aporte de SHAP & Sistema con vectores SHAP & Sistema sin SHAP (solo
scores) & VD2: cobertura top-k, Likert & H1b \\
E3 & Aporte de RAG & LLM + RAG (anclado en SHAP) & LLM libre (sin RAG) &
VD3: rúbrica 5D, ROUGE-L & H1c \\
E4 & Sistema integrado vs.~aislado & Pipeline completo de 4 capas &
Salidas técnicas aisladas por módulo & VD4: tiempo, Likert; VD5:
trazabilidad & H1, H1d \\
E5 & Ablation study & Configuraciones parciales (E5a, E5b, E5c, E5d) &
--- & VD1 + VD5 por configuración & Contribución por capa \\
\end{longtable}

Variantes del ablation study (E5): - \textbf{E5a}: Solo Capa 2 (sin
predicción, sin SHAP, sin RAG) --- baseline mínimo. - \textbf{E5b}:
Capas 1 + 2 + 4 (sin SHAP) --- evalúa el aporte de SHAP al pipeline. -
\textbf{E5c}: Capas 1 + 2 + 3 + LLM libre (sin RAG) --- evalúa el aporte
del anclaje RAG. - \textbf{E5d}: Pipeline completo de 4 capas ---
referencia experimental.

\subsubsection{3.3.4 Pruebas estadísticas y mapa hipótesis →
experimento}\label{pruebas-estaduxedsticas-y-mapa-hipuxf3tesis-experimento}

Cada sub-hipótesis se contrasta con una prueba estadística específica,
seleccionada según el tipo de variable y el diseño:

\begin{longtable}[]{@{}
  >{\raggedright\arraybackslash}p{(\columnwidth - 10\tabcolsep) * \real{0.1667}}
  >{\raggedright\arraybackslash}p{(\columnwidth - 10\tabcolsep) * \real{0.1667}}
  >{\raggedright\arraybackslash}p{(\columnwidth - 10\tabcolsep) * \real{0.1667}}
  >{\raggedright\arraybackslash}p{(\columnwidth - 10\tabcolsep) * \real{0.1667}}
  >{\raggedright\arraybackslash}p{(\columnwidth - 10\tabcolsep) * \real{0.1667}}
  >{\raggedright\arraybackslash}p{(\columnwidth - 10\tabcolsep) * \real{0.1667}}@{}}
\toprule\noalign{}
\begin{minipage}[b]{\linewidth}\raggedright
Sub-hipótesis
\end{minipage} & \begin{minipage}[b]{\linewidth}\raggedright
Comparación
\end{minipage} & \begin{minipage}[b]{\linewidth}\raggedright
Variable
\end{minipage} & \begin{minipage}[b]{\linewidth}\raggedright
Prueba estadística
\end{minipage} & \begin{minipage}[b]{\linewidth}\raggedright
α
\end{minipage} & \begin{minipage}[b]{\linewidth}\raggedright
Tamaño de efecto
\end{minipage} \\
\midrule\noalign{}
\endhead
\bottomrule\noalign{}
\endlastfoot
H1a & Ensemble vs.~detector único & PR-AUC sobre 6 semillas & Wilcoxon
signed-rank (no paramétrica, apareada) & 0.05 & Hedges' g \\
H1b & SHAP vs.~sin SHAP & Likert comprensión (1--5) & Mann-Whitney U
(escala ordinal, muestras independientes) & 0.05 & r de rangos \\
H1c & RAG vs.~sin RAG & Rúbrica de reportes (1--5) & t de Student
apareado o Wilcoxon según Shapiro-Wilk & 0.05 & Cohen's d \\
H1d & Sistema integrado vs.~aislado & Tiempo-a-decisión (s) & t de
Student apareado (within-subjects) & 0.05 & Cohen's dz \\
\end{longtable}

Para todas las pruebas se verifica previamente el supuesto de normalidad
con Shapiro-Wilk; ante violación, se aplica la prueba no paramétrica
equivalente. Se reporta intervalo de confianza al 95\% para cada métrica
y se calcula el tamaño de efecto como medida complementaria a la
significancia estadística.

\subsubsection{3.3.5 Estudio de usabilidad: tamaño y selección de
muestra}\label{estudio-de-usabilidad-tamauxf1o-y-selecciuxf3n-de-muestra}

El estudio de usabilidad para VD4 adopta un \textbf{diseño
within-subjects con N ≥ 15 participantes}, contrabalanceado en orden de
presentación (mitad evalúa primero el sistema integrado, mitad evalúa
primero el aislado). Con este tamaño se reportan resultados como
exploratorios, con cálculo de tamaño de efecto Cohen's dz y intervalo de
confianza al 95\%, sin afirmar significancia estadística con potencia
plena. Para detectar un efecto medio (dz = 0.5) con potencia 0.80 y α =
0.05 se requieren N = 27 participantes; este tamaño se considera meta
deseable y, en caso de no alcanzarse, se reporta el limitante en §5.2
(Limitaciones).

\textbf{Criterios de inclusión de participantes}: - Estudiantes
avanzados de Ingeniería de Sistemas, Industrial o Agronomía (≥ séptimo
semestre), o - Profesionales con ≥ 1 año de experiencia en supervisión
operativa, control de calidad o auditoría interna.

\textbf{Criterios de exclusión}: - Participación previa en el diseño del
sistema o de cualquiera de sus capas. - Conflicto de interés directo con
empresas agroexportadoras evaluadas.

El protocolo detallado, formulario de consentimiento informado y
cuestionario post-tarea figuran en el Anexo A.

\subsubsection{3.3.6 Tuning de hiperparámetros y
reproducibilidad}\label{tuning-de-hiperparuxe1metros-y-reproducibilidad}

La selección de hiperparámetros se realiza con \textbf{Optuna} (TPE
sampler, 50 trials por modelo), optimizando PR-AUC sobre el validation
set. Los rangos de búsqueda se documentan en el Anexo B (Model Cards).
El código fuente, requirements.txt con versiones exactas, semillas y
notebooks de reproducción se publican en repositorio GitHub público con
licencia MIT al cierre del Hito 3.

\subsubsection{3.3.7 Comparación con
baselines}\label{comparaciuxf3n-con-baselines}

Para cada experimento, los resultados del sistema propuesto se comparan
con baselines documentados:

\begin{longtable}[]{@{}
  >{\raggedright\arraybackslash}p{(\columnwidth - 4\tabcolsep) * \real{0.3333}}
  >{\raggedright\arraybackslash}p{(\columnwidth - 4\tabcolsep) * \real{0.3333}}
  >{\raggedright\arraybackslash}p{(\columnwidth - 4\tabcolsep) * \real{0.3333}}@{}}
\toprule\noalign{}
\begin{minipage}[b]{\linewidth}\raggedright
\#
\end{minipage} & \begin{minipage}[b]{\linewidth}\raggedright
Baseline
\end{minipage} & \begin{minipage}[b]{\linewidth}\raggedright
Justificación
\end{minipage} \\
\midrule\noalign{}
\endhead
\bottomrule\noalign{}
\endlastfoot
B1 & Isolation Forest individual & Detector más simple y ampliamente
adoptado \\
B2 & Ensemble IF + LOF (sin ECOD) & Aislar el aporte de ECOD al
ensemble \\
B3 & XGBoost supervisado con etiqueta de anomalía & Upper bound
supervisado \\
B4 & LLM sin RAG y sin SHAP & Línea base de reporte automático \\
\end{longtable}

Los baselines se ejecutan sobre el mismo dataset, misma división y
mismas semillas para garantizar comparación justa.

\begin{center}\rule{0.5\linewidth}{0.5pt}\end{center}

\begin{center}\rule{0.5\linewidth}{0.5pt}\end{center}

\section{CAPÍTULO IV: RESULTADOS Y
DISCUSIÓN}\label{capuxedtulo-iv-resultados-y-discusiuxf3n}

\begin{quote}
\textbf{Estado:} Estructura completa lista --- las tablas y gráficos
numéricos se completarán tras ejecutar los experimentos E1--E5 (Hito 4
--- 2026-06-22).
\end{quote}

\subsection{4.1 Resultados Cuantitativos (Predicción y Detección ---
VD1)}\label{resultados-cuantitativos-predicciuxf3n-y-detecciuxf3n-vd1}

Esta sección presenta las métricas obtenidas por el módulo de predicción
tabular y el ensemble de detección de anomalías sobre el conjunto de
test del dataset sintético agroexportador (v1.0, 400 registros del
período 2025-05-01 a 2025-12-31).

\subsubsection{4.1.1 Tabla 4.1 --- Rendimiento de detección (Experimento
E1)}\label{tabla-4.1-rendimiento-de-detecciuxf3n-experimento-e1}

\begin{longtable}[]{@{}
  >{\raggedright\arraybackslash}p{(\columnwidth - 12\tabcolsep) * \real{0.1429}}
  >{\raggedright\arraybackslash}p{(\columnwidth - 12\tabcolsep) * \real{0.1429}}
  >{\raggedright\arraybackslash}p{(\columnwidth - 12\tabcolsep) * \real{0.1429}}
  >{\raggedright\arraybackslash}p{(\columnwidth - 12\tabcolsep) * \real{0.1429}}
  >{\raggedright\arraybackslash}p{(\columnwidth - 12\tabcolsep) * \real{0.1429}}
  >{\raggedright\arraybackslash}p{(\columnwidth - 12\tabcolsep) * \real{0.1429}}
  >{\raggedright\arraybackslash}p{(\columnwidth - 12\tabcolsep) * \real{0.1429}}@{}}
\toprule\noalign{}
\begin{minipage}[b]{\linewidth}\raggedright
Método
\end{minipage} & \begin{minipage}[b]{\linewidth}\raggedright
PR-AUC
\end{minipage} & \begin{minipage}[b]{\linewidth}\raggedright
ROC-AUC
\end{minipage} & \begin{minipage}[b]{\linewidth}\raggedright
F1
\end{minipage} & \begin{minipage}[b]{\linewidth}\raggedright
Precision
\end{minipage} & \begin{minipage}[b]{\linewidth}\raggedright
Recall
\end{minipage} & \begin{minipage}[b]{\linewidth}\raggedright
Tiempo inferencia
\end{minipage} \\
\midrule\noalign{}
\endhead
\bottomrule\noalign{}
\endlastfoot
Isolation Forest individual (baseline B1) & \emph{pendiente} &
\emph{pendiente} & \emph{pendiente} & \emph{pendiente} &
\emph{pendiente} & \emph{pendiente} \\
LOF individual & \emph{pendiente} & \emph{pendiente} & \emph{pendiente}
& \emph{pendiente} & \emph{pendiente} & \emph{pendiente} \\
ECOD individual & \emph{pendiente} & \emph{pendiente} & \emph{pendiente}
& \emph{pendiente} & \emph{pendiente} & \emph{pendiente} \\
Ensemble IF + LOF (B2) & \emph{pendiente} & \emph{pendiente} &
\emph{pendiente} & \emph{pendiente} & \emph{pendiente} &
\emph{pendiente} \\
\textbf{Ensemble IF + LOF + ECOD (propuesto)} & \emph{pendiente} &
\emph{pendiente} & \emph{pendiente} & \emph{pendiente} &
\emph{pendiente} & \emph{pendiente} \\
XGBoost supervisado (B3 --- upper bound) & \emph{pendiente} &
\emph{pendiente} & \emph{pendiente} & \emph{pendiente} &
\emph{pendiente} & \emph{pendiente} \\
\end{longtable}

\begin{quote}
Valores reportados como media ± desviación estándar sobre 6 semillas
(42--47).
\end{quote}

\subsubsection{4.1.2 Tabla 4.2 --- Recall por tipo de
anomalía}\label{tabla-4.2-recall-por-tipo-de-anomaluxeda}

\begin{longtable}[]{@{}
  >{\raggedright\arraybackslash}p{(\columnwidth - 6\tabcolsep) * \real{0.2500}}
  >{\raggedright\arraybackslash}p{(\columnwidth - 6\tabcolsep) * \real{0.2500}}
  >{\raggedright\arraybackslash}p{(\columnwidth - 6\tabcolsep) * \real{0.2500}}
  >{\raggedright\arraybackslash}p{(\columnwidth - 6\tabcolsep) * \real{0.2500}}@{}}
\toprule\noalign{}
\begin{minipage}[b]{\linewidth}\raggedright
Tipo de anomalía
\end{minipage} & \begin{minipage}[b]{\linewidth}\raggedright
Recall ensemble
\end{minipage} & \begin{minipage}[b]{\linewidth}\raggedright
Recall IF solo
\end{minipage} & \begin{minipage}[b]{\linewidth}\raggedright
Δ (puntos porcentuales)
\end{minipage} \\
\midrule\noalign{}
\endhead
\bottomrule\noalign{}
\endlastfoot
precio & \emph{pendiente} & \emph{pendiente} & \emph{pendiente} \\
volumen & \emph{pendiente} & \emph{pendiente} & \emph{pendiente} \\
clima & \emph{pendiente} & \emph{pendiente} & \emph{pendiente} \\
logistica & \emph{pendiente} & \emph{pendiente} & \emph{pendiente} \\
calidad & \emph{pendiente} & \emph{pendiente} & \emph{pendiente} \\
\end{longtable}

\subsection{4.2 Resultados Cualitativos (Explicabilidad y Reportes ---
VD2 y
VD3)}\label{resultados-cualitativos-explicabilidad-y-reportes-vd2-y-vd3}

\subsubsection{4.2.1 Tabla 4.3 --- Calidad de explicabilidad
(Experimento
E2)}\label{tabla-4.3-calidad-de-explicabilidad-experimento-e2}

\begin{longtable}[]{@{}
  >{\raggedright\arraybackslash}p{(\columnwidth - 6\tabcolsep) * \real{0.2500}}
  >{\raggedright\arraybackslash}p{(\columnwidth - 6\tabcolsep) * \real{0.2500}}
  >{\raggedright\arraybackslash}p{(\columnwidth - 6\tabcolsep) * \real{0.2500}}
  >{\raggedright\arraybackslash}p{(\columnwidth - 6\tabcolsep) * \real{0.2500}}@{}}
\toprule\noalign{}
\begin{minipage}[b]{\linewidth}\raggedright
Métrica
\end{minipage} & \begin{minipage}[b]{\linewidth}\raggedright
Sistema con SHAP
\end{minipage} & \begin{minipage}[b]{\linewidth}\raggedright
Sistema sin SHAP
\end{minipage} & \begin{minipage}[b]{\linewidth}\raggedright
p-value (Mann-Whitney U)
\end{minipage} \\
\midrule\noalign{}
\endhead
\bottomrule\noalign{}
\endlastfoot
Cobertura top-3 (mediana) & \emph{pendiente} & N/A & --- \\
Cobertura top-5 (mediana) & \emph{pendiente} & N/A & --- \\
Consistencia ρ (Spearman) & \emph{pendiente} & N/A & --- \\
Claridad operativa (Likert 1--5, promedio) & \emph{pendiente} &
\emph{pendiente} & \emph{pendiente} \\
\end{longtable}

\subsubsection{4.2.2 Tabla 4.4 --- Calidad de reportes generados
(Experimento
E3)}\label{tabla-4.4-calidad-de-reportes-generados-experimento-e3}

\begin{longtable}[]{@{}
  >{\raggedright\arraybackslash}p{(\columnwidth - 8\tabcolsep) * \real{0.2000}}
  >{\raggedright\arraybackslash}p{(\columnwidth - 8\tabcolsep) * \real{0.2000}}
  >{\raggedright\arraybackslash}p{(\columnwidth - 8\tabcolsep) * \real{0.2000}}
  >{\raggedright\arraybackslash}p{(\columnwidth - 8\tabcolsep) * \real{0.2000}}
  >{\raggedright\arraybackslash}p{(\columnwidth - 8\tabcolsep) * \real{0.2000}}@{}}
\toprule\noalign{}
\begin{minipage}[b]{\linewidth}\raggedright
Dimensión
\end{minipage} & \begin{minipage}[b]{\linewidth}\raggedright
LLM + RAG (propuesto)
\end{minipage} & \begin{minipage}[b]{\linewidth}\raggedright
LLM libre (control)
\end{minipage} & \begin{minipage}[b]{\linewidth}\raggedright
Kappa Cohen
\end{minipage} & \begin{minipage}[b]{\linewidth}\raggedright
p-value
\end{minipage} \\
\midrule\noalign{}
\endhead
\bottomrule\noalign{}
\endlastfoot
Completitud & \emph{pendiente} & \emph{pendiente} & \emph{pendiente} &
\emph{pendiente} \\
Consistencia numérica & \emph{pendiente} & \emph{pendiente} &
\emph{pendiente} & \emph{pendiente} \\
Accionabilidad & \emph{pendiente} & \emph{pendiente} & \emph{pendiente}
& \emph{pendiente} \\
Coherencia textual & \emph{pendiente} & \emph{pendiente} &
\emph{pendiente} & \emph{pendiente} \\
Correspondencia con evidencias & \emph{pendiente} & \emph{pendiente} &
\emph{pendiente} & \emph{pendiente} \\
\textbf{Promedio total} & \emph{pendiente} & \emph{pendiente} &
\emph{pendiente} & \emph{pendiente} \\
ROUGE-L (subset con referencia) & \emph{pendiente} & \emph{pendiente} &
--- & \emph{pendiente} \\
\end{longtable}

\subsubsection{4.2.3 Ejemplo de reporte generado (alerta tipo
``calidad'')}\label{ejemplo-de-reporte-generado-alerta-tipo-calidad}

\begin{quote}
Espacio reservado para insertar un reporte real de muestra que ilustre
el patrón anclado: dato → modelo → score → umbral → SHAP top-5 → fuente
RAG → recomendación operativa.
\end{quote}

\subsection{4.3 Resultados del Estudio de Usabilidad (VD4 y
VD5)}\label{resultados-del-estudio-de-usabilidad-vd4-y-vd5}

\subsubsection{4.3.1 Tabla 4.5 --- Tiempo-a-decisión y comprensión
(Experimento
E4)}\label{tabla-4.5-tiempo-a-decisiuxf3n-y-comprensiuxf3n-experimento-e4}

\begin{longtable}[]{@{}
  >{\raggedright\arraybackslash}p{(\columnwidth - 10\tabcolsep) * \real{0.1667}}
  >{\raggedright\arraybackslash}p{(\columnwidth - 10\tabcolsep) * \real{0.1667}}
  >{\raggedright\arraybackslash}p{(\columnwidth - 10\tabcolsep) * \real{0.1667}}
  >{\raggedright\arraybackslash}p{(\columnwidth - 10\tabcolsep) * \real{0.1667}}
  >{\raggedright\arraybackslash}p{(\columnwidth - 10\tabcolsep) * \real{0.1667}}
  >{\raggedright\arraybackslash}p{(\columnwidth - 10\tabcolsep) * \real{0.1667}}@{}}
\toprule\noalign{}
\begin{minipage}[b]{\linewidth}\raggedright
Métrica
\end{minipage} & \begin{minipage}[b]{\linewidth}\raggedright
Sistema integrado
\end{minipage} & \begin{minipage}[b]{\linewidth}\raggedright
Componentes aislados
\end{minipage} & \begin{minipage}[b]{\linewidth}\raggedright
Δ relativo
\end{minipage} & \begin{minipage}[b]{\linewidth}\raggedright
p-value
\end{minipage} & \begin{minipage}[b]{\linewidth}\raggedright
Cohen's dz
\end{minipage} \\
\midrule\noalign{}
\endhead
\bottomrule\noalign{}
\endlastfoot
Tiempo-a-decisión (s, mediana) & \emph{pendiente} & \emph{pendiente} &
\emph{pendiente} & \emph{pendiente} & \emph{pendiente} \\
Comprensión (Likert 1--5, media) & \emph{pendiente} & \emph{pendiente} &
\emph{pendiente} & \emph{pendiente} & \emph{pendiente} \\
Decisión correcta (\%) & \emph{pendiente} & \emph{pendiente} &
\emph{pendiente} & \emph{pendiente} & --- \\
SUS Score (0--100) & \emph{pendiente} & \emph{pendiente} &
\emph{pendiente} & \emph{pendiente} & \emph{pendiente} \\
\end{longtable}

\begin{quote}
N participantes = \emph{pendiente} --- diseño within-subject
contrabalanceado.
\end{quote}

\subsubsection{4.3.2 Tabla 4.6 --- Trazabilidad documental
(VD5)}\label{tabla-4.6-trazabilidad-documental-vd5}

\begin{longtable}[]{@{}
  >{\raggedright\arraybackslash}p{(\columnwidth - 4\tabcolsep) * \real{0.3333}}
  >{\raggedright\arraybackslash}p{(\columnwidth - 4\tabcolsep) * \real{0.3333}}
  >{\raggedright\arraybackslash}p{(\columnwidth - 4\tabcolsep) * \real{0.3333}}@{}}
\toprule\noalign{}
\begin{minipage}[b]{\linewidth}\raggedright
Configuración
\end{minipage} & \begin{minipage}[b]{\linewidth}\raggedright
\% alertas completas
\end{minipage} & \begin{minipage}[b]{\linewidth}\raggedright
Campos faltantes más frecuentes
\end{minipage} \\
\midrule\noalign{}
\endhead
\bottomrule\noalign{}
\endlastfoot
Sistema integrado (E5d) & \emph{pendiente} (meta ≥95\%) &
\emph{pendiente} \\
Ablation sin SHAP (E5b) & \emph{pendiente} & \emph{pendiente} \\
Ablation sin RAG (E5c) & \emph{pendiente} & \emph{pendiente} \\
Componentes aislados (control) & \emph{pendiente} (esperado \textless{}
30\%) & \emph{pendiente} \\
\end{longtable}

\subsubsection{4.3.3 Tabla 4.7 --- Ablation study (Experimento
E5)}\label{tabla-4.7-ablation-study-experimento-e5}

\begin{longtable}[]{@{}
  >{\raggedright\arraybackslash}p{(\columnwidth - 14\tabcolsep) * \real{0.1250}}
  >{\raggedright\arraybackslash}p{(\columnwidth - 14\tabcolsep) * \real{0.1250}}
  >{\raggedright\arraybackslash}p{(\columnwidth - 14\tabcolsep) * \real{0.1250}}
  >{\raggedright\arraybackslash}p{(\columnwidth - 14\tabcolsep) * \real{0.1250}}
  >{\raggedright\arraybackslash}p{(\columnwidth - 14\tabcolsep) * \real{0.1250}}
  >{\raggedright\arraybackslash}p{(\columnwidth - 14\tabcolsep) * \real{0.1250}}
  >{\raggedright\arraybackslash}p{(\columnwidth - 14\tabcolsep) * \real{0.1250}}
  >{\raggedright\arraybackslash}p{(\columnwidth - 14\tabcolsep) * \real{0.1250}}@{}}
\toprule\noalign{}
\begin{minipage}[b]{\linewidth}\raggedright
Configuración
\end{minipage} & \begin{minipage}[b]{\linewidth}\raggedright
Capa 1
\end{minipage} & \begin{minipage}[b]{\linewidth}\raggedright
Capa 2
\end{minipage} & \begin{minipage}[b]{\linewidth}\raggedright
Capa 3
\end{minipage} & \begin{minipage}[b]{\linewidth}\raggedright
Capa 4
\end{minipage} & \begin{minipage}[b]{\linewidth}\raggedright
PR-AUC
\end{minipage} & \begin{minipage}[b]{\linewidth}\raggedright
Trazabilidad \%
\end{minipage} & \begin{minipage}[b]{\linewidth}\raggedright
Likert claridad
\end{minipage} \\
\midrule\noalign{}
\endhead
\bottomrule\noalign{}
\endlastfoot
E5a --- solo detección & ✗ & ✓ & ✗ & ✗ & \emph{pendiente} &
\emph{pendiente} & \emph{pendiente} \\
E5b --- sin SHAP & ✓ & ✓ & ✗ & ✓ & \emph{pendiente} & \emph{pendiente} &
\emph{pendiente} \\
E5c --- sin RAG & ✓ & ✓ & ✓ & LLM libre & \emph{pendiente} &
\emph{pendiente} & \emph{pendiente} \\
\textbf{E5d --- pipeline completo} & ✓ & ✓ & ✓ & ✓ & \emph{pendiente} &
\emph{pendiente} & \emph{pendiente} \\
\end{longtable}

\begin{center}\rule{0.5\linewidth}{0.5pt}\end{center}

\subsection{4.4 Discusión Detallada y Cruce de Datos
Comparativo}\label{discusiuxf3n-detallada-y-cruce-de-datos-comparativo}

\subsubsection{4.4.1 Propósito de la
discusión}\label{propuxf3sito-de-la-discusiuxf3n}

Esta sección triangula los resultados experimentales propios con tres
bloques externos: (a) la literatura comparable revisada en el Capítulo
II §2.2, (b) las hipótesis declaradas en el Capítulo I §1.4, y (c) las
variables operacionalizadas formalmente en
\texttt{variables-operacionalizadas.md}. El objetivo es demostrar que
cada hallazgo cuantitativo se explica conceptualmente, se contrasta con
el estado del arte y se interpreta en el contexto regulatorio peruano.

\subsubsection{4.4.2 Cruce 1 --- Resultados propios versus literatura
comparable (Tabla
4.8)}\label{cruce-1-resultados-propios-versus-literatura-comparable-tabla-4.8}

La Tabla 4.8 sitúa el rendimiento del sistema propuesto frente a los
cinco trabajos más cercanos identificados en la búsqueda sistemática
(\texttt{busqueda-sistematica-gap.md}). La comparación se realiza
considerando que cada trabajo opera en un dominio y dataset distintos,
por lo que los valores absolutos no son directamente comparables; lo que
se compara es la cobertura modular y la consistencia direccional de los
resultados.

\begin{longtable}[]{@{}
  >{\raggedright\arraybackslash}p{(\columnwidth - 12\tabcolsep) * \real{0.1429}}
  >{\raggedright\arraybackslash}p{(\columnwidth - 12\tabcolsep) * \real{0.1429}}
  >{\raggedright\arraybackslash}p{(\columnwidth - 12\tabcolsep) * \real{0.1429}}
  >{\raggedright\arraybackslash}p{(\columnwidth - 12\tabcolsep) * \real{0.1429}}
  >{\raggedright\arraybackslash}p{(\columnwidth - 12\tabcolsep) * \real{0.1429}}
  >{\raggedright\arraybackslash}p{(\columnwidth - 12\tabcolsep) * \real{0.1429}}
  >{\raggedright\arraybackslash}p{(\columnwidth - 12\tabcolsep) * \real{0.1429}}@{}}
\toprule\noalign{}
\begin{minipage}[b]{\linewidth}\raggedright
Atributo
\end{minipage} & \begin{minipage}[b]{\linewidth}\raggedright
Esta tesis
\end{minipage} & \begin{minipage}[b]{\linewidth}\raggedright
AuditCopilot (Kadir et al., 2025)
\end{minipage} & \begin{minipage}[b]{\linewidth}\raggedright
Park (2024)
\end{minipage} & \begin{minipage}[b]{\linewidth}\raggedright
JRFM (2025)
\end{minipage} & \begin{minipage}[b]{\linewidth}\raggedright
Almalki \& Masud (2025)
\end{minipage} & \begin{minipage}[b]{\linewidth}\raggedright
AuditMAI (Waltersdorfer et al., 2024)
\end{minipage} \\
\midrule\noalign{}
\endhead
\bottomrule\noalign{}
\endlastfoot
Predicción tabular GBDT & XGBoost + LightGBM & No reportada & No (solo
LLMs) & Stacking GBDT & Stacking GBDT & No \\
Ensemble de detección & IF + LOF + ECOD & Parcial & No & No & No & No \\
Explicabilidad SHAP estructurada & TreeSHAP top-k & No & No & SHAP &
SHAP & No \\
Generación LLM bajo RAG & RAG anclado en SHAP & LLM libre & Multi-agente
& No & No & No \\
Restricción anti-alucinación & Sí (validación numérica posterior) & No
documentada & No & N/A & N/A & N/A \\
PR-AUC reportado & \emph{pendiente} & No reporta & No reporta & 0.93 &
\emph{pendiente} verificar & No \\
Dominio & Agroexportador peruano & Asientos contables & S\&P 500 &
Fraude financiero & Fraude financiero & Auditoría de IA \\
Contexto regulatorio & D.S. 115-2025-PCM + NIST AI RMF & No especificado
& No & No & No & No \\
Evaluación con usuarios & Sí (N ≥ 15) & No & No & No & No & No \\
Dataset abierto disponible & Sí (CC BY 4.0) & No & No & No & Parcial &
N/A \\
\end{longtable}

\textbf{Lectura del cruce 1}: - La cobertura modular del sistema
propuesto (cuatro capas con restricción anti-alucinación) es
estrictamente mayor que cualquier trabajo individual de la literatura
revisada. - El trabajo más cercano en cobertura es JRFM (2025) que
combina GBDT y SHAP en fraude financiero, pero carece de detección no
supervisada, módulo LLM y contexto regulatorio. - AuditCopilot (Kadir et
al., 2025) es el único que combina detección con generación LLM, pero
opera en asientos contables sin SHAP estructurado ni restricción
anti-alucinación.

\subsubsection{4.4.3 Cruce 2 --- Contraste de hipótesis (Tabla
4.9)}\label{cruce-2-contraste-de-hipuxf3tesis-tabla-4.9}

\begin{longtable}[]{@{}
  >{\raggedright\arraybackslash}p{(\columnwidth - 6\tabcolsep) * \real{0.2500}}
  >{\raggedright\arraybackslash}p{(\columnwidth - 6\tabcolsep) * \real{0.2500}}
  >{\raggedright\arraybackslash}p{(\columnwidth - 6\tabcolsep) * \real{0.2500}}
  >{\raggedright\arraybackslash}p{(\columnwidth - 6\tabcolsep) * \real{0.2500}}@{}}
\toprule\noalign{}
\begin{minipage}[b]{\linewidth}\raggedright
Sub-hipótesis
\end{minipage} & \begin{minipage}[b]{\linewidth}\raggedright
Predicción del Capítulo I
\end{minipage} & \begin{minipage}[b]{\linewidth}\raggedright
Evidencia empírica obtenida
\end{minipage} & \begin{minipage}[b]{\linewidth}\raggedright
Decisión
\end{minipage} \\
\midrule\noalign{}
\endhead
\bottomrule\noalign{}
\endlastfoot
H1a --- Ensemble supera al detector individual & PR-AUC ensemble
\textgreater{} PR-AUC IF con tamaño de efecto Hedges' g ≥ 0.5 &
\emph{pendiente E1} & \emph{Aceptar / Rechazar / Inconcluso} \\
H1b --- SHAP mejora la comprensión & Likert claridad ≥ 4.0 con SHAP;
\textless{} 4.0 sin SHAP & \emph{pendiente E2} & \emph{pendiente} \\
H1c --- RAG mejora la calidad del reporte & Rúbrica promedio ≥ 4.0 y
Kappa Cohen ≥ 0.60 & \emph{pendiente E3} & \emph{pendiente} \\
H1d --- Sistema integrado reduce tiempo-a-decisión & Reducción ≥ 20\%
del tiempo, p \textless{} 0.05 & \emph{pendiente E4} &
\emph{pendiente} \\
H1 (general) --- Sistema integrado mejora trazabilidad & ≥ 95\% alertas
completas en integrado vs.~\textless{} 30\% en aislado & \emph{pendiente
E4 + E5} & \emph{pendiente} \\
\end{longtable}

\textbf{Lectura del cruce 2}: Esta tabla materializa la
operacionalización del Capítulo I y permite al jurado verificar que cada
hipótesis está contrastada empíricamente. La columna ``Decisión'' se
completa con: (a) ``Aceptar H1x'' si p \textless{} 0.05 y dirección
esperada, (b) ``Rechazar'' si p \textless{} 0.05 en dirección opuesta,
(c) ``Inconcluso'' si p ≥ 0.05.

\subsubsection{4.4.4 Cruce 3 --- Variables operacionalizadas versus
indicadores observados (Tabla
4.10)}\label{cruce-3-variables-operacionalizadas-versus-indicadores-observados-tabla-4.10}

\begin{longtable}[]{@{}
  >{\raggedright\arraybackslash}p{(\columnwidth - 6\tabcolsep) * \real{0.2500}}
  >{\raggedright\arraybackslash}p{(\columnwidth - 6\tabcolsep) * \real{0.2500}}
  >{\raggedright\arraybackslash}p{(\columnwidth - 6\tabcolsep) * \real{0.2500}}
  >{\raggedright\arraybackslash}p{(\columnwidth - 6\tabcolsep) * \real{0.2500}}@{}}
\toprule\noalign{}
\begin{minipage}[b]{\linewidth}\raggedright
Variable
\end{minipage} & \begin{minipage}[b]{\linewidth}\raggedright
Criterio de aceptación declarado
\end{minipage} & \begin{minipage}[b]{\linewidth}\raggedright
Valor observado
\end{minipage} & \begin{minipage}[b]{\linewidth}\raggedright
Cumple
\end{minipage} \\
\midrule\noalign{}
\endhead
\bottomrule\noalign{}
\endlastfoot
VD1 --- Rendimiento de detección & PR-AUC integrado \textgreater{}
PR-AUC B1 (p\textless0.05, g≥0.5) & \emph{pendiente} &
\emph{pendiente} \\
VD2 --- Calidad de explicabilidad & Cobertura top-5 ≥ 80\% en ≥70\% de
alertas; Likert ≥ 4.0 & \emph{pendiente} & \emph{pendiente} \\
VD3 --- Calidad de reportes & Promedio rúbrica ≥ 4.0/5; Kappa ≥ 0.60;
ROUGE-L ≥ 0.40 & \emph{pendiente} & \emph{pendiente} \\
VD4 --- Comprensión y tiempo & Reducción ≥ 20\% tiempo; Likert ≥ 4.0;
tasa correcta ≥ 0.80 & \emph{pendiente} & \emph{pendiente} \\
VD5 --- Trazabilidad & ≥ 95\% alertas con 8 campos completos &
\emph{pendiente} & \emph{pendiente} \\
\end{longtable}

\subsubsection{4.4.5 Cruce 4 --- Mapa principio regulatorio → componente
arquitectónico → métrica observada (Tabla
4.11)}\label{cruce-4-mapa-principio-regulatorio-componente-arquitectuxf3nico-muxe9trica-observada-tabla-4.11}

\begin{longtable}[]{@{}
  >{\raggedright\arraybackslash}p{(\columnwidth - 4\tabcolsep) * \real{0.3333}}
  >{\raggedright\arraybackslash}p{(\columnwidth - 4\tabcolsep) * \real{0.3333}}
  >{\raggedright\arraybackslash}p{(\columnwidth - 4\tabcolsep) * \real{0.3333}}@{}}
\toprule\noalign{}
\begin{minipage}[b]{\linewidth}\raggedright
Principio (D.S. 115-2025-PCM / NIST AI RMF / EU AI Act Art. 13)
\end{minipage} & \begin{minipage}[b]{\linewidth}\raggedright
Componente arquitectónico responsable
\end{minipage} & \begin{minipage}[b]{\linewidth}\raggedright
Métrica observada
\end{minipage} \\
\midrule\noalign{}
\endhead
\bottomrule\noalign{}
\endlastfoot
Transparencia & Capa 3 SHAP + Anexo B Model Cards & VD2 cobertura top-k
+ claridad Likert \\
Explicabilidad & TreeSHAP top-5 + reporte LLM+RAG & VD2 + VD3 \\
Supervisión humana & Revisión obligatoria del reporte antes de uso &
Cuestionario VD4 SUS pregunta 9 (confianza) \\
Gestión de riesgos & Umbrales calibrados + validación cruzada por
semilla & VD1 PR-AUC media ± DE \\
Documentación & Datasheet (Anexo C) + Model Cards (Anexo B) & Logs de
versión + Anexo D \\
Trazabilidad & Estructura de 8 campos por alerta & VD5 \% alertas
completas \\
Anti-alucinación & RAG anclado en SHAP + validación numérica posterior &
VD3 dimensión consistencia numérica \\
\end{longtable}

\textbf{Lectura del cruce 4}: Esta tabla evidencia que cada principio
regulatorio relevante tiene un componente concreto que lo materializa y
una métrica empírica que lo verifica. Es la evidencia operativa de la
``conformidad de diseño'' declarada en §2.3.8.

\subsubsection{4.4.6 Cruce 5 --- Errores por tipo de anomalía y posibles
causas (Tabla
4.12)}\label{cruce-5-errores-por-tipo-de-anomaluxeda-y-posibles-causas-tabla-4.12}

\begin{longtable}[]{@{}
  >{\raggedright\arraybackslash}p{(\columnwidth - 6\tabcolsep) * \real{0.2500}}
  >{\raggedright\arraybackslash}p{(\columnwidth - 6\tabcolsep) * \real{0.2500}}
  >{\raggedright\arraybackslash}p{(\columnwidth - 6\tabcolsep) * \real{0.2500}}
  >{\raggedright\arraybackslash}p{(\columnwidth - 6\tabcolsep) * \real{0.2500}}@{}}
\toprule\noalign{}
\begin{minipage}[b]{\linewidth}\raggedright
Tipo anomalía
\end{minipage} & \begin{minipage}[b]{\linewidth}\raggedright
Recall obs.
\end{minipage} & \begin{minipage}[b]{\linewidth}\raggedright
Mecanismo probable de fallos
\end{minipage} & \begin{minipage}[b]{\linewidth}\raggedright
Recomendación de mejora
\end{minipage} \\
\midrule\noalign{}
\endhead
\bottomrule\noalign{}
\endlastfoot
precio & \emph{pendiente} & Outliers de cola gruesa pueden caer dentro
del rango plausible si la variación es estacional & Incorporar
covariable mes y desviación respecto a la media móvil del producto \\
volumen & \emph{pendiente} & Posible confusión con eventos
extraordinarios reales (campaña pico) & Añadir feature
\texttt{dia\_pico\_campania} \\
clima & \emph{pendiente} & LOF sensible a vecindario; pocas
observaciones con sequía + calor extremo & Agregar densidad estacional
regional \\
logistica & \emph{pendiente} & Anomalía compuesta (dos condiciones) ---
IF puede subestimar & Reforzar con regla determinista complementaria \\
calidad & \emph{pendiente} & Pocas instancias inyectadas (10\%) & Subir
a 15\% en v1.1 si recall \textless{} 0.6 \\
\end{longtable}

\subsubsection{4.4.7 Interpretación
conjunta}\label{interpretaciuxf3n-conjunta}

\begin{quote}
\textbf{Lectura integrada (a completar con resultados reales)}: Si los
cinco cruces (Tabla 4.8--4.12) muestran consistencia direccional con las
predicciones del Capítulo I, esta tesis sostiene que la arquitectura
propuesta no solo es novedosa por su integración modular, sino también
empíricamente superior bajo las condiciones evaluadas. En caso de
inconsistencias, se documentan en §4.5 (Limitaciones de los resultados)
y se proponen líneas de mejora en §5.3 (Trabajo futuro).
\end{quote}

\subsection{4.5 Limitaciones de los
Resultados}\label{limitaciones-de-los-resultados}

Los resultados de este capítulo deben interpretarse considerando:

\begin{enumerate}
\def\labelenumi{\arabic{enumi}.}
\tightlist
\item
  \textbf{Naturaleza sintética del dataset}: las distribuciones reflejan
  rangos plausibles documentados, pero no replican la variabilidad
  operativa real de empresas agroexportadoras (ver §1.12.1).
\item
  \textbf{Tamaño de muestra del estudio de usabilidad}: N ≥ 15 permite
  reportar resultados exploratorios; para conclusiones con potencia 0.80
  se requiere N = 27.
\item
  \textbf{Variabilidad estocástica del LLM}: cada reporte se genera con
  temperature = 0.2 y se reporta como promedio sobre 3 generaciones;
  sigue siendo sensible a cambios de versión del modelo base.
\item
  \textbf{Comparación con literatura}: los valores absolutos no son
  directamente comparables entre trabajos por diferencias de dominio y
  dataset; se contrasta cobertura modular y direccionalidad.
\end{enumerate}

\subsection{4.6 Síntesis del Capítulo
IV}\label{suxedntesis-del-capuxedtulo-iv}

\begin{enumerate}
\def\labelenumi{\arabic{enumi}.}
\tightlist
\item
  El ensemble IF + LOF + ECOD \emph{supera/no supera} al detector
  individual con tamaño de efecto \emph{g}=\emph{pendiente} (H1a).
\item
  SHAP \emph{mejora/no mejora} la comprensión de las alertas (H1b).
\item
  RAG anclado \emph{mejora/no mejora} la calidad de los reportes (H1c).
\item
  El sistema integrado \emph{reduce/no reduce} el tiempo-a-decisión
  (H1d).
\item
  La trazabilidad documental alcanza \emph{XX}\% de alertas completas en
  la condición integrada.
\item
  El cruce con literatura confirma que la integración de los cuatro
  módulos con restricción anti-alucinación es una contribución original
  verificable.
\end{enumerate}

\begin{center}\rule{0.5\linewidth}{0.5pt}\end{center}

\begin{center}\rule{0.5\linewidth}{0.5pt}\end{center}

\section{CAPÍTULO V: CONCLUSIONES Y TRABAJOS
FUTUROS}\label{capuxedtulo-v-conclusiones-y-trabajos-futuros}

\begin{quote}
\textbf{Pendiente:} Capítulo V --- depende de los resultados del
Capítulo IV
\end{quote}

\subsection{5.1 Conclusiones}\label{conclusiones}

\emph{(Esqueleto para la síntesis final: el sistema integrado logró los
objetivos propuestos, manteniendo el balance entre vanguardia
tecnológica y rigor legal. Incluir: conclusión sobre el gap cerrado,
métricas alcanzadas vs.~objetivos, validación de hipótesis H1a--H1d,
aporte al contexto regulatorio peruano).}

\subsection{5.2 Limitaciones de la
Investigación}\label{limitaciones-de-la-investigaciuxf3n-1}

\emph{(Abordar: limitaciones del dataset sintético agroexportador;
granularidad y disponibilidad de fuentes públicas; dependencia de la
calidad de datos documentada en Datasheets for Datasets (Gebru et al.,
2021); deuda técnica de mantenimiento del pipeline MLOps (Sculley et
al., 2015); limitaciones del tamaño de la muestra en la evaluación de
comprensión; restricciones de los LLMs actuales en precisión de cálculo
aritmético (Maynez et al., 2026)).}

\subsection{5.3 Trabajos Futuros}\label{trabajos-futuros}

\emph{(Propuestas: integración de GraphRAG para recuperación semántica
más rica sobre conocimiento agroexportador; extensión del ensemble con
ECOD (Li et al., 2022) y modelos de concept drift para supervisión en
stream; exploración de Chronos (Ansari et al., 2024) para forecasting de
horizonte largo; prueba piloto en una empresa agroexportadora peruana;
evaluación de sesgos y limitaciones según Datasheets for Datasets (Gebru
et al., 2021)).}

\begin{center}\rule{0.5\linewidth}{0.5pt}\end{center}

\begin{center}\rule{0.5\linewidth}{0.5pt}\end{center}

\section{CRONOGRAMA DE ACTIVIDADES}\label{cronograma-de-actividades-1}

\begin{quote}
\textbf{Pendiente:} Conclusiones/Conclusions --- depende de Cap IV y V
\end{quote}

\emph{(Ver Tabla 1.2 en el Capítulo I, §1.11)}

\begin{center}\rule{0.5\linewidth}{0.5pt}\end{center}

\section{CONCLUSIONES}\label{conclusiones-1}

\emph{(Por completar con los resultados finales de la investigación.
Estructura sugerida:)}

\begin{enumerate}
\def\labelenumi{\arabic{enumi}.}
\item
  \emph{(Conclusión sobre el gap cerrado: el sistema integrado de cuatro
  capas constituye la primera propuesta académica en el Perú que unifica
  predicción GBDT, detección de anomalías ensemble, explicabilidad SHAP
  y generación de reportes LLM+RAG con trazabilidad regulatoria
  verificable para el contexto agroexportador.)}
\item
  \emph{(Conclusión sobre métricas alcanzadas: los objetivos específicos
  OE1--OE5 fueron alcanzados según las métricas definidas --- PR-AUC ≥
  0.92, SHAP Coverage ≥ 70\%, ROUGE-1 ≥ 0.50, reducción
  tiempo-a-decisión ≥ 30\%\ldots)}
\item
  \emph{(Conclusión sobre validación de hipótesis H1a--H1d.)}
\item
  \emph{(Conclusión sobre gobernanza: el diseño incorpora principios de
  trazabilidad, documentación y supervisión humana alineados con el D.S.
  N° 115-2025-PCM, NIST AI RMF y buenas prácticas de gestión de riesgo
  de modelos.)}
\item
  \emph{(Conclusión sobre aporte al campo: el sistema integrado
  proporciona mayor trazabilidad y usabilidad que los sistemas con
  componentes aislados, validado mediante experimentos comparativos
  controlados.)}
\end{enumerate}

\begin{center}\rule{0.5\linewidth}{0.5pt}\end{center}

\section{CONCLUSIONS}\label{conclusions}

\emph{(Por completar --- versión en inglés de las conclusiones
principales)}

\begin{enumerate}
\def\labelenumi{\arabic{enumi}.}
\item
  \emph{(The integrated four-layer system constitutes the first academic
  proposal in Peru that unifies GBDT prediction, ensemble anomaly
  detection, SHAP explainability, and LLM+RAG report generation with
  verifiable regulatory traceability.)}
\item
  \emph{(Technical objectives OE1--OE5 were achieved according to
  defined metrics --- PR-AUC ≥ 0.92, SHAP Coverage ≥ 70\%, ROUGE-1 ≥
  0.50, decision-time reduction ≥ 30\%.)}
\item
  \emph{(The integrated system produces superior traceability and
  usability compared to isolated-component systems, as validated through
  controlled comparative experiments.)}
\item
  \emph{(The design incorporates traceability, documentation, and human
  oversight principles aligned with D.S. N° 115-2025-PCM, NIST AI RMF,
  and model risk management good practices.)}
\end{enumerate}

\begin{center}\rule{0.5\linewidth}{0.5pt}\end{center}

\begin{center}\rule{0.5\linewidth}{0.5pt}\end{center}

\section{RECOMENDACIONES}\label{recomendaciones}

\begin{enumerate}
\def\labelenumi{\arabic{enumi}.}
\item
  \textbf{Para implementadores}: Se recomienda iniciar el despliegue del
  sistema con el módulo de predicción GBDT y el módulo de explicabilidad
  SHAP antes de integrar el componente LLM+RAG, siguiendo el principio
  de implementación incremental que reduce la deuda técnica (Sculley et
  al., 2015) y permite validar cada capa de forma independiente.
\item
  \textbf{Para empresas agroexportadoras}: Antes de adoptar el sistema
  en producción, se recomienda elaborar Datasheets for Datasets (Gebru
  et al., 2021) para todos los datasets de entrenamiento y Model Cards
  (Mitchell et al., 2019) para los modelos XGBoost, detectores de
  anomalías y LLM+RAG.
\item
  \textbf{Para futuros investigadores}: Se recomienda extender la
  evaluación del sistema con un diseño experimental longitudinal que
  capture el efecto del concept drift en precios, volúmenes, clima y
  comportamiento exportador, utilizando ventanas temporales y fuentes
  agroexportadoras reales.
\item
  \textbf{Para entidades públicas y sectoriales}: Se recomienda promover
  guías técnicas de IA explicable y trazabilidad para sistemas de
  supervisión en cadenas productivas, tomando como referencia marcos
  nacionales e internacionales de gobernanza de IA.
\item
  \textbf{Para la academia}: Se recomienda replicar el estudio con datos
  reales de una empresa agroexportadora colaboradora (bajo acuerdo de
  confidencialidad), ampliar la muestra de evaluación con supervisores
  operativos y responsables de calidad, e incorporar métricas de sesgo y
  robustez según las dimensiones de evaluación de ADBench (Han et al.,
  2022).
\end{enumerate}

\begin{center}\rule{0.5\linewidth}{0.5pt}\end{center}

\begin{center}\rule{0.5\linewidth}{0.5pt}\end{center}

\section{GLOSARIO DE TÉRMINOS}\label{glosario-de-tuxe9rminos}

\textbf{ADBench} (\emph{Anomaly Detection Benchmark}): Benchmark
sistemático para evaluación comparativa de algoritmos de detección de
anomalías, propuesto por Han et al.~(2022), que cubre 57 datasets y 30
algoritmos bajo tres niveles de supervisión.

\textbf{Alucinación (LLM)}: Fenómeno en el que un modelo de lenguaje
genera texto coherente en forma pero incorrecto en contenido, incluyendo
afirmaciones factuales erróneas, citas inexistentes o cifras fabricadas.

\textbf{AUC-PR} (\emph{Area Under the Precision-Recall Curve}): Métrica
de evaluación para clasificadores en datasets desbalanceados; a
diferencia de AUC-ROC, es sensible a la distribución de clases y
penaliza los falsos positivos de forma más relevante en contextos de
fraude y anomalías raras.

\textbf{BAF Benchmark} (\emph{Bank Account Fraud}): Dataset tabular de
referencia para fraude bancario con drift temporal y desbalance de
clases, publicado por Jesus et al.~(2022). En esta tesis se considera
únicamente como benchmark metodológico complementario para evaluar
robustez tabular, no como validación directa del dominio agroexportador.

\textbf{CatBoost}: Algoritmo GBDT desarrollado por Yandex (Prokhorenkova
et al., 2018) que resuelve el problema de target leakage en variables
categóricas mediante Ordered Boosting.

\textbf{Concept Drift}: Cambio en la distribución estadística de los
datos a lo largo del tiempo que degrada el rendimiento de modelos
entrenados con datos históricos; particularmente relevante en detección
de fraude y anomalías operativas.

\textbf{Deep SVDD} (\emph{Deep Support Vector Data Description}): Método
de detección de anomalías basado en redes neuronales que aprende una
hipersfera mínima en el espacio latente que contiene los datos normales
(Ruff et al., 2018).

\textbf{ECOD} (\emph{Empirical Cumulative distribution functions Outlier
Detection}): Detector de anomalías sin parámetros basado en
distribuciones empíricas acumuladas (Li et al., 2022), notable por su
ausencia de hiperparámetros a calibrar.

\textbf{EU AI Act} (\emph{Reglamento (UE) 2024/1689}): Reglamento
europeo de inteligencia artificial que clasifica los sistemas de IA por
nivel de riesgo y establece obligaciones de transparencia y
explicabilidad en el Artículo 13 para sistemas de alto riesgo.

\textbf{GBDT} (\emph{Gradient Boosting Decision Trees}): Familia de
algoritmos de aprendizaje supervisado que construyen modelos predictivos
mediante la combinación secuencial de árboles de decisión débiles,
minimizando una función de pérdida mediante descenso de gradiente
funcional.

\textbf{Isolation Forest}: Algoritmo de detección de anomalías no
supervisado (Liu et al., 2008) que aísla anomalías mediante
particionamiento aleatorio del espacio de datos, con complejidad
computacional O(n).

\textbf{LightGBM}: Algoritmo GBDT de Microsoft (Ke et al., 2017) que
acelera el entrenamiento hasta 20× mediante Gradient-based One-Side
Sampling y estructuras de datos basadas en histogramas.

\textbf{LLM} (\emph{Large Language Model}): Modelo de lenguaje de gran
tamaño entrenado en corpus masivos de texto para aprender distribuciones
probabilísticas sobre secuencias de tokens, capaz de realizar tareas de
generación, resumen y razonamiento en lenguaje natural.

\textbf{LOF} (\emph{Local Outlier Factor}): Detector de anomalías basado
en densidad local relativa (Breunig et al., 2000) que cuantifica el
grado de anomalía de un punto comparando su densidad con la de sus
vecinos k-NN.

\textbf{MLOps} (\emph{Machine Learning Operations}): Conjunto de
prácticas para gestionar el ciclo de vida completo de modelos ML en
producción, incluyendo CI/CD, monitoreo de drift, automatización de
reentrenamiento y trazabilidad de versiones.

\textbf{NIST AI RMF}: Framework de gestión de riesgo para sistemas de IA
publicado por el Instituto Nacional de Estándares y Tecnología de EE.UU.
(2023), organizado en cuatro funciones: Govern, Map, Measure y Manage.

\textbf{PR-AUC}: Ver AUC-PR.

\textbf{PyOD} (\emph{Python Outlier Detection}): Librería de Python que
implementa más de 40 algoritmos de detección de outliers con una API
estandarizada compatible con scikit-learn (Zhao et al., 2019).

\textbf{RAG} (\emph{Retrieval-Augmented Generation}): Arquitectura para
modelos de lenguaje que separa el conocimiento factual del modelo
generativo, recuperando documentos relevantes de una base de
conocimiento externa para anclar las respuestas y mitigar alucinaciones.

\textbf{Resolución SBS N° 053-2023}: Resolución de la Superintendencia
de Banca, Seguros y AFP del Perú que establece lineamientos de gestión
de riesgos de modelos para entidades supervisadas. En esta tesis se
utiliza como referencia nacional de buenas prácticas para trazabilidad,
validación y monitoreo, no como obligación directa para empresas
agroexportadoras.

\textbf{ROUGE} (\emph{Recall-Oriented Understudy for Gisting
Evaluation}): Conjunto de métricas para evaluación automática de
resúmenes y textos generados mediante comparación de superposición de
n-gramas con un texto de referencia humano.

\textbf{SHAP} (\emph{SHapley Additive exPlanations}): Marco de
explicabilidad post-hoc que asigna a cada feature una contribución
marginal promediada sobre todas las coaliciones posibles, garantizando
consistencia axiomática (Lundberg \& Lee, 2017).

\textbf{SHAP Stability Index}: Métrica de coherencia de explicaciones
SHAP entre instancias similares, que certifica que el modelo asigna
importancias consistentes a features semejantes, requisito en contextos
forenses.

\textbf{TFT} (\emph{Temporal Fusion Transformer}): Arquitectura de
Transformer para forecasting multi-horizonte interpretable con
covariables exógenas, mecanismo de gating y predicción por cuantiles
(Lim et al., 2021).

\textbf{TreeSHAP}: Algoritmo exacto para el cálculo de valores SHAP en
modelos basados en árboles, con complejidad O(TLD²), que hace viable la
explicabilidad en GBDT de producción con millones de transacciones.

\textbf{XGBoost} (\emph{eXtreme Gradient Boosting}): Implementación
escalable de gradient boosting (Chen \& Guestrin, 2016) con
regularización L1/L2, manejo nativo de valores faltantes y
paralelización por columnas; baseline universal en competencias de ML
con datos tabulares.

\begin{center}\rule{0.5\linewidth}{0.5pt}\end{center}

\begin{center}\rule{0.5\linewidth}{0.5pt}\end{center}

\section{REFERENCIAS
BIBLIOGRÁFICAS}\label{referencias-bibliogruxe1ficas}

Almalki, F., \& Masud, M. (2025). \emph{Financial fraud detection using
explainable AI and stacking ensemble methods}. arXiv preprint
arXiv:2505.10050.

Ansari, A. F., Stella, L., Turkmen, C., Zhang, X., Mercado, P., Sinha,
R., \& Bergmeir, C. (2024). \emph{Chronos: Learning the language of time
series}. arXiv preprint arXiv:2403.07815.

Arik, S. O., \& Pfister, T. (2021). TabNet: Attentive interpretable
tabular learning. \emph{Proceedings of the AAAI Conference on Artificial
Intelligence}, \emph{35}(8), 6679--6687.
https://doi.org/10.1609/aaai.v35i8.16842

Barclays Research. (2025). \emph{Beyond the black box: Interpretability
of LLMs in finance}. arXiv preprint arXiv:2505.24650.

Breunig, M. M., Kriegel, H.-P., Ng, R. T., \& Sander, J. (2000). LOF:
Identifying density-based local outliers. \emph{ACM SIGMOD Record},
\emph{29}(2), 93--104. https://doi.org/10.1145/342009.335388

Center for Audit Quality. (2024). \emph{Auditing in the age of
generative AI}. CAQ.
https://thecaq.org/wp-content/uploads/2024/04/caq\_auditing-in-the-age-of-generative-ai\_\_2024-04.pdf

Challu, C., Olivares, K. G., Oreshkin, B., Garza, F.,
Mergenthaler-Canseco, M., \& Dubrawski, A. (2022). N-HiTS: Neural
hierarchical interpolation for time series forecasting. \emph{Advances
in Neural Information Processing Systems}, \emph{35}. (arXiv:2201.12886)

Chen, T., \& Guestrin, C. (2016). XGBoost: A scalable tree boosting
system. \emph{Proceedings of the 22nd ACM SIGKDD International
Conference on Knowledge Discovery and Data Mining}, 785--794.
https://doi.org/10.1145/2939672.2939785

Gebru, T., Morgenstern, J., Vecchione, B., Vaughan, J. W., Wallach, H.,
Daumé III, H., \& Crawford, K. (2021). Datasheets for datasets.
\emph{Communications of the ACM}, \emph{64}(12), 86--92.
https://doi.org/10.1145/3458723

Gorishniy, Y., Rubashevskiy, I., Khrulkov, V., \& Babenko, A. (2021).
Revisiting deep learning models for tabular data. \emph{Advances in
Neural Information Processing Systems 34 (NeurIPS)}, 8946--8959.

Grinsztajn, L., Oyallon, E., \& Varoquaux, G. (2022). Why do tree-based
models still outperform deep learning on tabular data? \emph{Advances in
Neural Information Processing Systems 35 (NeurIPS)}, 507--520.

Han, X., Hu, Y., Venevdev, L., Liu, M., Wen, Q., \& Zhang, Y. (2022).
ADBench: Anomaly detection benchmark. \emph{Advances in Neural
Information Processing Systems}, \emph{35}.

Hegselmann, S., Buendia, A., Lang, H., Agrawal, M., Jiang, X., \&
Sontag, D. (2022). TabLLM: Few-shot classification of tabular data with
large language models. arXiv preprint arXiv:2210.10723.

Hyndman, R. J., \& Khandakar, Y. (2008). Automatic time series
forecasting: The forecast package for R. \emph{Journal of Statistical
Software}, \emph{27}(3), 1--22. https://doi.org/10.18637/jss.v027.i03

Jesus, S., Pombal, J., Alves, D., Cruz, A., Saleiro, P., Ribeiro, R. P.,
Gama, J., \& Bizarro, P. (2022). \emph{Turning the tables: Biased,
imbalanced, dynamic tabular datasets for ML evaluation}. arXiv preprint
arXiv:2211.13358. {[}NeurIPS 2022{]}

Kadir, M. A., Macharla Vasu, S. S., Nair, S. S., \& Sonntag, D. (2025).
AuditCopilot: Leveraging LLMs for fraud detection in double-entry
bookkeeping. arXiv preprint arXiv:2512.02726.
https://doi.org/10.48550/arXiv.2512.02726

Ke, G., Meng, Q., Finley, T., Wang, T., Chen, W., Ma, W., Ye, Q., \&
Liu, T.-Y. (2017). LightGBM: A highly efficient gradient boosting
decision tree. \emph{Advances in Neural Information Processing Systems
30 (NeurIPS)}, 3146--3154.

Kreuzberger, D., Kühl, N., \& Hirschl, G. (2022). MLOps: Overview,
definition, and architecture. \emph{IEEE Access}, \emph{10},
86995--87010. https://doi.org/10.1109/ACCESS.2022.3197550

Leocádio, D., et al.~(2024). \emph{Continuous auditing artificial
intelligence framework}. {[}Pending venue publication{]}

Lewis, P., Perez, E., Piktus, A., Schwenk, H., Schwab, C., Yeh, C., \&
Kiela, D. (2020). Retrieval-augmented generation for knowledge-intensive
NLP tasks. \emph{Proceedings of the 2020 Conference on Empirical Methods
in Natural Language Processing (EMNLP)}, 9457--9474.
https://doi.org/10.18653/v1/2020.emnlp-main.727

Li, Z., Zhao, Y., Hu, X., Botta, N., Ionescu, C., \& Chen, G. H. (2022).
ECOD: Unsupervised outlier detection using empirical cumulative
distribution functions. \emph{IEEE Transactions on Knowledge and Data
Engineering}, \emph{35}(12), 12181--12193.
https://doi.org/10.1109/TKDE.2022.3159580

Lim, B., Arik, S. O., Loeff, N., \& Pfister, T. (2021). Temporal fusion
transformers for interpretable multi-horizon time series forecasting.
\emph{International Conference on Learning Representations (ICLR)}.
(arXiv:1912.09300)

Lin, C.-Y. (2004). ROUGE: A package for automatic evaluation of
summaries. \emph{Proceedings of the Workshop on Text Summarization
Branches Out (ACL 2004)}, 74--81.

Liu, F. T., Ting, K. M., \& Zhou, Z.-H. (2008). Isolation forest.
\emph{Proceedings of the 8th IEEE International Conference on Data
Mining (ICDM)}, 413--422. https://doi.org/10.1109/ICDM.2008.17

Liu, Y., Hu, T., Zhang, H., Wu, H., Wang, S., Ma, L., \& Long, M.
(2024). iTransformer: Inverted transformers are effective for time
series forecasting. \emph{International Conference on Learning
Representations (ICLR)}.

Lundberg, S. M., \& Lee, S.-I. (2017). A unified approach to
interpreting model predictions. \emph{Advances in Neural Information
Processing Systems 30 (NeurIPS)}, 4765--4774.

Mitchell, M., Wu, S., Zaldivar, A., Barnes, P., Vasserman, L.,
Hutchinson, B., Spitzer, E., Raji, I. D., \& Gebru, T. (2019). Model
cards for model reporting. \emph{Proceedings of the Conference on
Fairness, Accountability, and Transparency (FAccT)}, 220--229.
https://doi.org/10.1145/3287560.3287596

National Institute of Standards and Technology. (2023). \emph{Artificial
intelligence risk management framework (AI RMF 1.0)}. NIST.
https://doi.org/10.6028/NIST.AI.600-1

Nie, Y., Nguyen, N. H., Sinthong, P., \& Kalagnanam, J. (2023). A time
series is worth 64 words: Long-term forecasting with transformers.
\emph{International Conference on Learning Representations (ICLR)}.

Oreshkin, B. N., Carpov, D., Chapados, N., \& Bengio, Y. (2020).
N-BEATS: Neural basis expansion analysis for interpretable time series
forecasting. \emph{Proceedings of the 37th International Conference on
Machine Learning (ICML)}, 4799--4808.
https://doi.org/10.5555/3524938.3525255

Papineni, K., Roukos, S., Ward, T., \& Zhu, W.-J. (2002). BLEU: A method
for automatic evaluation of machine translation. \emph{Proceedings of
the 40th Annual Meeting of the Association for Computational Linguistics
(ACL)}, 311--318. https://doi.org/10.3115/1073083.1073135

Park, S. (2024). \emph{LLMs for anomaly validation and report generation
in financial systems}. arXiv preprint arXiv:2403.19735.

Patel, R., Khan, F., Silva, B., \& Shaturaev, J. (2024). AI-driven
continuous auditing and real-time financial monitoring.
\emph{International Research Journal of Modernization in Engineering
Technology and Science}.

Parlamento Europeo y Consejo de la Unión Europea. (2024).
\emph{Reglamento (UE) 2024/1689 por el que se establecen normas
armonizadas en materia de inteligencia artificial (Ley de Inteligencia
Artificial)}. Diario Oficial de la Unión Europea.
https://eur-lex.europa.eu/eli/reg/2024/1689

Prenio, J., \& Yong, J. (2024). \emph{Managing explanations: How
regulators can address AI explainability}. Bank for International
Settlements, Financial Stability Institute.
https://www.bis.org/fsi/fsipapers24.pdf

Presidencia del Consejo de Ministros. (2025, septiembre). \emph{Decreto
Supremo N° 115-2025-PCM: Reglamento de la Ley N° 31814}. Lima, Perú.
https://busquedas.elperuano.pe

Prokhorenkova, L., Gusev, G., Vorobev, A., Dorogush, A. V., \& Gulin, A.
(2018). CatBoost: Unbiased boosting with categorical features.
\emph{Advances in Neural Information Processing Systems 31 (NeurIPS)},
6638--6648. https://doi.org/10.5555/3327757.3327790

Ribeiro, M. T., Singh, S., \& Guestrin, C. (2016). ``Why should I trust
you?'': Explaining the predictions of any classifier. \emph{Proceedings
of the 22nd ACM SIGKDD International Conference on Knowledge Discovery
and Data Mining}, 1135--1144. https://doi.org/10.1145/2939672.2939778

Ruff, L., Kauffmann, J. R., Vandermeulen, R. A., Montavon, G., Samek,
W., Kloft, M., Dickhaus, T., \& Müller, K.-R. (2018). Deep one-class
classification. \emph{Proceedings of the 35th International Conference
on Machine Learning (ICML)}, 4393--4402.

Schneider, J., et al.~(2025). Retrieval-augmented generation (RAG).
\emph{Business \& Information Systems Engineering}.
https://doi.org/10.1007/s12599-025-00945-3

Sculley, D., Holt, G., Golovin, D., Davydov, E., Phillips, T., Ebner,
D., Chaudhary, V., \& Young, M. (2015). Hidden technical debt in machine
learning systems. \emph{Proceedings of the 28th International Conference
on Neural Information Processing Systems (NIPS) Workshop}.

Superintendencia de Banca, Seguros y AFP. (2023). \emph{Resolución SBS
N° 053-2023: Reglamento de gestión de riesgos de modelo}. Lima, Perú.
https://www.sbs.gob.pe

Taylor, S. J., \& Letham, B. (2017). \emph{Forecasting at scale}. PeerJ
Preprints. https://doi.org/10.7287/peerj.preprints.3190v2

Thanathamathee, P., et al.~(2024). SHAP-instance weighting and anchor
explainable AI: Enhancing XGBoost for financial fraud detection.
\emph{Emerging Science Journal}, \emph{8}(6).
https://doi.org/10.28991/ESJ-2024-08-06-024

Tsai, C.-P., et al.~(2025). \emph{LLM-based anomaly detection in tabular
data}. {[}Pending venue publication{]}

Varios autores. (2025). \emph{Hallucination detection and mitigation in
large language models}. arXiv preprint arXiv:2601.09929.

Varios autores. (2025). Financial statement fraud detection through an
integrated machine learning and explainable AI framework. \emph{Journal
of Risk and Financial Management}, \emph{19}(1), 13.
https://doi.org/10.3390/jrfm19010013

Varios autores. (2025). Explainable AI for forensic analysis: A
comparative study of SHAP and LIME in intrusion detection models.
\emph{Applied Sciences}, \emph{15}(13), 7329.
https://doi.org/10.3390/app15137329

Waltersdorfer, L., Ekaputra, F. J., Miksa, T., \& Sabou, M. (2024).
AuditMAI: Towards an infrastructure for continuous AI auditing. arXiv
preprint arXiv:2406.14243. https://doi.org/10.48550/arXiv.2406.14243

Zeng, A., Chen, M., Zhang, L., \& Xu, Q. (2023). Are transformers
effective for time series forecasting? \emph{Proceedings of the AAAI
Conference on Artificial Intelligence}, \emph{37}(9), 11121--11128.
https://doi.org/10.1609/aaai.v37i9.26317

Zhao, Y., Nasrullah, Z., \& Li, Z. (2019). PyOD: A Python toolbox for
scalable outlier detection. \emph{Journal of Machine Learning Research},
\emph{20}(96), 1--7.

\begin{center}\rule{0.5\linewidth}{0.5pt}\end{center}

\section{ANEXOS}\label{anexos}

\begin{center}\rule{0.5\linewidth}{0.5pt}\end{center}

\section{ANEXOS}\label{anexos-1}

\subsection{Anexo A --- Protocolo de Evaluación de
Usabilidad}\label{anexo-a-protocolo-de-evaluaciuxf3n-de-usabilidad}

\begin{quote}
\textbf{Versión 1.0 --- 2026-05-17.} Aprobar antes de iniciar
reclutamiento de participantes.
\end{quote}

\subsubsection{A.1 Objetivo del
experimento}\label{a.1-objetivo-del-experimento}

El experimento de usabilidad mide el impacto del sistema integrado de
supervisión operativa en la \textbf{eficiencia (VD4 ---
tiempo-a-decisión)}, \textbf{comprensión (VD4 --- Likert)} y
\textbf{trazabilidad documental (VD5)} frente al uso de componentes
aislados. Constituye la fuente principal de evidencia para contrastar
las sub-hipótesis H1b y H1d (Capítulo I §1.4).

\subsubsection{A.2 Diseño
experimental}\label{a.2-diseuxf1o-experimental}

\textbf{Tipo}: Cuasi-experimental con diseño within-subjects (apareado)
y orden contrabalanceado.

Cada participante ejecuta las mismas tareas en dos condiciones: -
\textbf{Condición A --- Sistema integrado}: pipeline de 4 capas con
alerta + score + vector SHAP top-5 + reporte LLM+RAG. -
\textbf{Condición B --- Componentes aislados}: alerta + score crudo del
detector, sin SHAP, sin reporte narrativo (solo tabla y visualización
técnica).

La mitad de los participantes inicia con A y la otra mitad con B,
asignación aleatorizada con \texttt{np.random.seed(42)}. Entre ambas
condiciones se intercala un descanso de 5 minutos y una tarea
distractora (sopa de letras) para reducir efectos de arrastre.

\subsubsection{A.3 Tareas evaluadas}\label{a.3-tareas-evaluadas}

Cada condición presenta el mismo bloque de 10 alertas (5 positivas
reales, 3 negativas reales, 2 ambiguas) extraídas del conjunto de test
del dataset sintético. Para cada alerta, el participante:

\begin{enumerate}
\def\labelenumi{\arabic{enumi}.}
\tightlist
\item
  \textbf{Tarea T1 --- Clasificación}: decide si la alerta corresponde a
  una anomalía operativa real (sí/no/dudoso).
\item
  \textbf{Tarea T2 --- Justificación}: en una oración (máx. 200
  caracteres) indica la variable que justifica la decisión.
\item
  \textbf{Tarea T3 --- Comprensión}: responde Likert 1--5 sobre cuán
  comprensible resultó la alerta.
\end{enumerate}

Al finalizar el bloque de 10 alertas, completa un cuestionario
post-bloque (SUS adaptado + preguntas abiertas).

\subsubsection{A.4 Métricas registradas
automáticamente}\label{a.4-muxe9tricas-registradas-automuxe1ticamente}

\begin{longtable}[]{@{}
  >{\raggedright\arraybackslash}p{(\columnwidth - 4\tabcolsep) * \real{0.3333}}
  >{\raggedright\arraybackslash}p{(\columnwidth - 4\tabcolsep) * \real{0.3333}}
  >{\raggedright\arraybackslash}p{(\columnwidth - 4\tabcolsep) * \real{0.3333}}@{}}
\toprule\noalign{}
\begin{minipage}[b]{\linewidth}\raggedright
Métrica
\end{minipage} & \begin{minipage}[b]{\linewidth}\raggedright
Fuente
\end{minipage} & \begin{minipage}[b]{\linewidth}\raggedright
Resolución
\end{minipage} \\
\midrule\noalign{}
\endhead
\bottomrule\noalign{}
\endlastfoot
Tiempo apertura alerta → decisión & Log JavaScript de la plataforma &
milisegundos \\
Tiempo total bloque & Log JavaScript & segundos \\
Clasificación del participante & Formulario & sí / no / dudoso \\
Justificación textual & Formulario & texto libre \\
Likert comprensión & Formulario & 1--5 \\
Versión del sistema & Variable de configuración & string \\
Identificador de alerta & Variable & string \\
\end{longtable}

\subsubsection{A.5 Criterios de inclusión y exclusión de
participantes}\label{a.5-criterios-de-inclusiuxf3n-y-exclusiuxf3n-de-participantes}

\textbf{Inclusión}: - (a) Estudiantes de pregrado ≥ 7° semestre en
Ingeniería de Sistemas, Industrial, Agronomía o Administración; o - (b)
Profesionales con ≥ 1 año de experiencia en supervisión operativa,
control de calidad, auditoría interna o cargos análogos. - Mayores de 18
años. - Aceptación de consentimiento informado.

\textbf{Exclusión}: - Participación previa en el diseño, desarrollo o
entrenamiento de cualquier capa del sistema evaluado. - Conflicto de
interés directo declarado. - Discapacidad visual no corregible que
impida la lectura del dashboard.

\subsubsection{A.6 Tamaño de muestra y
reclutamiento}\label{a.6-tamauxf1o-de-muestra-y-reclutamiento}

\textbf{Tamaño meta}: N ≥ 15 participantes (reportar resultados como
exploratorios). Tamaño deseable: N = 27 para potencia 0.80 con Cohen's
dz = 0.5 (efecto medio). N ≤ 14 obliga a marcar resultados como
preliminares y sin pretensión de significancia estadística.

\textbf{Reclutamiento}: difusión a través de UNSA (Escuela de Ingeniería
de Sistemas, Industrial y Agronomía); contactos directos con empresas
agroexportadoras de Arequipa, Ica y La Libertad mediante invitación
remitida por el asesor.

\subsubsection{A.7 Procedimiento detallado (sesión por participante,
\textasciitilde45
minutos)}\label{a.7-procedimiento-detallado-sesiuxf3n-por-participante-45-minutos}

\begin{longtable}[]{@{}lll@{}}
\toprule\noalign{}
Paso & Duración & Contenido \\
\midrule\noalign{}
\endhead
\bottomrule\noalign{}
\endlastfoot
1 & 5 min & Bienvenida + consentimiento informado firmado \\
2 & 5 min & Tutorial guiado de la plataforma (alerta de ejemplo) \\
3 & 12 min & Bloque 1 --- Condición A o B (según contrabalanceo) \\
4 & 5 min & Descanso + tarea distractora \\
5 & 12 min & Bloque 2 --- Condición contraria \\
6 & 5 min & Cuestionario final SUS + preguntas abiertas \\
7 & 1 min & Cierre + agradecimiento \\
\end{longtable}

\subsubsection{A.8 Consentimiento informado (texto
base)}\label{a.8-consentimiento-informado-texto-base}

\begin{verbatim}
Por la presente confirmo que:
1. He sido informado sobre el propósito del estudio: evaluar la usabilidad de
   un sistema de supervisión operativa con IA explicable.
2. Comprendo que mi participación es voluntaria y puedo retirarme en cualquier
   momento sin justificación ni consecuencia.
3. Comprendo que mis respuestas son anónimas. Solo el investigador principal
   accederá a los datos, que se almacenarán cifrados y se eliminarán al
   finalizar la tesis (julio 2027).
4. Acepto que se registre el tiempo de mis respuestas y las opciones
   seleccionadas para análisis estadístico agregado.
5. Comprendo que no recibiré evaluación individual ni se compartirán mis
   resultados con terceros.
6. Acepto participar en una sesión de aproximadamente 45 minutos.

Nombre: ____________________  Fecha: __________  Firma: __________
\end{verbatim}

\subsubsection{A.9 Cuestionario
post-bloque}\label{a.9-cuestionario-post-bloque}

\textbf{Bloque I --- Comprensión percibida (Likert 1--5)}

\begin{longtable}[]{@{}
  >{\raggedright\arraybackslash}p{(\columnwidth - 12\tabcolsep) * \real{0.1429}}
  >{\raggedright\arraybackslash}p{(\columnwidth - 12\tabcolsep) * \real{0.1429}}
  >{\raggedright\arraybackslash}p{(\columnwidth - 12\tabcolsep) * \real{0.1429}}
  >{\raggedright\arraybackslash}p{(\columnwidth - 12\tabcolsep) * \real{0.1429}}
  >{\raggedright\arraybackslash}p{(\columnwidth - 12\tabcolsep) * \real{0.1429}}
  >{\raggedright\arraybackslash}p{(\columnwidth - 12\tabcolsep) * \real{0.1429}}
  >{\raggedright\arraybackslash}p{(\columnwidth - 12\tabcolsep) * \real{0.1429}}@{}}
\toprule\noalign{}
\begin{minipage}[b]{\linewidth}\raggedright
\#
\end{minipage} & \begin{minipage}[b]{\linewidth}\raggedright
Ítem
\end{minipage} & \begin{minipage}[b]{\linewidth}\raggedright
1
\end{minipage} & \begin{minipage}[b]{\linewidth}\raggedright
2
\end{minipage} & \begin{minipage}[b]{\linewidth}\raggedright
3
\end{minipage} & \begin{minipage}[b]{\linewidth}\raggedright
4
\end{minipage} & \begin{minipage}[b]{\linewidth}\raggedright
5
\end{minipage} \\
\midrule\noalign{}
\endhead
\bottomrule\noalign{}
\endlastfoot
1 & Entendí claramente por qué el sistema marcó cada alerta & & & & & \\
2 & La información presentada me ayudó a tomar una decisión & & & & & \\
3 & Las variables explicativas fueron suficientes & & & & & \\
4 & El reporte generado fue útil para justificar mi decisión & & & &
& \\
5 & Confío en la decisión del sistema & & & & & \\
\end{longtable}

\textbf{Bloque II --- Tiempo y carga cognitiva (Likert 1--5)}

\begin{longtable}[]{@{}
  >{\raggedright\arraybackslash}p{(\columnwidth - 12\tabcolsep) * \real{0.1429}}
  >{\raggedright\arraybackslash}p{(\columnwidth - 12\tabcolsep) * \real{0.1429}}
  >{\raggedright\arraybackslash}p{(\columnwidth - 12\tabcolsep) * \real{0.1429}}
  >{\raggedright\arraybackslash}p{(\columnwidth - 12\tabcolsep) * \real{0.1429}}
  >{\raggedright\arraybackslash}p{(\columnwidth - 12\tabcolsep) * \real{0.1429}}
  >{\raggedright\arraybackslash}p{(\columnwidth - 12\tabcolsep) * \real{0.1429}}
  >{\raggedright\arraybackslash}p{(\columnwidth - 12\tabcolsep) * \real{0.1429}}@{}}
\toprule\noalign{}
\begin{minipage}[b]{\linewidth}\raggedright
\#
\end{minipage} & \begin{minipage}[b]{\linewidth}\raggedright
Ítem
\end{minipage} & \begin{minipage}[b]{\linewidth}\raggedright
1
\end{minipage} & \begin{minipage}[b]{\linewidth}\raggedright
2
\end{minipage} & \begin{minipage}[b]{\linewidth}\raggedright
3
\end{minipage} & \begin{minipage}[b]{\linewidth}\raggedright
4
\end{minipage} & \begin{minipage}[b]{\linewidth}\raggedright
5
\end{minipage} \\
\midrule\noalign{}
\endhead
\bottomrule\noalign{}
\endlastfoot
6 & Me tomó mucho tiempo entender cada alerta (1) --- Fui rápido (5) & &
& & & \\
7 & Tuve que esforzarme mucho mentalmente & & & & & \\
8 & Me sentí seguro al decidir & & & & & \\
\end{longtable}

\textbf{Bloque III --- Adaptación SUS (10 ítems Likert 1--5)}

\begin{verbatim}
1. Creo que me gustaría usar este sistema frecuentemente.
2. Encontré el sistema innecesariamente complejo.
3. Pensé que el sistema fue fácil de usar.
4. Necesitaría el apoyo de un experto técnico para usar este sistema.
5. Las funciones del sistema están bien integradas.
6. Hay demasiada inconsistencia en este sistema.
7. La mayoría aprendería a usar este sistema rápidamente.
8. Encontré el sistema muy engorroso.
9. Me sentí muy confiado al usar el sistema.
10. Necesitaría aprender muchas cosas antes de usar el sistema.
\end{verbatim}

\textbf{Bloque IV --- Preguntas abiertas}

\begin{enumerate}
\def\labelenumi{\arabic{enumi}.}
\tightlist
\item
  ¿Qué fue lo más útil del sistema en esta sesión?
\item
  ¿Qué información agregaría o quitaría?
\item
  ¿En qué situación operativa real este sistema sería más valioso?
\end{enumerate}

\subsubsection{A.10 Variables registradas para
análisis}\label{a.10-variables-registradas-para-anuxe1lisis}

\begin{longtable}[]{@{}lll@{}}
\toprule\noalign{}
Variable & Tipo & Origen \\
\midrule\noalign{}
\endhead
\bottomrule\noalign{}
\endlastfoot
\texttt{participant\_id} & string anónimo & Generado \\
\texttt{order} & \{AB, BA\} & Aleatorización \\
\texttt{condition} & \{integrated, isolated\} & Variable
independiente \\
\texttt{alert\_id} & string & Dataset \\
\texttt{gt\_label} & \{0, 1\} & Dataset (oculto al participante) \\
\texttt{user\_decision} & \{yes, no, dunno\} & Formulario \\
\texttt{time\_to\_decision\_ms} & int & Log JavaScript \\
\texttt{likert\_comprehension} & 1--5 & Cuestionario \\
\texttt{justification\_text} & string & Formulario \\
\texttt{sus\_score} & 0--100 & Cálculo SUS \\
\end{longtable}

\subsubsection{A.11 Plan de análisis
estadístico}\label{a.11-plan-de-anuxe1lisis-estaduxedstico}

\begin{enumerate}
\def\labelenumi{\arabic{enumi}.}
\tightlist
\item
  \textbf{Tiempo-a-decisión (VD4-a)}: t de Student apareado integrado
  vs.~aislado; Wilcoxon si Shapiro-Wilk rechaza normalidad. Reportar
  media ± DE, IC95\%, Cohen's dz.
\item
  \textbf{Comprensión Likert (VD4-b)}: Wilcoxon signed-rank apareado.
\item
  \textbf{Decisión correcta (VD4-c)}: McNemar sobre pares
  concordantes/discordantes.
\item
  \textbf{SUS}: comparación de medias con Mann-Whitney U.
\item
  \textbf{Análisis cualitativo de respuestas abiertas}: análisis
  temático con doble codificación independiente.
\end{enumerate}

\subsubsection{A.12 Almacenamiento y privacidad de
datos}\label{a.12-almacenamiento-y-privacidad-de-datos}

\begin{itemize}
\tightlist
\item
  Datos almacenados en archivo CSV cifrado con clave conocida solo por
  el investigador.
\item
  Identificadores anónimos, sin nombre, correo ni datos demográficos
  sensibles.
\item
  Backup en disco duro institucional UNSA.
\item
  Eliminación de datos crudos al cierre del proyecto (julio 2027).
\item
  Resultados agregados publicados en la tesis y en el repositorio
  GitHub.
\end{itemize}

\subsubsection{A.13 Aprobación ética}\label{a.13-aprobaciuxf3n-uxe9tica}

El protocolo se somete a revisión del asesor de tesis (Dr.~Víctor Manuel
Cornejo Aparicio) y, si la Escuela de Ingeniería de Sistemas dispone de
comité de ética, a su aprobación formal antes del reclutamiento.

\begin{center}\rule{0.5\linewidth}{0.5pt}\end{center}

\emph{Anexo A --- versión 1.0 --- 2026-05-17. Sometido a revisión final
antes de ejecutar el estudio.}

\begin{center}\rule{0.5\linewidth}{0.5pt}\end{center}

\subsection{Anexo B --- Model Cards del
Sistema}\label{anexo-b-model-cards-del-sistema}

\begin{quote}
\textbf{Estándar aplicado}: Mitchell et al.~(2019) --- Model Cards for
Model Reporting (Mitchell et al., 2019). \textbf{Estado}: 📐 Plantillas
listas, métricas pendientes de completar tras experimentos E1--E5.
\end{quote}

\begin{quote}
Cada Model Card sigue las 9 secciones de Mitchell et al.: detalles del
modelo, uso previsto, factores, métricas, datos de evaluación, datos de
entrenamiento, análisis cuantitativos, consideraciones éticas,
advertencias y recomendaciones.
\end{quote}

\begin{center}\rule{0.5\linewidth}{0.5pt}\end{center}

\subsubsection{B.1 Model Card --- Módulo de Predicción Tabular (XGBoost
/
LightGBM)}\label{b.1-model-card-muxf3dulo-de-predicciuxf3n-tabular-xgboost-lightgbm}

\textbf{1. Detalles del modelo} - Nombre: \texttt{module1\_prediction} -
Tipo: Ensemble de Gradient Boosting Decision Trees (XGBoost + LightGBM)
- Versión: 1.0 - Autor: Yoset Cozco Mauri (UNSA) - Licencia del modelo
entrenado: MIT - Citación recomendada: Cozco Mauri (2026), \emph{Tesis
UNSA}. - Hiperparámetros: definidos mediante Optuna TPE con 50 trials,
optimizando PR-AUC.

\textbf{2. Uso previsto} - Estimar el riesgo operativo (valor esperado
de score) de cada registro agroexportador. - Usuario primario: módulo de
detección de anomalías (capa 2) y supervisor operativo a través del
dashboard. - Uso fuera de alcance: clasificación de decisiones legales o
financieras automatizadas; no debe usarse en contextos distintos al
sector agroexportador peruano sin recalibración.

\textbf{3. Factores} - Producto, zona, mes, destino. - Variables
climáticas (temperatura, precipitación, humedad). - Cumplimiento
fitosanitario y días logísticos.

\textbf{4. Métricas} - PR-AUC, ROC-AUC, F1, Precision, Recall. -
Reportadas como media ± DE sobre 6 semillas en el test set.

\textbf{5. Datos de evaluación} - Dataset sintético agroexportador v1.0
(Anexo C), conjunto de test (20\% cronológicamente posterior).

\textbf{6. Datos de entrenamiento} - Dataset sintético v1.0, conjunto de
train (70\%) + validation (10\%).

\textbf{7. Análisis cuantitativos} \textbar{} Métrica \textbar{} Valor
(mean ± SD) \textbar{} \textbar---\textbar---\textbar{} \textbar{}
PR-AUC \textbar{} \emph{pendiente} \textbar{} \textbar{} ROC-AUC
\textbar{} \emph{pendiente} \textbar{} \textbar{} F1 \textbar{}
\emph{pendiente} \textbar{} \textbar{} Precision \textbar{}
\emph{pendiente} \textbar{} \textbar{} Recall \textbar{}
\emph{pendiente} \textbar{}

Análisis por subgrupos (fairness): \textbar{} Subgrupo \textbar{} PR-AUC
\textbar{} F1 \textbar{} \textbar---\textbar---\textbar---\textbar{}
\textbar{} Producto = arándano \textbar{} \emph{pendiente} \textbar{}
\emph{pendiente} \textbar{} \textbar{} Producto = uva \textbar{}
\emph{pendiente} \textbar{} \emph{pendiente} \textbar{} \textbar{}
Producto = palta \textbar{} \emph{pendiente} \textbar{} \emph{pendiente}
\textbar{} \textbar{} Producto = cacao \textbar{} \emph{pendiente}
\textbar{} \emph{pendiente} \textbar{} \textbar{} Producto = espárrago
\textbar{} \emph{pendiente} \textbar{} \emph{pendiente} \textbar{}
\textbar{} Zona = Ica \textbar{} \emph{pendiente} \textbar{}
\emph{pendiente} \textbar{} \textbar{} Zona = La Libertad \textbar{}
\emph{pendiente} \textbar{} \emph{pendiente} \textbar{}

\textbf{8. Consideraciones éticas} - Datos sintéticos: cero riesgo de
exposición de información personal o empresarial. - Posible sesgo
geográfico: solo 5 departamentos modelados. - Mitigación: documentación
explícita del alcance y limitaciones de generalización.

\textbf{9. Advertencias y recomendaciones} - Recalibrar antes de uso
operativo real con datos de empresa. - Monitorear data drift mensual (KS
test sobre distribuciones de input). - No usar como única fuente de
decisión; siempre con revisión humana.

\begin{center}\rule{0.5\linewidth}{0.5pt}\end{center}

\subsubsection{B.2 Model Card --- Módulo de Detección de Anomalías
(Ensemble IF + LOF +
ECOD)}\label{b.2-model-card-muxf3dulo-de-detecciuxf3n-de-anomaluxedas-ensemble-if-lof-ecod}

\textbf{1. Detalles del modelo} - Nombre: \texttt{module2\_anomaly} -
Tipo: Ensemble no supervisado (Isolation Forest + LOF + ECOD) - Versión:
1.0 - Infraestructura: PyOD (Zhao et al., 2019) - Estrategia de
combinación: promedio normalizado de scores (alternativa: voto por
mayoría con umbral por detector).

\textbf{2. Uso previsto} - Identificar registros operativos atípicos en
el dataset agroexportador. - Salida: score continuo (0--1) y bandera
binaria (≥ umbral → anomalía).

\textbf{3. Factores} - Mismas variables que módulo 1, con
preprocesamiento StandardScaler para LOF.

\textbf{4. Métricas} - PR-AUC, ROC-AUC, F1 con umbral óptimo,
Specificity. - Tasa de falsos positivos como métrica operativa (cuántas
alertas se generan al día).

\textbf{5. Datos de evaluación} - Conjunto de test del dataset sintético
v1.0.

\textbf{6. Datos de entrenamiento} - Conjunto de train (no se usan
etiquetas durante el fit --- entrenamiento no supervisado).

\textbf{7. Análisis cuantitativos} \textbar{} Detector \textbar{} PR-AUC
\textbar{} ROC-AUC \textbar{} F1 \textbar{} FPR \textbar{}
\textbar---\textbar---\textbar---\textbar---\textbar---\textbar{}
\textbar{} Isolation Forest individual \textbar{} \emph{pendiente}
\textbar{} \emph{pendiente} \textbar{} \emph{pendiente} \textbar{}
\emph{pendiente} \textbar{} \textbar{} LOF individual \textbar{}
\emph{pendiente} \textbar{} \emph{pendiente} \textbar{} \emph{pendiente}
\textbar{} \emph{pendiente} \textbar{} \textbar{} ECOD individual
\textbar{} \emph{pendiente} \textbar{} \emph{pendiente} \textbar{}
\emph{pendiente} \textbar{} \emph{pendiente} \textbar{} \textbar{}
Ensemble IF+LOF+ECOD \textbar{} \emph{pendiente} \textbar{}
\emph{pendiente} \textbar{} \emph{pendiente} \textbar{} \emph{pendiente}
\textbar{}

Por tipo de anomalía inyectada: \textbar{} Tipo \textbar{} Recall
ensemble \textbar{} Recall IF solo \textbar{}
\textbar---\textbar---\textbar---\textbar{} \textbar{} precio \textbar{}
\emph{pendiente} \textbar{} \emph{pendiente} \textbar{} \textbar{}
volumen \textbar{} \emph{pendiente} \textbar{} \emph{pendiente}
\textbar{} \textbar{} clima \textbar{} \emph{pendiente} \textbar{}
\emph{pendiente} \textbar{} \textbar{} logistica \textbar{}
\emph{pendiente} \textbar{} \emph{pendiente} \textbar{} \textbar{}
calidad \textbar{} \emph{pendiente} \textbar{} \emph{pendiente}
\textbar{}

\textbf{8. Consideraciones éticas} - Riesgo de sobre-alerta sobre
productos o zonas específicas si la distribución del dataset está
sesgada. - Mitigación: análisis de tasa de alerta por subgrupo y
calibración de umbral por categoría.

\textbf{9. Advertencias y recomendaciones} - Ajustar umbral según costo
operativo de falsos positivos en cada empresa. - No usar sin la capa de
explicabilidad SHAP (las alertas sin contexto generan ruido).

\begin{center}\rule{0.5\linewidth}{0.5pt}\end{center}

\subsubsection{B.3 Model Card --- Módulo de Explicabilidad
(TreeSHAP)}\label{b.3-model-card-muxf3dulo-de-explicabilidad-treeshap}

\textbf{1. Detalles del modelo} - Nombre: \texttt{module3\_shap} - Tipo:
TreeSHAP (cálculo exacto de valores de Shapley en árboles) - Versión:
1.0 (librería \texttt{shap}) - Fundamento: Lundberg \& Lee (2017);
TreeSHAP --- Lundberg et al.~(2020)

\textbf{2. Uso previsto} - Generar vectores de contribución por variable
para cada alerta. - Alimentar la capa LLM+RAG con evidencia cuantitativa
estructurada.

\textbf{3. Factores} - Las mismas variables del modelo predictor; las
contribuciones se reportan en la escala del logit del XGBoost.

\textbf{4. Métricas} - Cobertura top-k (k=3, k=5). - Consistencia ρ
entre alertas del mismo tipo (Spearman). - Claridad operativa (Likert
1--5 evaluada en estudio de usabilidad).

\textbf{5. Datos de evaluación} - Subconjunto de 100 alertas
seleccionadas aleatoriamente del test set.

\textbf{6. Datos de entrenamiento} - N/A --- SHAP es un método post-hoc,
no se entrena.

\textbf{7. Análisis cuantitativos} \textbar{} Métrica \textbar{} Valor
\textbar{} \textbar---\textbar---\textbar{} \textbar{} Cobertura top-5
(mediana) \textbar{} \emph{pendiente} \textbar{} \textbar{} Cobertura
top-3 (mediana) \textbar{} \emph{pendiente} \textbar{} \textbar{}
Consistencia ρ (mediana) \textbar{} \emph{pendiente} \textbar{}
\textbar{} Likert claridad (mediana) \textbar{} \emph{pendiente}
\textbar{}

\textbf{8. Consideraciones éticas} - SHAP puede inducir sobre-confianza
si se interpreta como causalidad. Es una atribución, no una causa. -
Mitigación: documentación explícita en el dashboard y en el reporte
generado.

\textbf{9. Advertencias y recomendaciones} - TreeSHAP es exacto solo
para árboles; no aplicar a modelos no-árbol. - El orden y magnitud de
las contribuciones depende del modelo predictor; cambiar el modelo
invalida explicaciones previas.

\begin{center}\rule{0.5\linewidth}{0.5pt}\end{center}

\subsubsection{B.4 Model Card --- Módulo de Generación de Reportes (LLM
+
RAG)}\label{b.4-model-card-muxf3dulo-de-generaciuxf3n-de-reportes-llm-rag}

\textbf{1. Detalles del modelo} - Nombre: \texttt{module4\_rag} -
Componentes: - LLM base: Anthropic Claude Sonnet 4.6 (o alternativa
local Llama 3.1 8B Instruct). - Retriever: BM25 sobre fuentes
agroexportadoras + vectores SHAP estructurados. - Prompt template:
estructurado con campos obligatorios (dato, modelo, score, umbral, SHAP,
fuente, recomendación). - Versión: 1.0 - Parámetros: temperature = 0.2,
max\_tokens = 800.

\textbf{2. Uso previsto} - Generar reporte narrativo de cada alerta
operativa, anclado en evidencias SHAP y fuentes recuperadas. - Usuario
primario: supervisor operativo, auditor interno.

\textbf{3. Factores} - Vector SHAP top-5 de la alerta. - Score y umbral
aplicado. - Fragmentos recuperados de la base de conocimiento (fuentes
MIDAGRI, SENAMHI, etc.).

\textbf{4. Métricas} - Rúbrica de 5 dimensiones (completitud,
consistencia numérica, accionabilidad, coherencia textual,
correspondencia con evidencias). - ROUGE-1, ROUGE-L cuando exista
referencia humana. - Kappa de Cohen entre dos revisores.

\textbf{5. Datos de evaluación} - 20 reportes generados a partir de
alertas seleccionadas aleatoriamente del test set. - Evaluación por 2
revisores independientes con la rúbrica del Capítulo III §3.3.

\textbf{6. Datos de entrenamiento} - N/A --- LLM no se fine-tunea. El
conocimiento operativo se inyecta vía RAG en tiempo de inferencia.

\textbf{7. Análisis cuantitativos} \textbar{} Dimensión \textbar{} Valor
(media de 2 revisores) \textbar{} Kappa Cohen \textbar{}
\textbar---\textbar---\textbar---\textbar{} \textbar{} Completitud
\textbar{} \emph{pendiente} \textbar{} \emph{pendiente} \textbar{}
\textbar{} Consistencia numérica \textbar{} \emph{pendiente} \textbar{}
\emph{pendiente} \textbar{} \textbar{} Accionabilidad \textbar{}
\emph{pendiente} \textbar{} \emph{pendiente} \textbar{} \textbar{}
Coherencia textual \textbar{} \emph{pendiente} \textbar{}
\emph{pendiente} \textbar{} \textbar{} Correspondencia con evidencias
\textbar{} \emph{pendiente} \textbar{} \emph{pendiente} \textbar{}
\textbar{} \textbf{Promedio total} \textbar{} \emph{pendiente}
\textbar{} \emph{pendiente} \textbar{} \textbar{} ROUGE-L (subset con
referencia humana) \textbar{} \emph{pendiente} \textbar{} --- \textbar{}

\textbf{8. Consideraciones éticas} - Riesgo de alucinación numérica
residual (intrinsic hallucination): un porcentaje de reportes puede
contener números no presentes en SHAP. - Mitigación: validación
posterior automática (regex extrae números del reporte y compara con
vector SHAP). - Dependencia de API comercial: variabilidad por versión
de modelo.

\textbf{9. Advertencias y recomendaciones} - Ningún reporte debe usarse
sin revisión humana en contextos operativos críticos. - Documentar
siempre la versión exacta del LLM al generar el reporte. - Verificar
coincidencia numérica entre reporte y SHAP antes de publicar.

\begin{center}\rule{0.5\linewidth}{0.5pt}\end{center}

\emph{Anexo B --- versión 1.0 --- 2026-05-17. Plantillas completas;
métricas se completan tras experimentos E1--E5.}

\begin{center}\rule{0.5\linewidth}{0.5pt}\end{center}

\subsection{Anexo C --- Datasheet del Dataset Sintético
Agroexportador}\label{anexo-c-datasheet-del-dataset-sintuxe9tico-agroexportador}

\begin{quote}
\textbf{Estándar aplicado}: Gebru et al.~(2021) --- Datasheets for
Datasets (Gebru et al., 2021). \textbf{Versión del dataset}: v1.0
\textbf{Fecha de creación prevista}: 2026-06-01 (Hito 2 del plan
general) \textbf{Estado actual}: 📐 Especificación completa lista para
generación.
\end{quote}

\begin{center}\rule{0.5\linewidth}{0.5pt}\end{center}

\subsubsection{C.1 Motivación}\label{c.1-motivaciuxf3n}

\textbf{¿Para qué se creó el dataset?} Para entrenar y evaluar de forma
reproducible el sistema integrado de supervisión operativa con IA
explicable propuesto en esta tesis, sin depender de datos privados de
empresas agroexportadoras. La opción sintética permite (a) controlar la
distribución de anomalías, (b) garantizar reproducibilidad mediante
semilla fija, y (c) publicar el dataset junto al paper.

\textbf{¿Quién lo creó?} Yoset Cozco Mauri, Escuela Profesional de
Ingeniería de Sistemas, Universidad Nacional de San Agustín de Arequipa
(UNSA). Bajo asesoría del Dr.~Víctor Manuel Cornejo Aparicio.

\textbf{¿Quién financió la creación?} Sin financiamiento externo.
Desarrollo en el marco de una tesis de pregrado.

\begin{center}\rule{0.5\linewidth}{0.5pt}\end{center}

\subsubsection{C.2 Composición}\label{c.2-composiciuxf3n}

\textbf{¿Qué instancias representa?} Cada fila representa un evento
operativo agroexportador diario por combinación de producto, zona y
destino. Una instancia agrupa indicadores de producción,
comercialización, clima, sanidad, logística y calidad para un día
específico.

\textbf{¿Cuántas instancias?} Mínimo 2,000 --- Recomendado 5,000 ---
Tope 10,000. La versión inicial v1.0 generará 2,000 instancias.

\textbf{¿Es muestra o universo?} Es una muestra sintética generada para
cubrir las distribuciones plausibles del sector durante el período
2022-01-01 a 2025-12-31.

\textbf{¿Qué datos contiene cada instancia?}

\begin{longtable}[]{@{}
  >{\raggedright\arraybackslash}p{(\columnwidth - 10\tabcolsep) * \real{0.1667}}
  >{\raggedright\arraybackslash}p{(\columnwidth - 10\tabcolsep) * \real{0.1667}}
  >{\raggedright\arraybackslash}p{(\columnwidth - 10\tabcolsep) * \real{0.1667}}
  >{\raggedright\arraybackslash}p{(\columnwidth - 10\tabcolsep) * \real{0.1667}}
  >{\raggedright\arraybackslash}p{(\columnwidth - 10\tabcolsep) * \real{0.1667}}
  >{\raggedright\arraybackslash}p{(\columnwidth - 10\tabcolsep) * \real{0.1667}}@{}}
\toprule\noalign{}
\begin{minipage}[b]{\linewidth}\raggedright
Variable
\end{minipage} & \begin{minipage}[b]{\linewidth}\raggedright
Tipo
\end{minipage} & \begin{minipage}[b]{\linewidth}\raggedright
Descripción
\end{minipage} & \begin{minipage}[b]{\linewidth}\raggedright
Rango plausible
\end{minipage} & \begin{minipage}[b]{\linewidth}\raggedright
Fuente para rangos
\end{minipage} & \begin{minipage}[b]{\linewidth}\raggedright
Distribución
\end{minipage} \\
\midrule\noalign{}
\endhead
\bottomrule\noalign{}
\endlastfoot
\texttt{id} & int & Identificador único & 1..N & --- & Secuencial \\
\texttt{fecha} & datetime & Día del evento & 2022-01-01 a 2025-12-31 &
--- & Uniforme \\
\texttt{producto} & category & Producto agroexportador & \{arándano,
uva, palta, cacao, espárrago\} & MIDAGRI & Ponderada por
participación \\
\texttt{zona} & category & Departamento productor & \{Ica, La Libertad,
Piura, Arequipa, Lima\} & MIDAGRI & Ponderada por área cultivada \\
\texttt{volumen\_kg} & float & Volumen del día (kg) & 500--50,000 &
MIDAGRI & LogNormal(μ=8, σ=1.2) \\
\texttt{precio\_kg\_usd} & float & Precio FOB por kg & 0.5--12.0 &
MIDAGRI / SUNAT & Normal por producto \\
\texttt{temperatura\_max\_c} & float & Temperatura máxima del día (°C) &
15--38 & SENAMHI & Normal por zona/mes \\
\texttt{temperatura\_min\_c} & float & Temperatura mínima del día (°C) &
5--22 & SENAMHI & Normal por zona/mes \\
\texttt{precipitacion\_mm} & float & Precipitación diaria (mm) & 0--200
& SENAMHI & Gamma(α=0.5, β=10) \\
\texttt{humedad\_pct} & float & Humedad relativa (\%) & 40--95 & SENAMHI
& Beta(α=8, β=3) \\
\texttt{destino\_mercado} & category & Mercado destino & \{EEUU, UE,
Asia, Otro\} & SUNAT & Ponderada por exportaciones \\
\texttt{cumplimiento\_fitosanitario} & binary & Lote cumple SENASA &
\{0, 1\} & SENASA & Bernoulli(p=0.92) \\
\texttt{dias\_logisticos} & int & Días desde cosecha a embarque & 3--45
& Estimado & LogNormal(μ=2.3, σ=0.5) \\
\texttt{merma\_pct} & float & Pérdida del lote (\%) & 0--30 & Estimado &
Beta(α=2, β=10) \\
\texttt{costo\_logistico\_usd\_kg} & float & Costo logístico unitario &
0.05--1.2 & Estimado & LogNormal \\
\texttt{tipo\_cambio\_pen\_usd} & float & Tipo de cambio del día &
3.5--4.2 & BCRP & Random walk con tendencia \\
\texttt{etiqueta\_anomalia} & binary & Es anomalía operativa & \{0, 1\}
& Inyección controlada & Bernoulli(p=0.12) \\
\texttt{tipo\_anomalia} & category & Tipo de anomalía & \{precio,
volumen, clima, logistica, calidad, none\} & Inyección controlada &
Definida por reglas \\
\end{longtable}

\textbf{Total}: 17 columnas; 2,000 filas en v1.0.

\textbf{¿Las etiquetas son confiables?} Sí, porque el dataset es
sintético y las etiquetas se asignan según las reglas de inyección
documentadas (§C.4). No hay ambigüedad por error humano de etiquetado.

\textbf{¿Falta algún dato?} Sí, se inyectan valores faltantes en el 3\%
de las filas para simular registros parciales (campos: humedad\_pct,
dias\_logisticos, costo\_logistico\_usd\_kg). Esto evalúa robustez del
modelo ante datos incompletos típicos de fuentes operativas.

\textbf{¿Las relaciones entre instancias son explícitas?} Las instancias
están temporalmente ordenadas. No hay relación de identidad (no es panel
longitudinal con seguimiento individual). Una misma combinación
(producto, zona, fecha) puede aparecer solo una vez.

\textbf{¿División train/test recomendada?} División cronológica: -
Train: 2022-01-01 a 2024-12-31 (70\%) - Validation: 2025-01-01 a
2025-04-30 (10\%) - Test: 2025-05-01 a 2025-12-31 (20\%)

Esta división evita data leakage temporal y simula el caso realista de
aplicar el modelo a períodos futuros.

\textbf{¿Hay datos sensibles?} No.~No hay datos personales, no se
identifican empresas reales, no se referencian transacciones
específicas. Es completamente sintético.

\begin{center}\rule{0.5\linewidth}{0.5pt}\end{center}

\subsubsection{C.3 Proceso de
recolección}\label{c.3-proceso-de-recolecciuxf3n}

\textbf{¿Cómo se generaron los datos?} Mediante un script Python
(\texttt{src/generate\_synthetic\_dataset.py}) que muestrea cada columna
según las distribuciones de la tabla §C.2 condicionadas a (producto,
zona, mes). La generación se realiza en cuatro pasos: 1. Muestreo de
variables base (fecha, producto, zona). 2. Muestreo condicional de
variables dependientes (clima por zona+mes; precio por producto+mes). 3.
Inyección de correlaciones plausibles (volumen ↑ → merma ↓ por economías
de escala; precipitacion ↑ → merma ↑ por daño post-cosecha). 4.
Inyección controlada de anomalías según reglas §C.4.

\textbf{¿Quién recopiló los datos?} N/A --- Datos sintéticos generados
por el autor con \texttt{numpy.random.default\_rng(seed=42)}.

\textbf{¿En qué período se generaron?} Generación prevista: 2026-05-31 a
2026-06-01.

\begin{center}\rule{0.5\linewidth}{0.5pt}\end{center}

\subsubsection{C.4 Inyección de
anomalías}\label{c.4-inyecciuxf3n-de-anomaluxedas}

\textbf{Distribución de tipos de anomalía} (sobre 12\% de filas marcadas
como anómalas):

\begin{longtable}[]{@{}
  >{\raggedright\arraybackslash}p{(\columnwidth - 6\tabcolsep) * \real{0.2500}}
  >{\raggedright\arraybackslash}p{(\columnwidth - 6\tabcolsep) * \real{0.2500}}
  >{\raggedright\arraybackslash}p{(\columnwidth - 6\tabcolsep) * \real{0.2500}}
  >{\raggedright\arraybackslash}p{(\columnwidth - 6\tabcolsep) * \real{0.2500}}@{}}
\toprule\noalign{}
\begin{minipage}[b]{\linewidth}\raggedright
Tipo
\end{minipage} & \begin{minipage}[b]{\linewidth}\raggedright
Proporción
\end{minipage} & \begin{minipage}[b]{\linewidth}\raggedright
Mecanismo de inyección
\end{minipage} & \begin{minipage}[b]{\linewidth}\raggedright
Variables afectadas
\end{minipage} \\
\midrule\noalign{}
\endhead
\bottomrule\noalign{}
\endlastfoot
\texttt{precio} & 30\% & \texttt{precio\_kg\_usd} \textgreater{}
percentil 99 o \textless{} percentil 1 del producto & precio\_kg\_usd \\
\texttt{volumen} & 25\% & \texttt{volumen\_kg} \textgreater{} media +
3·DE para producto/zona & volumen\_kg \\
\texttt{clima} & 20\% & \texttt{temperatura\_max\_c} \textgreater{} 38°C
∧ \texttt{precipitacion\_mm} \textless{} 1 mm (sequía con calor extremo)
& temperatura\_max\_c, precipitacion\_mm \\
\texttt{logistica} & 15\% & \texttt{dias\_logisticos} \textgreater{}
percentil 95 ∧ \texttt{cumplimiento\_fitosanitario} = 1 (demora a pesar
de cumplir) & dias\_logisticos \\
\texttt{calidad} & 10\% & \texttt{merma\_pct} \textgreater{} 25\% ∧
\texttt{precio\_kg\_usd} \textgreater{} mediana (pérdida con precio
alto) & merma\_pct \\
\end{longtable}

\textbf{Verificabilidad}: Para cada anomalía inyectada se registra
\texttt{tipo\_anomalia} y el campo \texttt{regla\_inyeccion} que
documenta los valores exactos que activaron la regla. Esto permite
auditoría posterior del proceso de generación.

\begin{center}\rule{0.5\linewidth}{0.5pt}\end{center}

\subsubsection{C.5 Preprocesamiento, limpieza y
etiquetado}\label{c.5-preprocesamiento-limpieza-y-etiquetado}

\textbf{¿Se aplicó preprocesamiento?} El dataset v1.0 se publica sin
preprocesamiento adicional. El script de pipeline
(\texttt{src/pipeline.py}) aplica: 1. Imputación de valores faltantes
(mediana para numéricas, moda para categóricas). 2. Codificación one-hot
de variables categóricas. 3. Escalamiento StandardScaler para variables
numéricas (solo para LOF; XGBoost no lo requiere). 4. Construcción de
features derivadas:
\texttt{temperatura\_rango\ =\ temperatura\_max\_c\ -\ temperatura\_min\_c},
\texttt{precio\_unitario\_zscore} por producto.

\textbf{¿Los datos crudos se conservan?} Sí. El CSV crudo
(\texttt{data/dataset\_agro\_sintetico\_v1.csv}) se versiona en Git LFS
y queda como referencia de entrada al pipeline.

\begin{center}\rule{0.5\linewidth}{0.5pt}\end{center}

\subsubsection{C.6 Usos previstos y no
previstos}\label{c.6-usos-previstos-y-no-previstos}

\textbf{¿Para qué se usará el dataset?} - Entrenar y evaluar el sistema
integrado de supervisión operativa de esta tesis. - Comparar baselines
(Isolation Forest individual, ensembles parciales). - Publicar como
benchmark abierto en el repositorio GitHub asociado al paper.

\textbf{¿Para qué NO debe usarse?} - No representa datos reales de
empresas específicas; no debe usarse para toma de decisiones operativas
en ninguna empresa real. - No se diseñó para análisis económico del
sector agroexportador peruano (los rangos son plausibles pero no son
estadísticas oficiales). - No es un benchmark estandarizado de la
comunidad de detección de anomalías.

\textbf{¿Existen tareas para las que el dataset sería inadecuado?} -
Análisis causal: no hay manipulación experimental real, solo inyección
de correlaciones plausibles. - Modelos de pronóstico de mercado: el
componente de precio no refleja la volatilidad real de los mercados
internacionales. - Sesgo geográfico/demográfico: solo se modelan 5
productos y 5 departamentos peruanos.

\begin{center}\rule{0.5\linewidth}{0.5pt}\end{center}

\subsubsection{C.7 Distribución y
licencia}\label{c.7-distribuciuxf3n-y-licencia}

\textbf{¿El dataset se distribuirá?} Sí, publicación bajo licencia
\textbf{CC BY 4.0} en el repositorio GitHub asociado a esta tesis. Esto
permite uso académico y comercial citando la fuente.

\textbf{¿Cómo se citará?}

\begin{verbatim}
Cozco Mauri, Y. (2026). Dataset Sintético Agroexportador Peruano v1.0 [Data set].
  Universidad Nacional de San Agustín de Arequipa.
  URL: [repositorio GitHub al publicar]
\end{verbatim}

\textbf{¿Habrá DOI?} Se solicitará DOI en Zenodo al cierre de la tesis
para garantizar persistencia.

\begin{center}\rule{0.5\linewidth}{0.5pt}\end{center}

\subsubsection{C.8 Mantenimiento}\label{c.8-mantenimiento}

\textbf{¿Quién mantendrá el dataset?} Yoset Cozco Mauri hasta diciembre
2027. Después, el repositorio queda como archivo histórico.

\textbf{¿Habrá actualizaciones?} - v1.0 --- 2,000 filas --- versión base
de la tesis. - v1.1 --- posible si se requieren ajustes durante
experimentos. - v2.0 --- versión extendida (5,000 filas) para
publicación en paper.

\textbf{¿Cómo se notificarán los cambios?} Mediante CHANGELOG.md en el
repositorio y release notes versionadas por Git.

\begin{center}\rule{0.5\linewidth}{0.5pt}\end{center}

\subsubsection{C.9 Consideraciones
éticas}\label{c.9-consideraciones-uxe9ticas}

\textbf{¿Hay riesgo de daño a personas o grupos?} No, dado que el
dataset es sintético y no contiene información de individuos ni empresas
reales. La elección de 5 productos y 5 departamentos cubre los
principales del sector peruano pero no excluye intencionalmente a otros
actores.

\textbf{¿Hay riesgo de uso indebido?} Riesgo bajo. El dataset es
claramente etiquetado como sintético en cada archivo (header, README,
paper). Cualquier uso operativo real sería contrario a la documentación.

\begin{center}\rule{0.5\linewidth}{0.5pt}\end{center}

\subsubsection{C.10 Documentación complementaria --- Fuentes públicas
utilizadas}\label{c.10-documentaciuxf3n-complementaria-fuentes-puxfablicas-utilizadas}

El dataset se calibra con rangos plausibles tomados de las siguientes
fuentes públicas. Para reproducir o ampliar el dataset, consultar:

\begin{longtable}[]{@{}
  >{\raggedright\arraybackslash}p{(\columnwidth - 6\tabcolsep) * \real{0.2500}}
  >{\raggedright\arraybackslash}p{(\columnwidth - 6\tabcolsep) * \real{0.2500}}
  >{\raggedright\arraybackslash}p{(\columnwidth - 6\tabcolsep) * \real{0.2500}}
  >{\raggedright\arraybackslash}p{(\columnwidth - 6\tabcolsep) * \real{0.2500}}@{}}
\toprule\noalign{}
\begin{minipage}[b]{\linewidth}\raggedright
Fuente
\end{minipage} & \begin{minipage}[b]{\linewidth}\raggedright
URL
\end{minipage} & \begin{minipage}[b]{\linewidth}\raggedright
Variable calibrada
\end{minipage} & \begin{minipage}[b]{\linewidth}\raggedright
Acceso
\end{minipage} \\
\midrule\noalign{}
\endhead
\bottomrule\noalign{}
\endlastfoot
MIDAGRI & https://www.gob.pe/midagri & volumen, precio, producto, zona &
Público \\
SENAMHI & https://www.senamhi.gob.pe & temperatura, precipitación,
humedad & Público \\
SENASA & https://www.gob.pe/senasa & cumplimiento\_fitosanitario &
Público \\
SUNAT & https://www.sunat.gob.pe & destino\_mercado, valor exportado &
Público \\
INEI & https://www.inei.gob.pe & tipo\_cambio, indicadores económicos &
Público \\
FAOSTAT & https://www.fao.org/faostat & producción nacional comparativa
& Público \\
UN Comtrade & https://comtradeplus.un.org & exportaciones por país
destino & Público \\
\end{longtable}

\begin{center}\rule{0.5\linewidth}{0.5pt}\end{center}

\subsubsection{C.11 Benchmark
complementario}\label{c.11-benchmark-complementario}

El \textbf{BAF Benchmark} (Jesus et al., 2022) se utiliza únicamente
como referencia metodológica para validación cruzada de la arquitectura
en datos tabulares desbalanceados con drift temporal, NO como evidencia
del dominio agroexportador. Se documenta su uso en §3.2 con esa
restricción.

\begin{center}\rule{0.5\linewidth}{0.5pt}\end{center}

\emph{Anexo C --- versión 1.0 --- 2026-05-17. Datasheet completo, listo
para generación de v1.0 del dataset.}

\begin{center}\rule{0.5\linewidth}{0.5pt}\end{center}

\subsection{Anexo D --- Registro de Uso de Herramientas de
IA}\label{anexo-d-registro-de-uso-de-herramientas-de-ia}

La presente investigación utilizó herramientas de inteligencia
artificial generativa como apoyo en las siguientes actividades: revisión
bibliográfica exploratoria, verificación de coherencia de argumentos,
corrección de estilo académico y generación de borradores de secciones
específicas. Todas las referencias bibliográficas fueron verificadas
manualmente en las fuentes originales. Las decisiones de diseño, la
interpretación de resultados y las conclusiones son responsabilidad
exclusiva del investigador.

\emph{(Adjuntar registro detallado de las sesiones de uso según los
requerimientos de transparencia de la UNSA)}

\begin{center}\rule{0.5\linewidth}{0.5pt}\end{center}

\emph{(Documento elaborado con apoyo de herramientas de IA --- UNSA
Arequipa, 2026)}
